\QHeading{第1题　用小数表示下图的涂色部分。}

\par\noindent\makebox[\linewidth][c]{
\begin{tikzpicture}[x=1bp,y=-1bp,color=black,draw=black,text=black]
\begin{scope}\clip (56.0,113.0) rectangle (518.0,261.33);
\hdhead{46}{小数的意义与性质——信息窗1}
\hdsection{1}{135}{用小数表示下图的涂色部分。}
\begin{scope}[scale=.6666667,draw=black,text=black]\path[pattern=north east lines,pattern color=black] (143,178) rectangle (200.6,323);\end{scope}
\begin{scope}[scale=.6666667,draw=black,text=black]\draw[draw=black,line width=.4pt] (143,178) rectangle (287,323);\foreach \i in {1,...,9}{\draw[draw=black,line width=.4pt] ({143+14.4*\i},178)--({143+14.4*\i},323);}\foreach \j in {1,...,0}{\draw[draw=black,line width=.4pt] (143,{178+145.0*\j})--(287,{178+145.0*\j});}\end{scope}
\begin{scope}[scale=.6666667,draw=black,text=black]\path[pattern=north east lines,pattern color=black] (389,178) rectangle (432.2,322);\end{scope}
\begin{scope}[scale=.6666667,draw=black,text=black]\path[pattern=north east lines,pattern color=black] (432.2,278.8) rectangle (446.59999999999997,322.0);\end{scope}
\begin{scope}[scale=.6666667,draw=black,text=black]\draw[draw=black,line width=.35pt] (389,178) rectangle (533,322);\foreach \i in {1,...,9}{\draw[draw=black,line width=.35pt] ({389+14.4*\i},178)--({389+14.4*\i},322);}\foreach \j in {1,...,9}{\draw[draw=black,line width=.35pt] (389,{178+14.4*\j})--(533,{178+14.4*\j});}\end{scope}
\begin{scope}[scale=.6666667,draw=black,text=black]\fill[white] (636.0000,319.0000)--(736.0000,319.0000)--(736.0000,219.0000)--(636.0000,219.0000)--cycle;\end{scope}
\begin{scope}[scale=.6666667,draw=black,text=black]\fill[white] (636.0000,219.0000)--(736.0000,219.0000)--(769.0000,179.0000)--(669.0000,179.0000)--cycle;\end{scope}
\begin{scope}[scale=.6666667,draw=black,text=black]\path[pattern=north east lines,pattern color=black] (736.0000,319.0000)--(769.0000,279.0000)--(769.0000,179.0000)--(736.0000,219.0000)--cycle;\end{scope}
\begin{scope}[scale=.6666667,draw=black,text=black]\path[pattern=north east lines,pattern color=black] (706.0000,249.0000)--(716.0000,249.0000)--(716.0000,239.0000)--(706.0000,239.0000)--cycle;\end{scope}
\begin{scope}[scale=.6666667,draw=black,text=black]\path[pattern=north east lines,pattern color=black] (706.0000,239.0000)--(716.0000,239.0000)--(716.0000,229.0000)--(706.0000,229.0000)--cycle;\end{scope}
\begin{scope}[scale=.6666667,draw=black,text=black]\path[pattern=north east lines,pattern color=black] (706.0000,229.0000)--(716.0000,229.0000)--(716.0000,219.0000)--(706.0000,219.0000)--cycle;\end{scope}
\begin{scope}[scale=.6666667,draw=black,text=black]\path[pattern=north east lines,pattern color=black] (706.0000,219.0000)--(716.0000,219.0000)--(719.3000,215.0000)--(709.3000,215.0000)--cycle;\end{scope}
\begin{scope}[scale=.6666667,draw=black,text=black]\path[pattern=north east lines,pattern color=black] (716.0000,319.0000)--(726.0000,319.0000)--(726.0000,309.0000)--(716.0000,309.0000)--cycle;\end{scope}
\begin{scope}[scale=.6666667,draw=black,text=black]\path[pattern=north east lines,pattern color=black] (716.0000,309.0000)--(726.0000,309.0000)--(726.0000,299.0000)--(716.0000,299.0000)--cycle;\end{scope}
\begin{scope}[scale=.6666667,draw=black,text=black]\path[pattern=north east lines,pattern color=black] (716.0000,299.0000)--(726.0000,299.0000)--(726.0000,289.0000)--(716.0000,289.0000)--cycle;\end{scope}
\begin{scope}[scale=.6666667,draw=black,text=black]\path[pattern=north east lines,pattern color=black] (716.0000,289.0000)--(726.0000,289.0000)--(726.0000,279.0000)--(716.0000,279.0000)--cycle;\end{scope}
\begin{scope}[scale=.6666667,draw=black,text=black]\path[pattern=north east lines,pattern color=black] (716.0000,279.0000)--(726.0000,279.0000)--(726.0000,269.0000)--(716.0000,269.0000)--cycle;\end{scope}
\begin{scope}[scale=.6666667,draw=black,text=black]\path[pattern=north east lines,pattern color=black] (716.0000,269.0000)--(726.0000,269.0000)--(726.0000,259.0000)--(716.0000,259.0000)--cycle;\end{scope}
\begin{scope}[scale=.6666667,draw=black,text=black]\path[pattern=north east lines,pattern color=black] (716.0000,259.0000)--(726.0000,259.0000)--(726.0000,249.0000)--(716.0000,249.0000)--cycle;\end{scope}
\begin{scope}[scale=.6666667,draw=black,text=black]\path[pattern=north east lines,pattern color=black] (716.0000,249.0000)--(726.0000,249.0000)--(726.0000,239.0000)--(716.0000,239.0000)--cycle;\end{scope}
\begin{scope}[scale=.6666667,draw=black,text=black]\path[pattern=north east lines,pattern color=black] (716.0000,239.0000)--(726.0000,239.0000)--(726.0000,229.0000)--(716.0000,229.0000)--cycle;\end{scope}
\begin{scope}[scale=.6666667,draw=black,text=black]\path[pattern=north east lines,pattern color=black] (716.0000,229.0000)--(726.0000,229.0000)--(726.0000,219.0000)--(716.0000,219.0000)--cycle;\end{scope}
\begin{scope}[scale=.6666667,draw=black,text=black]\path[pattern=north east lines,pattern color=black] (716.0000,219.0000)--(726.0000,219.0000)--(729.3000,215.0000)--(719.3000,215.0000)--cycle;\end{scope}
\begin{scope}[scale=.6666667,draw=black,text=black]\path[pattern=north east lines,pattern color=black] (719.3000,215.0000)--(729.3000,215.0000)--(732.6000,211.0000)--(722.6000,211.0000)--cycle;\end{scope}
\begin{scope}[scale=.6666667,draw=black,text=black]\path[pattern=north east lines,pattern color=black] (722.6000,211.0000)--(732.6000,211.0000)--(735.9000,207.0000)--(725.9000,207.0000)--cycle;\end{scope}
\begin{scope}[scale=.6666667,draw=black,text=black]\path[pattern=north east lines,pattern color=black] (725.9000,207.0000)--(735.9000,207.0000)--(739.2000,203.0000)--(729.2000,203.0000)--cycle;\end{scope}
\begin{scope}[scale=.6666667,draw=black,text=black]\path[pattern=north east lines,pattern color=black] (729.2000,203.0000)--(739.2000,203.0000)--(742.5000,199.0000)--(732.5000,199.0000)--cycle;\end{scope}
\begin{scope}[scale=.6666667,draw=black,text=black]\path[pattern=north east lines,pattern color=black] (732.5000,199.0000)--(742.5000,199.0000)--(745.8000,195.0000)--(735.8000,195.0000)--cycle;\end{scope}
\begin{scope}[scale=.6666667,draw=black,text=black]\path[pattern=north east lines,pattern color=black] (735.8000,195.0000)--(745.8000,195.0000)--(749.1000,191.0000)--(739.1000,191.0000)--cycle;\end{scope}
\begin{scope}[scale=.6666667,draw=black,text=black]\path[pattern=north east lines,pattern color=black] (739.1000,191.0000)--(749.1000,191.0000)--(752.4000,187.0000)--(742.4000,187.0000)--cycle;\end{scope}
\begin{scope}[scale=.6666667,draw=black,text=black]\path[pattern=north east lines,pattern color=black] (742.4000,187.0000)--(752.4000,187.0000)--(755.7000,183.0000)--(745.7000,183.0000)--cycle;\end{scope}
\begin{scope}[scale=.6666667,draw=black,text=black]\path[pattern=north east lines,pattern color=black] (745.7000,183.0000)--(755.7000,183.0000)--(759.0000,179.0000)--(749.0000,179.0000)--cycle;\end{scope}
\begin{scope}[scale=.6666667,draw=black,text=black]\path[pattern=north east lines,pattern color=black] (726.0000,319.0000)--(736.0000,319.0000)--(736.0000,309.0000)--(726.0000,309.0000)--cycle;\end{scope}
\begin{scope}[scale=.6666667,draw=black,text=black]\path[pattern=north east lines,pattern color=black] (726.0000,309.0000)--(736.0000,309.0000)--(736.0000,299.0000)--(726.0000,299.0000)--cycle;\end{scope}
\begin{scope}[scale=.6666667,draw=black,text=black]\path[pattern=north east lines,pattern color=black] (726.0000,299.0000)--(736.0000,299.0000)--(736.0000,289.0000)--(726.0000,289.0000)--cycle;\end{scope}
\begin{scope}[scale=.6666667,draw=black,text=black]\path[pattern=north east lines,pattern color=black] (726.0000,289.0000)--(736.0000,289.0000)--(736.0000,279.0000)--(726.0000,279.0000)--cycle;\end{scope}
\begin{scope}[scale=.6666667,draw=black,text=black]\path[pattern=north east lines,pattern color=black] (726.0000,279.0000)--(736.0000,279.0000)--(736.0000,269.0000)--(726.0000,269.0000)--cycle;\end{scope}
\begin{scope}[scale=.6666667,draw=black,text=black]\path[pattern=north east lines,pattern color=black] (726.0000,269.0000)--(736.0000,269.0000)--(736.0000,259.0000)--(726.0000,259.0000)--cycle;\end{scope}
\begin{scope}[scale=.6666667,draw=black,text=black]\path[pattern=north east lines,pattern color=black] (726.0000,259.0000)--(736.0000,259.0000)--(736.0000,249.0000)--(726.0000,249.0000)--cycle;\end{scope}
\begin{scope}[scale=.6666667,draw=black,text=black]\path[pattern=north east lines,pattern color=black] (726.0000,249.0000)--(736.0000,249.0000)--(736.0000,239.0000)--(726.0000,239.0000)--cycle;\end{scope}
\begin{scope}[scale=.6666667,draw=black,text=black]\path[pattern=north east lines,pattern color=black] (726.0000,239.0000)--(736.0000,239.0000)--(736.0000,229.0000)--(726.0000,229.0000)--cycle;\end{scope}
\begin{scope}[scale=.6666667,draw=black,text=black]\path[pattern=north east lines,pattern color=black] (726.0000,229.0000)--(736.0000,229.0000)--(736.0000,219.0000)--(726.0000,219.0000)--cycle;\end{scope}
\begin{scope}[scale=.6666667,draw=black,text=black]\path[pattern=north east lines,pattern color=black] (726.0000,219.0000)--(736.0000,219.0000)--(739.3000,215.0000)--(729.3000,215.0000)--cycle;\end{scope}
\begin{scope}[scale=.6666667,draw=black,text=black]\path[pattern=north east lines,pattern color=black] (729.3000,215.0000)--(739.3000,215.0000)--(742.6000,211.0000)--(732.6000,211.0000)--cycle;\end{scope}
\begin{scope}[scale=.6666667,draw=black,text=black]\path[pattern=north east lines,pattern color=black] (732.6000,211.0000)--(742.6000,211.0000)--(745.9000,207.0000)--(735.9000,207.0000)--cycle;\end{scope}
\begin{scope}[scale=.6666667,draw=black,text=black]\path[pattern=north east lines,pattern color=black] (735.9000,207.0000)--(745.9000,207.0000)--(749.2000,203.0000)--(739.2000,203.0000)--cycle;\end{scope}
\begin{scope}[scale=.6666667,draw=black,text=black]\path[pattern=north east lines,pattern color=black] (739.2000,203.0000)--(749.2000,203.0000)--(752.5000,199.0000)--(742.5000,199.0000)--cycle;\end{scope}
\begin{scope}[scale=.6666667,draw=black,text=black]\path[pattern=north east lines,pattern color=black] (742.5000,199.0000)--(752.5000,199.0000)--(755.8000,195.0000)--(745.8000,195.0000)--cycle;\end{scope}
\begin{scope}[scale=.6666667,draw=black,text=black]\path[pattern=north east lines,pattern color=black] (745.8000,195.0000)--(755.8000,195.0000)--(759.1000,191.0000)--(749.1000,191.0000)--cycle;\end{scope}
\begin{scope}[scale=.6666667,draw=black,text=black]\path[pattern=north east lines,pattern color=black] (749.1000,191.0000)--(759.1000,191.0000)--(762.4000,187.0000)--(752.4000,187.0000)--cycle;\end{scope}
\begin{scope}[scale=.6666667,draw=black,text=black]\path[pattern=north east lines,pattern color=black] (752.4000,187.0000)--(762.4000,187.0000)--(765.7000,183.0000)--(755.7000,183.0000)--cycle;\end{scope}
\begin{scope}[scale=.6666667,draw=black,text=black]\path[pattern=north east lines,pattern color=black] (755.7000,183.0000)--(765.7000,183.0000)--(769.0000,179.0000)--(759.0000,179.0000)--cycle;\end{scope}
\begin{scope}[scale=.6666667,draw=black,text=black]\draw[draw=black!85,line width=.32pt] (636.000,319.000)--(636.000,219.000);\end{scope}
\begin{scope}[scale=.6666667,draw=black,text=black]\draw[draw=black!85,line width=.32pt] (636.000,319.000)--(736.000,319.000);\end{scope}
\begin{scope}[scale=.6666667,draw=black,text=black]\draw[draw=black!85,line width=.32pt] (636.000,219.000)--(669.000,179.000);\end{scope}
\begin{scope}[scale=.6666667,draw=black,text=black]\draw[draw=black!85,line width=.32pt] (636.000,219.000)--(736.000,219.000);\end{scope}
\begin{scope}[scale=.6666667,draw=black,text=black]\draw[draw=black!85,line width=.32pt] (736.000,319.000)--(736.000,219.000);\end{scope}
\begin{scope}[scale=.6666667,draw=black,text=black]\draw[draw=black!85,line width=.32pt] (736.000,319.000)--(769.000,279.000);\end{scope}
\begin{scope}[scale=.6666667,draw=black,text=black]\draw[draw=black!85,line width=.32pt] (646.000,319.000)--(646.000,219.000);\end{scope}
\begin{scope}[scale=.6666667,draw=black,text=black]\draw[draw=black!85,line width=.32pt] (636.000,309.000)--(736.000,309.000);\end{scope}
\begin{scope}[scale=.6666667,draw=black,text=black]\draw[draw=black!85,line width=.32pt] (646.000,219.000)--(679.000,179.000);\end{scope}
\begin{scope}[scale=.6666667,draw=black,text=black]\draw[draw=black!85,line width=.32pt] (639.300,215.000)--(739.300,215.000);\end{scope}
\begin{scope}[scale=.6666667,draw=black,text=black]\draw[draw=black!85,line width=.32pt] (739.300,315.000)--(739.300,215.000);\end{scope}
\begin{scope}[scale=.6666667,draw=black,text=black]\draw[draw=black!85,line width=.32pt] (736.000,309.000)--(769.000,269.000);\end{scope}
\begin{scope}[scale=.6666667,draw=black,text=black]\draw[draw=black!85,line width=.32pt] (656.000,319.000)--(656.000,219.000);\end{scope}
\begin{scope}[scale=.6666667,draw=black,text=black]\draw[draw=black!85,line width=.32pt] (636.000,299.000)--(736.000,299.000);\end{scope}
\begin{scope}[scale=.6666667,draw=black,text=black]\draw[draw=black!85,line width=.32pt] (656.000,219.000)--(689.000,179.000);\end{scope}
\begin{scope}[scale=.6666667,draw=black,text=black]\draw[draw=black!85,line width=.32pt] (642.600,211.000)--(742.600,211.000);\end{scope}
\begin{scope}[scale=.6666667,draw=black,text=black]\draw[draw=black!85,line width=.32pt] (742.600,311.000)--(742.600,211.000);\end{scope}
\begin{scope}[scale=.6666667,draw=black,text=black]\draw[draw=black!85,line width=.32pt] (736.000,299.000)--(769.000,259.000);\end{scope}
\begin{scope}[scale=.6666667,draw=black,text=black]\draw[draw=black!85,line width=.32pt] (666.000,319.000)--(666.000,219.000);\end{scope}
\begin{scope}[scale=.6666667,draw=black,text=black]\draw[draw=black!85,line width=.32pt] (636.000,289.000)--(736.000,289.000);\end{scope}
\begin{scope}[scale=.6666667,draw=black,text=black]\draw[draw=black!85,line width=.32pt] (666.000,219.000)--(699.000,179.000);\end{scope}
\begin{scope}[scale=.6666667,draw=black,text=black]\draw[draw=black!85,line width=.32pt] (645.900,207.000)--(745.900,207.000);\end{scope}
\begin{scope}[scale=.6666667,draw=black,text=black]\draw[draw=black!85,line width=.32pt] (745.900,307.000)--(745.900,207.000);\end{scope}
\begin{scope}[scale=.6666667,draw=black,text=black]\draw[draw=black!85,line width=.32pt] (736.000,289.000)--(769.000,249.000);\end{scope}
\begin{scope}[scale=.6666667,draw=black,text=black]\draw[draw=black!85,line width=.32pt] (676.000,319.000)--(676.000,219.000);\end{scope}
\begin{scope}[scale=.6666667,draw=black,text=black]\draw[draw=black!85,line width=.32pt] (636.000,279.000)--(736.000,279.000);\end{scope}
\begin{scope}[scale=.6666667,draw=black,text=black]\draw[draw=black!85,line width=.32pt] (676.000,219.000)--(709.000,179.000);\end{scope}
\begin{scope}[scale=.6666667,draw=black,text=black]\draw[draw=black!85,line width=.32pt] (649.200,203.000)--(749.200,203.000);\end{scope}
\begin{scope}[scale=.6666667,draw=black,text=black]\draw[draw=black!85,line width=.32pt] (749.200,303.000)--(749.200,203.000);\end{scope}
\begin{scope}[scale=.6666667,draw=black,text=black]\draw[draw=black!85,line width=.32pt] (736.000,279.000)--(769.000,239.000);\end{scope}
\begin{scope}[scale=.6666667,draw=black,text=black]\draw[draw=black!85,line width=.32pt] (686.000,319.000)--(686.000,219.000);\end{scope}
\begin{scope}[scale=.6666667,draw=black,text=black]\draw[draw=black!85,line width=.32pt] (636.000,269.000)--(736.000,269.000);\end{scope}
\begin{scope}[scale=.6666667,draw=black,text=black]\draw[draw=black!85,line width=.32pt] (686.000,219.000)--(719.000,179.000);\end{scope}
\begin{scope}[scale=.6666667,draw=black,text=black]\draw[draw=black!85,line width=.32pt] (652.500,199.000)--(752.500,199.000);\end{scope}
\begin{scope}[scale=.6666667,draw=black,text=black]\draw[draw=black!85,line width=.32pt] (752.500,299.000)--(752.500,199.000);\end{scope}
\begin{scope}[scale=.6666667,draw=black,text=black]\draw[draw=black!85,line width=.32pt] (736.000,269.000)--(769.000,229.000);\end{scope}
\begin{scope}[scale=.6666667,draw=black,text=black]\draw[draw=black!85,line width=.32pt] (696.000,319.000)--(696.000,219.000);\end{scope}
\begin{scope}[scale=.6666667,draw=black,text=black]\draw[draw=black!85,line width=.32pt] (636.000,259.000)--(736.000,259.000);\end{scope}
\begin{scope}[scale=.6666667,draw=black,text=black]\draw[draw=black!85,line width=.32pt] (696.000,219.000)--(729.000,179.000);\end{scope}
\begin{scope}[scale=.6666667,draw=black,text=black]\draw[draw=black!85,line width=.32pt] (655.800,195.000)--(755.800,195.000);\end{scope}
\begin{scope}[scale=.6666667,draw=black,text=black]\draw[draw=black!85,line width=.32pt] (755.800,295.000)--(755.800,195.000);\end{scope}
\begin{scope}[scale=.6666667,draw=black,text=black]\draw[draw=black!85,line width=.32pt] (736.000,259.000)--(769.000,219.000);\end{scope}
\begin{scope}[scale=.6666667,draw=black,text=black]\draw[draw=black!85,line width=.32pt] (706.000,319.000)--(706.000,219.000);\end{scope}
\begin{scope}[scale=.6666667,draw=black,text=black]\draw[draw=black!85,line width=.32pt] (636.000,249.000)--(736.000,249.000);\end{scope}
\begin{scope}[scale=.6666667,draw=black,text=black]\draw[draw=black!85,line width=.32pt] (706.000,219.000)--(739.000,179.000);\end{scope}
\begin{scope}[scale=.6666667,draw=black,text=black]\draw[draw=black!85,line width=.32pt] (659.100,191.000)--(759.100,191.000);\end{scope}
\begin{scope}[scale=.6666667,draw=black,text=black]\draw[draw=black!85,line width=.32pt] (759.100,291.000)--(759.100,191.000);\end{scope}
\begin{scope}[scale=.6666667,draw=black,text=black]\draw[draw=black!85,line width=.32pt] (736.000,249.000)--(769.000,209.000);\end{scope}
\begin{scope}[scale=.6666667,draw=black,text=black]\draw[draw=black!85,line width=.32pt] (716.000,319.000)--(716.000,219.000);\end{scope}
\begin{scope}[scale=.6666667,draw=black,text=black]\draw[draw=black!85,line width=.32pt] (636.000,239.000)--(736.000,239.000);\end{scope}
\begin{scope}[scale=.6666667,draw=black,text=black]\draw[draw=black!85,line width=.32pt] (716.000,219.000)--(749.000,179.000);\end{scope}
\begin{scope}[scale=.6666667,draw=black,text=black]\draw[draw=black!85,line width=.32pt] (662.400,187.000)--(762.400,187.000);\end{scope}
\begin{scope}[scale=.6666667,draw=black,text=black]\draw[draw=black!85,line width=.32pt] (762.400,287.000)--(762.400,187.000);\end{scope}
\begin{scope}[scale=.6666667,draw=black,text=black]\draw[draw=black!85,line width=.32pt] (736.000,239.000)--(769.000,199.000);\end{scope}
\begin{scope}[scale=.6666667,draw=black,text=black]\draw[draw=black!85,line width=.32pt] (726.000,319.000)--(726.000,219.000);\end{scope}
\begin{scope}[scale=.6666667,draw=black,text=black]\draw[draw=black!85,line width=.32pt] (636.000,229.000)--(736.000,229.000);\end{scope}
\begin{scope}[scale=.6666667,draw=black,text=black]\draw[draw=black!85,line width=.32pt] (726.000,219.000)--(759.000,179.000);\end{scope}
\begin{scope}[scale=.6666667,draw=black,text=black]\draw[draw=black!85,line width=.32pt] (665.700,183.000)--(765.700,183.000);\end{scope}
\begin{scope}[scale=.6666667,draw=black,text=black]\draw[draw=black!85,line width=.32pt] (765.700,283.000)--(765.700,183.000);\end{scope}
\begin{scope}[scale=.6666667,draw=black,text=black]\draw[draw=black!85,line width=.32pt] (736.000,229.000)--(769.000,189.000);\end{scope}
\begin{scope}[scale=.6666667,draw=black,text=black]\draw[draw=black!85,line width=.32pt] (736.000,319.000)--(736.000,219.000);\end{scope}
\begin{scope}[scale=.6666667,draw=black,text=black]\draw[draw=black!85,line width=.32pt] (636.000,219.000)--(736.000,219.000);\end{scope}
\begin{scope}[scale=.6666667,draw=black,text=black]\draw[draw=black!85,line width=.32pt] (736.000,219.000)--(769.000,179.000);\end{scope}
\begin{scope}[scale=.6666667,draw=black,text=black]\draw[draw=black!85,line width=.32pt] (669.000,179.000)--(769.000,179.000);\end{scope}
\begin{scope}[scale=.6666667,draw=black,text=black]\draw[draw=black!85,line width=.32pt] (769.000,279.000)--(769.000,179.000);\end{scope}
\begin{scope}[scale=.6666667,draw=black,text=black]\draw[draw=black!85,line width=.32pt] (736.000,219.000)--(769.000,179.000);\end{scope}
\begin{scope}[scale=.6666667,draw=black,text=black]\draw[draw=black,line width=.55pt] (636.000,319.000)--(736.000,319.000)--(736.000,219.000)--(636.000,219.000)--cycle;\end{scope}
\begin{scope}[scale=.6666667,draw=black,text=black]\draw[draw=black,line width=.55pt] (636.000,219.000)--(736.000,219.000)--(769.000,179.000)--(669.000,179.000)--cycle;\end{scope}
\begin{scope}[scale=.6666667,draw=black,text=black]\draw[draw=black,line width=.55pt] (736.000,319.000)--(769.000,279.000)--(769.000,179.000)--(736.000,219.000)--cycle;\end{scope}
\begin{scope}[scale=.6666667,draw=black,text=black]\node[inner sep=0pt,font=\normalsize,text=black] at (123.000,189.000) {①};\end{scope}
\begin{scope}[scale=.6666667,draw=black,text=black]\node[inner sep=0pt,font=\normalsize,text=black] at (370.000,189.000) {②};\end{scope}
\begin{scope}[scale=.6666667,draw=black,text=black]\node[inner sep=0pt,font=\normalsize,text=black] at (617.000,190.000) {③};\end{scope}
\q{171}{339}{28}{13}{（}
\q{247}{339}{25}{13}{）}
\q{418}{339}{28}{13}{（}
\q{493}{339}{25}{13}{）}
\q{657}{338}{28}{13}{（}
\q{733}{338}{25}{13}{）}
\hdsection{2}{395}{给图形涂色表示它下面的小数。}
\begin{scope}[scale=.6666667,draw=black,text=black]\draw[draw=black,line width=.4pt] (143,444) rectangle (287,591);\foreach \i in {1,...,9}{\draw[draw=black,line width=.4pt] ({143+14.4*\i},444)--({143+14.4*\i},591);}\foreach \j in {1,...,0}{\draw[draw=black,line width=.4pt] (143,{444+147.0*\j})--(287,{444+147.0*\j});}\end{scope}
\begin{scope}[scale=.6666667,draw=black,text=black]\draw[draw=black,line width=.35pt] (389,444) rectangle (533,590);\foreach \i in {1,...,9}{\draw[draw=black,line width=.35pt] ({389+14.4*\i},444)--({389+14.4*\i},590);}\foreach \j in {1,...,9}{\draw[draw=black,line width=.35pt] (389,{444+14.6*\j})--(533,{444+14.6*\j});}\end{scope}
\begin{scope}[scale=.6666667,draw=black,text=black]\draw[draw=black,line width=.35pt] (636,444) rectangle (782,590);\foreach \i in {1,...,9}{\draw[draw=black,line width=.35pt] ({636+14.6*\i},444)--({636+14.6*\i},590);}\foreach \j in {1,...,9}{\draw[draw=black,line width=.35pt] (636,{444+14.6*\j})--(782,{444+14.6*\j});}\end{scope}
\begin{scope}[scale=.6666667,draw=black,text=black]\node[inner sep=0pt,font=\normalsize,text=black] at (123.000,451.000) {①};\end{scope}
\begin{scope}[scale=.6666667,draw=black,text=black]\node[inner sep=0pt,font=\normalsize,text=black] at (369.000,451.000) {②};\end{scope}
\begin{scope}[scale=.6666667,draw=black,text=black]\node[inner sep=0pt,font=\normalsize,text=black] at (616.000,451.000) {③};\end{scope}
\begin{scope}[scale=.6666667,draw=black,text=black]\node[inner sep=0pt,font=\normalsize,text=black] at (215.000,612.000) {$0.9$};\end{scope}
\begin{scope}[scale=.6666667,draw=black,text=black]\node[inner sep=0pt,font=\normalsize,text=black] at (461.000,612.000) {$0.09$};\end{scope}
\begin{scope}[scale=.6666667,draw=black,text=black]\node[inner sep=0pt,font=\normalsize,text=black] at (709.000,612.000) {$0.53$};\end{scope}
\hdsection{3}{657}{照样子填一填。}
\begin{scope}[scale=.6666667,draw=black,text=black]\draw[draw=black,line width=.55pt] (210.000,705.000)--(729.000,705.000)--(724.000,721.000)--(733.000,734.000)--(729.000,754.000)--(210.000,754.000)--cycle;\end{scope}
\begin{scope}[scale=.6666667,draw=black,text=black]\draw[draw=black,line width=.35pt] (219.000,705.000)--(219.000,719.000);\end{scope}
\begin{scope}[scale=.6666667,draw=black,text=black]\draw[draw=black,line width=.35pt] (223.767,705.000)--(223.767,711.000);\end{scope}
\begin{scope}[scale=.6666667,draw=black,text=black]\draw[draw=black,line width=.35pt] (228.533,705.000)--(228.533,711.000);\end{scope}
\begin{scope}[scale=.6666667,draw=black,text=black]\draw[draw=black,line width=.35pt] (233.300,705.000)--(233.300,711.000);\end{scope}
\begin{scope}[scale=.6666667,draw=black,text=black]\draw[draw=black,line width=.35pt] (238.067,705.000)--(238.067,711.000);\end{scope}
\begin{scope}[scale=.6666667,draw=black,text=black]\draw[draw=black,line width=.35pt] (242.833,705.000)--(242.833,714.000);\end{scope}
\begin{scope}[scale=.6666667,draw=black,text=black]\draw[draw=black,line width=.35pt] (247.600,705.000)--(247.600,711.000);\end{scope}
\begin{scope}[scale=.6666667,draw=black,text=black]\draw[draw=black,line width=.35pt] (252.367,705.000)--(252.367,711.000);\end{scope}
\begin{scope}[scale=.6666667,draw=black,text=black]\draw[draw=black,line width=.35pt] (257.133,705.000)--(257.133,711.000);\end{scope}
\begin{scope}[scale=.6666667,draw=black,text=black]\draw[draw=black,line width=.35pt] (261.900,705.000)--(261.900,711.000);\end{scope}
\begin{scope}[scale=.6666667,draw=black,text=black]\draw[draw=black,line width=.35pt] (266.667,705.000)--(266.667,719.000);\end{scope}
\begin{scope}[scale=.6666667,draw=black,text=black]\draw[draw=black,line width=.35pt] (271.433,705.000)--(271.433,711.000);\end{scope}
\begin{scope}[scale=.6666667,draw=black,text=black]\draw[draw=black,line width=.35pt] (276.200,705.000)--(276.200,711.000);\end{scope}
\begin{scope}[scale=.6666667,draw=black,text=black]\draw[draw=black,line width=.35pt] (280.967,705.000)--(280.967,711.000);\end{scope}
\begin{scope}[scale=.6666667,draw=black,text=black]\draw[draw=black,line width=.35pt] (285.733,705.000)--(285.733,711.000);\end{scope}
\begin{scope}[scale=.6666667,draw=black,text=black]\draw[draw=black,line width=.35pt] (290.500,705.000)--(290.500,714.000);\end{scope}
\begin{scope}[scale=.6666667,draw=black,text=black]\draw[draw=black,line width=.35pt] (295.267,705.000)--(295.267,711.000);\end{scope}
\begin{scope}[scale=.6666667,draw=black,text=black]\draw[draw=black,line width=.35pt] (300.033,705.000)--(300.033,711.000);\end{scope}
\begin{scope}[scale=.6666667,draw=black,text=black]\draw[draw=black,line width=.35pt] (304.800,705.000)--(304.800,711.000);\end{scope}
\begin{scope}[scale=.6666667,draw=black,text=black]\draw[draw=black,line width=.35pt] (309.567,705.000)--(309.567,711.000);\end{scope}
\begin{scope}[scale=.6666667,draw=black,text=black]\draw[draw=black,line width=.35pt] (314.333,705.000)--(314.333,719.000);\end{scope}
\begin{scope}[scale=.6666667,draw=black,text=black]\draw[draw=black,line width=.35pt] (319.100,705.000)--(319.100,711.000);\end{scope}
\begin{scope}[scale=.6666667,draw=black,text=black]\draw[draw=black,line width=.35pt] (323.867,705.000)--(323.867,711.000);\end{scope}
\begin{scope}[scale=.6666667,draw=black,text=black]\draw[draw=black,line width=.35pt] (328.633,705.000)--(328.633,711.000);\end{scope}
\begin{scope}[scale=.6666667,draw=black,text=black]\draw[draw=black,line width=.35pt] (333.400,705.000)--(333.400,711.000);\end{scope}
\begin{scope}[scale=.6666667,draw=black,text=black]\draw[draw=black,line width=.35pt] (338.167,705.000)--(338.167,714.000);\end{scope}
\begin{scope}[scale=.6666667,draw=black,text=black]\draw[draw=black,line width=.35pt] (342.933,705.000)--(342.933,711.000);\end{scope}
\begin{scope}[scale=.6666667,draw=black,text=black]\draw[draw=black,line width=.35pt] (347.700,705.000)--(347.700,711.000);\end{scope}
\begin{scope}[scale=.6666667,draw=black,text=black]\draw[draw=black,line width=.35pt] (352.467,705.000)--(352.467,711.000);\end{scope}
\begin{scope}[scale=.6666667,draw=black,text=black]\draw[draw=black,line width=.35pt] (357.233,705.000)--(357.233,711.000);\end{scope}
\begin{scope}[scale=.6666667,draw=black,text=black]\draw[draw=black,line width=.35pt] (362.000,705.000)--(362.000,719.000);\end{scope}
\begin{scope}[scale=.6666667,draw=black,text=black]\draw[draw=black,line width=.35pt] (366.767,705.000)--(366.767,711.000);\end{scope}
\begin{scope}[scale=.6666667,draw=black,text=black]\draw[draw=black,line width=.35pt] (371.533,705.000)--(371.533,711.000);\end{scope}
\begin{scope}[scale=.6666667,draw=black,text=black]\draw[draw=black,line width=.35pt] (376.300,705.000)--(376.300,711.000);\end{scope}
\begin{scope}[scale=.6666667,draw=black,text=black]\draw[draw=black,line width=.35pt] (381.067,705.000)--(381.067,711.000);\end{scope}
\begin{scope}[scale=.6666667,draw=black,text=black]\draw[draw=black,line width=.35pt] (385.833,705.000)--(385.833,714.000);\end{scope}
\begin{scope}[scale=.6666667,draw=black,text=black]\draw[draw=black,line width=.35pt] (390.600,705.000)--(390.600,711.000);\end{scope}
\begin{scope}[scale=.6666667,draw=black,text=black]\draw[draw=black,line width=.35pt] (395.367,705.000)--(395.367,711.000);\end{scope}
\begin{scope}[scale=.6666667,draw=black,text=black]\draw[draw=black,line width=.35pt] (400.133,705.000)--(400.133,711.000);\end{scope}
\begin{scope}[scale=.6666667,draw=black,text=black]\draw[draw=black,line width=.35pt] (404.900,705.000)--(404.900,711.000);\end{scope}
\begin{scope}[scale=.6666667,draw=black,text=black]\draw[draw=black,line width=.35pt] (409.667,705.000)--(409.667,719.000);\end{scope}
\begin{scope}[scale=.6666667,draw=black,text=black]\draw[draw=black,line width=.35pt] (414.433,705.000)--(414.433,711.000);\end{scope}
\begin{scope}[scale=.6666667,draw=black,text=black]\draw[draw=black,line width=.35pt] (419.200,705.000)--(419.200,711.000);\end{scope}
\begin{scope}[scale=.6666667,draw=black,text=black]\draw[draw=black,line width=.35pt] (423.967,705.000)--(423.967,711.000);\end{scope}
\begin{scope}[scale=.6666667,draw=black,text=black]\draw[draw=black,line width=.35pt] (428.733,705.000)--(428.733,711.000);\end{scope}
\begin{scope}[scale=.6666667,draw=black,text=black]\draw[draw=black,line width=.35pt] (433.500,705.000)--(433.500,714.000);\end{scope}
\begin{scope}[scale=.6666667,draw=black,text=black]\draw[draw=black,line width=.35pt] (438.267,705.000)--(438.267,711.000);\end{scope}
\begin{scope}[scale=.6666667,draw=black,text=black]\draw[draw=black,line width=.35pt] (443.033,705.000)--(443.033,711.000);\end{scope}
\begin{scope}[scale=.6666667,draw=black,text=black]\draw[draw=black,line width=.35pt] (447.800,705.000)--(447.800,711.000);\end{scope}
\begin{scope}[scale=.6666667,draw=black,text=black]\draw[draw=black,line width=.35pt] (452.567,705.000)--(452.567,711.000);\end{scope}
\begin{scope}[scale=.6666667,draw=black,text=black]\draw[draw=black,line width=.35pt] (457.333,705.000)--(457.333,719.000);\end{scope}
\begin{scope}[scale=.6666667,draw=black,text=black]\draw[draw=black,line width=.35pt] (462.100,705.000)--(462.100,711.000);\end{scope}
\begin{scope}[scale=.6666667,draw=black,text=black]\draw[draw=black,line width=.35pt] (466.867,705.000)--(466.867,711.000);\end{scope}
\begin{scope}[scale=.6666667,draw=black,text=black]\draw[draw=black,line width=.35pt] (471.633,705.000)--(471.633,711.000);\end{scope}
\begin{scope}[scale=.6666667,draw=black,text=black]\draw[draw=black,line width=.35pt] (476.400,705.000)--(476.400,711.000);\end{scope}
\begin{scope}[scale=.6666667,draw=black,text=black]\draw[draw=black,line width=.35pt] (481.167,705.000)--(481.167,714.000);\end{scope}
\begin{scope}[scale=.6666667,draw=black,text=black]\draw[draw=black,line width=.35pt] (485.933,705.000)--(485.933,711.000);\end{scope}
\begin{scope}[scale=.6666667,draw=black,text=black]\draw[draw=black,line width=.35pt] (490.700,705.000)--(490.700,711.000);\end{scope}
\begin{scope}[scale=.6666667,draw=black,text=black]\draw[draw=black,line width=.35pt] (495.467,705.000)--(495.467,711.000);\end{scope}
\begin{scope}[scale=.6666667,draw=black,text=black]\draw[draw=black,line width=.35pt] (500.233,705.000)--(500.233,711.000);\end{scope}
\begin{scope}[scale=.6666667,draw=black,text=black]\draw[draw=black,line width=.35pt] (505.000,705.000)--(505.000,719.000);\end{scope}
\begin{scope}[scale=.6666667,draw=black,text=black]\draw[draw=black,line width=.35pt] (509.767,705.000)--(509.767,711.000);\end{scope}
\begin{scope}[scale=.6666667,draw=black,text=black]\draw[draw=black,line width=.35pt] (514.533,705.000)--(514.533,711.000);\end{scope}
\begin{scope}[scale=.6666667,draw=black,text=black]\draw[draw=black,line width=.35pt] (519.300,705.000)--(519.300,711.000);\end{scope}
\begin{scope}[scale=.6666667,draw=black,text=black]\draw[draw=black,line width=.35pt] (524.067,705.000)--(524.067,711.000);\end{scope}
\begin{scope}[scale=.6666667,draw=black,text=black]\draw[draw=black,line width=.35pt] (528.833,705.000)--(528.833,714.000);\end{scope}
\begin{scope}[scale=.6666667,draw=black,text=black]\draw[draw=black,line width=.35pt] (533.600,705.000)--(533.600,711.000);\end{scope}
\begin{scope}[scale=.6666667,draw=black,text=black]\draw[draw=black,line width=.35pt] (538.367,705.000)--(538.367,711.000);\end{scope}
\begin{scope}[scale=.6666667,draw=black,text=black]\draw[draw=black,line width=.35pt] (543.133,705.000)--(543.133,711.000);\end{scope}
\begin{scope}[scale=.6666667,draw=black,text=black]\draw[draw=black,line width=.35pt] (547.900,705.000)--(547.900,711.000);\end{scope}
\begin{scope}[scale=.6666667,draw=black,text=black]\draw[draw=black,line width=.35pt] (552.667,705.000)--(552.667,719.000);\end{scope}
\begin{scope}[scale=.6666667,draw=black,text=black]\draw[draw=black,line width=.35pt] (557.433,705.000)--(557.433,711.000);\end{scope}
\begin{scope}[scale=.6666667,draw=black,text=black]\draw[draw=black,line width=.35pt] (562.200,705.000)--(562.200,711.000);\end{scope}
\begin{scope}[scale=.6666667,draw=black,text=black]\draw[draw=black,line width=.35pt] (566.967,705.000)--(566.967,711.000);\end{scope}
\begin{scope}[scale=.6666667,draw=black,text=black]\draw[draw=black,line width=.35pt] (571.733,705.000)--(571.733,711.000);\end{scope}
\begin{scope}[scale=.6666667,draw=black,text=black]\draw[draw=black,line width=.35pt] (576.500,705.000)--(576.500,714.000);\end{scope}
\begin{scope}[scale=.6666667,draw=black,text=black]\draw[draw=black,line width=.35pt] (581.267,705.000)--(581.267,711.000);\end{scope}
\begin{scope}[scale=.6666667,draw=black,text=black]\draw[draw=black,line width=.35pt] (586.033,705.000)--(586.033,711.000);\end{scope}
\begin{scope}[scale=.6666667,draw=black,text=black]\draw[draw=black,line width=.35pt] (590.800,705.000)--(590.800,711.000);\end{scope}
\begin{scope}[scale=.6666667,draw=black,text=black]\draw[draw=black,line width=.35pt] (595.567,705.000)--(595.567,711.000);\end{scope}
\begin{scope}[scale=.6666667,draw=black,text=black]\draw[draw=black,line width=.35pt] (600.333,705.000)--(600.333,719.000);\end{scope}
\begin{scope}[scale=.6666667,draw=black,text=black]\draw[draw=black,line width=.35pt] (605.100,705.000)--(605.100,711.000);\end{scope}
\begin{scope}[scale=.6666667,draw=black,text=black]\draw[draw=black,line width=.35pt] (609.867,705.000)--(609.867,711.000);\end{scope}
\begin{scope}[scale=.6666667,draw=black,text=black]\draw[draw=black,line width=.35pt] (614.633,705.000)--(614.633,711.000);\end{scope}
\begin{scope}[scale=.6666667,draw=black,text=black]\draw[draw=black,line width=.35pt] (619.400,705.000)--(619.400,711.000);\end{scope}
\begin{scope}[scale=.6666667,draw=black,text=black]\draw[draw=black,line width=.35pt] (624.167,705.000)--(624.167,714.000);\end{scope}
\begin{scope}[scale=.6666667,draw=black,text=black]\draw[draw=black,line width=.35pt] (628.933,705.000)--(628.933,711.000);\end{scope}
\begin{scope}[scale=.6666667,draw=black,text=black]\draw[draw=black,line width=.35pt] (633.700,705.000)--(633.700,711.000);\end{scope}
\begin{scope}[scale=.6666667,draw=black,text=black]\draw[draw=black,line width=.35pt] (638.467,705.000)--(638.467,711.000);\end{scope}
\begin{scope}[scale=.6666667,draw=black,text=black]\draw[draw=black,line width=.35pt] (643.233,705.000)--(643.233,711.000);\end{scope}
\begin{scope}[scale=.6666667,draw=black,text=black]\draw[draw=black,line width=.35pt] (648.000,705.000)--(648.000,719.000);\end{scope}
\begin{scope}[scale=.6666667,draw=black,text=black]\draw[draw=black,line width=.35pt] (652.767,705.000)--(652.767,711.000);\end{scope}
\begin{scope}[scale=.6666667,draw=black,text=black]\draw[draw=black,line width=.35pt] (657.533,705.000)--(657.533,711.000);\end{scope}
\begin{scope}[scale=.6666667,draw=black,text=black]\draw[draw=black,line width=.35pt] (662.300,705.000)--(662.300,711.000);\end{scope}
\begin{scope}[scale=.6666667,draw=black,text=black]\draw[draw=black,line width=.35pt] (667.067,705.000)--(667.067,711.000);\end{scope}
\begin{scope}[scale=.6666667,draw=black,text=black]\draw[draw=black,line width=.35pt] (671.833,705.000)--(671.833,714.000);\end{scope}
\begin{scope}[scale=.6666667,draw=black,text=black]\draw[draw=black,line width=.35pt] (676.600,705.000)--(676.600,711.000);\end{scope}
\begin{scope}[scale=.6666667,draw=black,text=black]\draw[draw=black,line width=.35pt] (681.367,705.000)--(681.367,711.000);\end{scope}
\begin{scope}[scale=.6666667,draw=black,text=black]\draw[draw=black,line width=.35pt] (686.133,705.000)--(686.133,711.000);\end{scope}
\begin{scope}[scale=.6666667,draw=black,text=black]\draw[draw=black,line width=.35pt] (690.900,705.000)--(690.900,711.000);\end{scope}
\begin{scope}[scale=.6666667,draw=black,text=black]\draw[draw=black,line width=.35pt] (695.667,705.000)--(695.667,719.000);\end{scope}
\begin{scope}[scale=.6666667,draw=black,text=black]\node[inner sep=0pt,font=\normalsize,text=black] at (219.000,738.000) {0};\end{scope}
\begin{scope}[scale=.6666667,draw=black,text=black]\node[inner sep=0pt,font=\normalsize,text=black] at (266.667,738.000) {1 cm};\end{scope}
\begin{scope}[scale=.6666667,draw=black,text=black]\node[inner sep=0pt,font=\normalsize,text=black] at (314.333,738.000) {2};\end{scope}
\begin{scope}[scale=.6666667,draw=black,text=black]\node[inner sep=0pt,font=\normalsize,text=black] at (362.000,738.000) {3};\end{scope}
\begin{scope}[scale=.6666667,draw=black,text=black]\node[inner sep=0pt,font=\normalsize,text=black] at (409.667,738.000) {4};\end{scope}
\begin{scope}[scale=.6666667,draw=black,text=black]\node[inner sep=0pt,font=\normalsize,text=black] at (457.333,738.000) {5};\end{scope}
\begin{scope}[scale=.6666667,draw=black,text=black]\node[inner sep=0pt,font=\normalsize,text=black] at (505.000,738.000) {6};\end{scope}
\begin{scope}[scale=.6666667,draw=black,text=black]\node[inner sep=0pt,font=\normalsize,text=black] at (552.667,738.000) {7};\end{scope}
\begin{scope}[scale=.6666667,draw=black,text=black]\node[inner sep=0pt,font=\normalsize,text=black] at (600.333,738.000) {8};\end{scope}
\begin{scope}[scale=.6666667,draw=black,text=black]\node[inner sep=0pt,font=\normalsize,text=black] at (648.000,738.000) {9};\end{scope}
\begin{scope}[scale=.6666667,draw=black,text=black]\node[inner sep=0pt,font=\normalsize,text=black] at (695.667,738.000) {10};\end{scope}
\begin{scope}[scale=.6666667,draw=black,text=black]\draw[draw=black,line width=.45pt,-{Stealth[length=3pt,width=2.4pt]}] (267.000,760.000)--(267.000,785.000);\end{scope}
\begin{scope}[scale=.6666667,draw=black,text=black]\draw[draw=black,line width=.45pt,-{Stealth[length=3pt,width=2.4pt]}] (362.000,760.000)--(362.000,785.000);\end{scope}
\begin{scope}[scale=.6666667,draw=black,text=black]\draw[draw=black,line width=.45pt,-{Stealth[length=3pt,width=2.4pt]}] (505.000,760.000)--(505.000,785.000);\end{scope}
\q{230}{790}{85}{12.5}{$\dfrac1{10}$分米}
\q{331}{790}{96}{12.5}{$\dfrac3{10}$分米}
\q{463}{804}{28}{13}{（}
\q{512}{804}{25}{13}{）}
\q{533}{804}{80}{12.5}{分米}
\q{225}{862}{103}{12.5}{0.1分米}
\q{324}{861}{28}{13}{（}
\q{378}{861}{25}{13}{）}
\q{399}{861}{76}{12.5}{分米}
\q{463}{861}{28}{13}{（}
\q{514}{861}{25}{13}{）}
\q{535}{861}{75}{12.5}{分米}
\hdsection{4}{921}{在计数器上表示下面各小数。}
\q{112}{967}{172}{13}{① 0.35}
\q{368}{967}{172}{13}{② 4.3}
\q{623}{967}{172}{13}{③ 0.732}
\begin{scope}[scale=.6666667,draw=black,text=black]\draw[line width=.5pt] (143,1015) rectangle (287,1079);\draw[line width=.5pt] (143,1019)--(287,1019);\end{scope}
\begin{scope}[scale=.6666667,draw=black,text=black]\draw[line width=.45pt] (165,1017)--(165,1079);\draw[line width=.45pt] (157,1015)..controls(166,995) and (171,1003)..(170,1015);\end{scope}
\begin{scope}[scale=.6666667,draw=black,text=black]\draw[line width=.45pt] (191,1017)--(191,1079);\draw[line width=.45pt] (183,1015)..controls(192,995) and (197,1003)..(196,1015);\end{scope}
\begin{scope}[scale=.6666667,draw=black,text=black]\draw[line width=.45pt] (224,1017)--(224,1079);\draw[line width=.45pt] (216,1015)..controls(225,995) and (230,1003)..(229,1015);\end{scope}
\begin{scope}[scale=.6666667,draw=black,text=black]\draw[line width=.45pt] (247,1017)--(247,1079);\draw[line width=.45pt] (239,1015)..controls(248,995) and (253,1003)..(252,1015);\end{scope}
\begin{scope}[scale=.6666667,draw=black,text=black]\draw[line width=.45pt] (273,1017)--(273,1079);\draw[line width=.45pt] (265,1015)..controls(274,995) and (279,1003)..(278,1015);\end{scope}
\begin{scope}[scale=.6666667,draw=black,text=black]\draw[draw=black,line width=.5pt,fill=white] (143.000,1078.000)--(151.000,1083.000)--(294.000,1083.000)--(287.000,1078.000)--cycle;\end{scope}
\begin{scope}[scale=.6666667,draw=black,text=black]\draw[draw=black,line width=.5pt,fill=white] (151.000,1083.000)--(294.000,1083.000)--(294.000,1147.000)--(151.000,1147.000)--cycle;\end{scope}
\begin{scope}[scale=.6666667,draw=black,text=black]\draw[draw=black,line width=.5pt,fill=white] (143.000,1078.000)--(151.000,1083.000)--(151.000,1147.000)--(143.000,1131.000)--cycle;\end{scope}
\begin{scope}[scale=.6666667,draw=black,text=black]\node[inner sep=0pt,font=\normalsize,text=black] at (165.000,1115.000) {\shortstack{十\\位}};\end{scope}
\begin{scope}[scale=.6666667,draw=black,text=black]\node[inner sep=0pt,font=\normalsize,text=black] at (191.000,1115.000) {\shortstack{个\\位}};\end{scope}
\begin{scope}[scale=.6666667,draw=black,text=black]\node[inner sep=0pt,font=\normalsize,text=black] at (224.000,1115.000) {\shortstack{十\\分\\位}};\end{scope}
\begin{scope}[scale=.6666667,draw=black,text=black]\node[inner sep=0pt,font=\normalsize,text=black] at (247.000,1115.000) {\shortstack{百\\分\\位}};\end{scope}
\begin{scope}[scale=.6666667,draw=black,text=black]\node[inner sep=0pt,font=\normalsize,text=black] at (273.000,1115.000) {\shortstack{千\\分\\位}};\end{scope}
\begin{scope}[scale=.6666667,draw=black,text=black]\draw[line width=.5pt] (398,1015) rectangle (542,1079);\draw[line width=.5pt] (398,1019)--(542,1019);\end{scope}
\begin{scope}[scale=.6666667,draw=black,text=black]\draw[line width=.45pt] (416,1017)--(416,1079);\draw[line width=.45pt] (408,1015)..controls(417,995) and (422,1003)..(421,1015);\end{scope}
\begin{scope}[scale=.6666667,draw=black,text=black]\draw[line width=.45pt] (441,1017)--(441,1079);\draw[line width=.45pt] (433,1015)..controls(442,995) and (447,1003)..(446,1015);\end{scope}
\begin{scope}[scale=.6666667,draw=black,text=black]\draw[line width=.45pt] (477,1017)--(477,1079);\draw[line width=.45pt] (469,1015)..controls(478,995) and (483,1003)..(482,1015);\end{scope}
\begin{scope}[scale=.6666667,draw=black,text=black]\draw[line width=.45pt] (502,1017)--(502,1079);\draw[line width=.45pt] (494,1015)..controls(503,995) and (508,1003)..(507,1015);\end{scope}
\begin{scope}[scale=.6666667,draw=black,text=black]\draw[line width=.45pt] (527,1017)--(527,1079);\draw[line width=.45pt] (519,1015)..controls(528,995) and (533,1003)..(532,1015);\end{scope}
\begin{scope}[scale=.6666667,draw=black,text=black]\draw[draw=black,line width=.5pt,fill=white] (398.000,1078.000)--(406.000,1083.000)--(549.000,1083.000)--(542.000,1078.000)--cycle;\end{scope}
\begin{scope}[scale=.6666667,draw=black,text=black]\draw[draw=black,line width=.5pt,fill=white] (406.000,1083.000)--(549.000,1083.000)--(549.000,1147.000)--(406.000,1147.000)--cycle;\end{scope}
\begin{scope}[scale=.6666667,draw=black,text=black]\draw[draw=black,line width=.5pt,fill=white] (398.000,1078.000)--(406.000,1083.000)--(406.000,1147.000)--(398.000,1131.000)--cycle;\end{scope}
\begin{scope}[scale=.6666667,draw=black,text=black]\node[inner sep=0pt,font=\normalsize,text=black] at (416.000,1115.000) {\shortstack{十\\位}};\end{scope}
\begin{scope}[scale=.6666667,draw=black,text=black]\node[inner sep=0pt,font=\normalsize,text=black] at (441.000,1115.000) {\shortstack{个\\位}};\end{scope}
\begin{scope}[scale=.6666667,draw=black,text=black]\node[inner sep=0pt,font=\normalsize,text=black] at (477.000,1115.000) {\shortstack{十\\分\\位}};\end{scope}
\begin{scope}[scale=.6666667,draw=black,text=black]\node[inner sep=0pt,font=\normalsize,text=black] at (502.000,1115.000) {\shortstack{百\\分\\位}};\end{scope}
\begin{scope}[scale=.6666667,draw=black,text=black]\node[inner sep=0pt,font=\normalsize,text=black] at (527.000,1115.000) {\shortstack{千\\分\\位}};\end{scope}
\begin{scope}[scale=.6666667,draw=black,text=black]\draw[line width=.5pt] (654,1015) rectangle (798,1079);\draw[line width=.5pt] (654,1019)--(798,1019);\end{scope}
\begin{scope}[scale=.6666667,draw=black,text=black]\draw[line width=.45pt] (675,1017)--(675,1079);\draw[line width=.45pt] (667,1015)..controls(676,995) and (681,1003)..(680,1015);\end{scope}
\begin{scope}[scale=.6666667,draw=black,text=black]\draw[line width=.45pt] (701,1017)--(701,1079);\draw[line width=.45pt] (693,1015)..controls(702,995) and (707,1003)..(706,1015);\end{scope}
\begin{scope}[scale=.6666667,draw=black,text=black]\draw[line width=.45pt] (731,1017)--(731,1079);\draw[line width=.45pt] (723,1015)..controls(732,995) and (737,1003)..(736,1015);\end{scope}
\begin{scope}[scale=.6666667,draw=black,text=black]\draw[line width=.45pt] (757,1017)--(757,1079);\draw[line width=.45pt] (749,1015)..controls(758,995) and (763,1003)..(762,1015);\end{scope}
\begin{scope}[scale=.6666667,draw=black,text=black]\draw[line width=.45pt] (781,1017)--(781,1079);\draw[line width=.45pt] (773,1015)..controls(782,995) and (787,1003)..(786,1015);\end{scope}
\begin{scope}[scale=.6666667,draw=black,text=black]\draw[draw=black,line width=.5pt,fill=white] (654.000,1078.000)--(662.000,1083.000)--(805.000,1083.000)--(798.000,1078.000)--cycle;\end{scope}
\begin{scope}[scale=.6666667,draw=black,text=black]\draw[draw=black,line width=.5pt,fill=white] (662.000,1083.000)--(805.000,1083.000)--(805.000,1147.000)--(662.000,1147.000)--cycle;\end{scope}
\begin{scope}[scale=.6666667,draw=black,text=black]\draw[draw=black,line width=.5pt,fill=white] (654.000,1078.000)--(662.000,1083.000)--(662.000,1147.000)--(654.000,1131.000)--cycle;\end{scope}
\begin{scope}[scale=.6666667,draw=black,text=black]\node[inner sep=0pt,font=\normalsize,text=black] at (675.000,1115.000) {\shortstack{十\\位}};\end{scope}
\begin{scope}[scale=.6666667,draw=black,text=black]\node[inner sep=0pt,font=\normalsize,text=black] at (701.000,1115.000) {\shortstack{个\\位}};\end{scope}
\begin{scope}[scale=.6666667,draw=black,text=black]\node[inner sep=0pt,font=\normalsize,text=black] at (731.000,1115.000) {\shortstack{十\\分\\位}};\end{scope}
\begin{scope}[scale=.6666667,draw=black,text=black]\node[inner sep=0pt,font=\normalsize,text=black] at (757.000,1115.000) {\shortstack{百\\分\\位}};\end{scope}
\begin{scope}[scale=.6666667,draw=black,text=black]\node[inner sep=0pt,font=\normalsize,text=black] at (781.000,1115.000) {\shortstack{千\\分\\位}};\end{scope}

\begin{scope}[color=answerred,draw=answerred,text=answerred]\ans{128.00}{226.67}{33.33}{10}{$\vphantom{0}0.4$}
\ans{294.00}{226.67}{33.33}{10}{$\vphantom{0}0.33$}
\ans{455.33}{225.33}{33.33}{10}{$\vphantom{0}0.203$}
\fill[answerred,opacity=.28](95.333,296.000) rectangle (181.733,394.000);
\fill[answerred,opacity=.28](259.333,296.000) rectangle (345.733,305.733);
\fill[answerred,opacity=.28](424.000,296.000) rectangle (472.667,393.333);
\fill[answerred,opacity=.28](472.667,296.000) rectangle (482.400,325.200);
\ans{228.67}{573.33}{48.00}{11}{$\vphantom{0}0.3$}
\ans{322.67}{534.00}{42.00}{11}{$\vphantom{0}\frac6{10}$}
\ans{321.33}{573.33}{48.00}{11}{$\vphantom{0}0.6$}
\filldraw[answerred,fill=answerred!65] (149.333,714.000) ellipse [x radius=4bp,y radius=2.2bp];
\filldraw[answerred,fill=answerred!65] (149.333,708.667) ellipse [x radius=4bp,y radius=2.2bp];
\filldraw[answerred,fill=answerred!65] (149.333,703.333) ellipse [x radius=4bp,y radius=2.2bp];
\filldraw[answerred,fill=answerred!65] (164.667,714.000) ellipse [x radius=4bp,y radius=2.2bp];
\filldraw[answerred,fill=answerred!65] (164.667,708.667) ellipse [x radius=4bp,y radius=2.2bp];
\filldraw[answerred,fill=answerred!65] (164.667,703.333) ellipse [x radius=4bp,y radius=2.2bp];
\filldraw[answerred,fill=answerred!65] (164.667,698.000) ellipse [x radius=4bp,y radius=2.2bp];
\filldraw[answerred,fill=answerred!65] (164.667,692.667) ellipse [x radius=4bp,y radius=2.2bp];
\filldraw[answerred,fill=answerred!65] (294.000,714.000) ellipse [x radius=4bp,y radius=2.2bp];
\filldraw[answerred,fill=answerred!65] (294.000,708.667) ellipse [x radius=4bp,y radius=2.2bp];
\filldraw[answerred,fill=answerred!65] (294.000,703.333) ellipse [x radius=4bp,y radius=2.2bp];
\filldraw[answerred,fill=answerred!65] (294.000,698.000) ellipse [x radius=4bp,y radius=2.2bp];
\filldraw[answerred,fill=answerred!65] (318.000,714.000) ellipse [x radius=4bp,y radius=2.2bp];
\filldraw[answerred,fill=answerred!65] (318.000,708.667) ellipse [x radius=4bp,y radius=2.2bp];
\filldraw[answerred,fill=answerred!65] (318.000,703.333) ellipse [x radius=4bp,y radius=2.2bp];
\filldraw[answerred,fill=answerred!65] (487.333,714.000) ellipse [x radius=4bp,y radius=2.2bp];
\filldraw[answerred,fill=answerred!65] (487.333,708.667) ellipse [x radius=4bp,y radius=2.2bp];
\filldraw[answerred,fill=answerred!65] (487.333,703.333) ellipse [x radius=4bp,y radius=2.2bp];
\filldraw[answerred,fill=answerred!65] (487.333,698.000) ellipse [x radius=4bp,y radius=2.2bp];
\filldraw[answerred,fill=answerred!65] (487.333,692.667) ellipse [x radius=4bp,y radius=2.2bp];
\filldraw[answerred,fill=answerred!65] (487.333,687.333) ellipse [x radius=4bp,y radius=2.2bp];
\filldraw[answerred,fill=answerred!65] (487.333,682.000) ellipse [x radius=4bp,y radius=2.2bp];
\filldraw[answerred,fill=answerred!65] (504.667,714.000) ellipse [x radius=4bp,y radius=2.2bp];
\filldraw[answerred,fill=answerred!65] (504.667,708.667) ellipse [x radius=4bp,y radius=2.2bp];
\filldraw[answerred,fill=answerred!65] (504.667,703.333) ellipse [x radius=4bp,y radius=2.2bp];
\filldraw[answerred,fill=answerred!65] (520.667,714.000) ellipse [x radius=4bp,y radius=2.2bp];
\filldraw[answerred,fill=answerred!65] (520.667,708.667) ellipse [x radius=4bp,y radius=2.2bp];
\end{scope}
\end{scope}\pgfresetboundingbox\path[use as bounding box] (56.0,113.0) rectangle (518.0,261.33);
\end{tikzpicture}}\par
\par\noindent




\par





\par\noindent

\QHeading{第2题　给图形涂色表示它下面的小数。}

\par\noindent\makebox[\linewidth][c]{
\begin{tikzpicture}[x=1bp,y=-1bp,color=black,draw=black,text=black]
\begin{scope}\clip (56.0,288.0) rectangle (527.0,436.0);
\hdhead{46}{小数的意义与性质——信息窗1}
\hdsection{1}{135}{用小数表示下图的涂色部分。}
\begin{scope}[scale=.6666667,draw=black,text=black]\path[pattern=north east lines,pattern color=black] (143,178) rectangle (200.6,323);\end{scope}
\begin{scope}[scale=.6666667,draw=black,text=black]\draw[draw=black,line width=.4pt] (143,178) rectangle (287,323);\foreach \i in {1,...,9}{\draw[draw=black,line width=.4pt] ({143+14.4*\i},178)--({143+14.4*\i},323);}\foreach \j in {1,...,0}{\draw[draw=black,line width=.4pt] (143,{178+145.0*\j})--(287,{178+145.0*\j});}\end{scope}
\begin{scope}[scale=.6666667,draw=black,text=black]\path[pattern=north east lines,pattern color=black] (389,178) rectangle (432.2,322);\end{scope}
\begin{scope}[scale=.6666667,draw=black,text=black]\path[pattern=north east lines,pattern color=black] (432.2,278.8) rectangle (446.59999999999997,322.0);\end{scope}
\begin{scope}[scale=.6666667,draw=black,text=black]\draw[draw=black,line width=.35pt] (389,178) rectangle (533,322);\foreach \i in {1,...,9}{\draw[draw=black,line width=.35pt] ({389+14.4*\i},178)--({389+14.4*\i},322);}\foreach \j in {1,...,9}{\draw[draw=black,line width=.35pt] (389,{178+14.4*\j})--(533,{178+14.4*\j});}\end{scope}
\begin{scope}[scale=.6666667,draw=black,text=black]\fill[white] (636.0000,319.0000)--(736.0000,319.0000)--(736.0000,219.0000)--(636.0000,219.0000)--cycle;\end{scope}
\begin{scope}[scale=.6666667,draw=black,text=black]\fill[white] (636.0000,219.0000)--(736.0000,219.0000)--(769.0000,179.0000)--(669.0000,179.0000)--cycle;\end{scope}
\begin{scope}[scale=.6666667,draw=black,text=black]\path[pattern=north east lines,pattern color=black] (736.0000,319.0000)--(769.0000,279.0000)--(769.0000,179.0000)--(736.0000,219.0000)--cycle;\end{scope}
\begin{scope}[scale=.6666667,draw=black,text=black]\path[pattern=north east lines,pattern color=black] (706.0000,249.0000)--(716.0000,249.0000)--(716.0000,239.0000)--(706.0000,239.0000)--cycle;\end{scope}
\begin{scope}[scale=.6666667,draw=black,text=black]\path[pattern=north east lines,pattern color=black] (706.0000,239.0000)--(716.0000,239.0000)--(716.0000,229.0000)--(706.0000,229.0000)--cycle;\end{scope}
\begin{scope}[scale=.6666667,draw=black,text=black]\path[pattern=north east lines,pattern color=black] (706.0000,229.0000)--(716.0000,229.0000)--(716.0000,219.0000)--(706.0000,219.0000)--cycle;\end{scope}
\begin{scope}[scale=.6666667,draw=black,text=black]\path[pattern=north east lines,pattern color=black] (706.0000,219.0000)--(716.0000,219.0000)--(719.3000,215.0000)--(709.3000,215.0000)--cycle;\end{scope}
\begin{scope}[scale=.6666667,draw=black,text=black]\path[pattern=north east lines,pattern color=black] (716.0000,319.0000)--(726.0000,319.0000)--(726.0000,309.0000)--(716.0000,309.0000)--cycle;\end{scope}
\begin{scope}[scale=.6666667,draw=black,text=black]\path[pattern=north east lines,pattern color=black] (716.0000,309.0000)--(726.0000,309.0000)--(726.0000,299.0000)--(716.0000,299.0000)--cycle;\end{scope}
\begin{scope}[scale=.6666667,draw=black,text=black]\path[pattern=north east lines,pattern color=black] (716.0000,299.0000)--(726.0000,299.0000)--(726.0000,289.0000)--(716.0000,289.0000)--cycle;\end{scope}
\begin{scope}[scale=.6666667,draw=black,text=black]\path[pattern=north east lines,pattern color=black] (716.0000,289.0000)--(726.0000,289.0000)--(726.0000,279.0000)--(716.0000,279.0000)--cycle;\end{scope}
\begin{scope}[scale=.6666667,draw=black,text=black]\path[pattern=north east lines,pattern color=black] (716.0000,279.0000)--(726.0000,279.0000)--(726.0000,269.0000)--(716.0000,269.0000)--cycle;\end{scope}
\begin{scope}[scale=.6666667,draw=black,text=black]\path[pattern=north east lines,pattern color=black] (716.0000,269.0000)--(726.0000,269.0000)--(726.0000,259.0000)--(716.0000,259.0000)--cycle;\end{scope}
\begin{scope}[scale=.6666667,draw=black,text=black]\path[pattern=north east lines,pattern color=black] (716.0000,259.0000)--(726.0000,259.0000)--(726.0000,249.0000)--(716.0000,249.0000)--cycle;\end{scope}
\begin{scope}[scale=.6666667,draw=black,text=black]\path[pattern=north east lines,pattern color=black] (716.0000,249.0000)--(726.0000,249.0000)--(726.0000,239.0000)--(716.0000,239.0000)--cycle;\end{scope}
\begin{scope}[scale=.6666667,draw=black,text=black]\path[pattern=north east lines,pattern color=black] (716.0000,239.0000)--(726.0000,239.0000)--(726.0000,229.0000)--(716.0000,229.0000)--cycle;\end{scope}
\begin{scope}[scale=.6666667,draw=black,text=black]\path[pattern=north east lines,pattern color=black] (716.0000,229.0000)--(726.0000,229.0000)--(726.0000,219.0000)--(716.0000,219.0000)--cycle;\end{scope}
\begin{scope}[scale=.6666667,draw=black,text=black]\path[pattern=north east lines,pattern color=black] (716.0000,219.0000)--(726.0000,219.0000)--(729.3000,215.0000)--(719.3000,215.0000)--cycle;\end{scope}
\begin{scope}[scale=.6666667,draw=black,text=black]\path[pattern=north east lines,pattern color=black] (719.3000,215.0000)--(729.3000,215.0000)--(732.6000,211.0000)--(722.6000,211.0000)--cycle;\end{scope}
\begin{scope}[scale=.6666667,draw=black,text=black]\path[pattern=north east lines,pattern color=black] (722.6000,211.0000)--(732.6000,211.0000)--(735.9000,207.0000)--(725.9000,207.0000)--cycle;\end{scope}
\begin{scope}[scale=.6666667,draw=black,text=black]\path[pattern=north east lines,pattern color=black] (725.9000,207.0000)--(735.9000,207.0000)--(739.2000,203.0000)--(729.2000,203.0000)--cycle;\end{scope}
\begin{scope}[scale=.6666667,draw=black,text=black]\path[pattern=north east lines,pattern color=black] (729.2000,203.0000)--(739.2000,203.0000)--(742.5000,199.0000)--(732.5000,199.0000)--cycle;\end{scope}
\begin{scope}[scale=.6666667,draw=black,text=black]\path[pattern=north east lines,pattern color=black] (732.5000,199.0000)--(742.5000,199.0000)--(745.8000,195.0000)--(735.8000,195.0000)--cycle;\end{scope}
\begin{scope}[scale=.6666667,draw=black,text=black]\path[pattern=north east lines,pattern color=black] (735.8000,195.0000)--(745.8000,195.0000)--(749.1000,191.0000)--(739.1000,191.0000)--cycle;\end{scope}
\begin{scope}[scale=.6666667,draw=black,text=black]\path[pattern=north east lines,pattern color=black] (739.1000,191.0000)--(749.1000,191.0000)--(752.4000,187.0000)--(742.4000,187.0000)--cycle;\end{scope}
\begin{scope}[scale=.6666667,draw=black,text=black]\path[pattern=north east lines,pattern color=black] (742.4000,187.0000)--(752.4000,187.0000)--(755.7000,183.0000)--(745.7000,183.0000)--cycle;\end{scope}
\begin{scope}[scale=.6666667,draw=black,text=black]\path[pattern=north east lines,pattern color=black] (745.7000,183.0000)--(755.7000,183.0000)--(759.0000,179.0000)--(749.0000,179.0000)--cycle;\end{scope}
\begin{scope}[scale=.6666667,draw=black,text=black]\path[pattern=north east lines,pattern color=black] (726.0000,319.0000)--(736.0000,319.0000)--(736.0000,309.0000)--(726.0000,309.0000)--cycle;\end{scope}
\begin{scope}[scale=.6666667,draw=black,text=black]\path[pattern=north east lines,pattern color=black] (726.0000,309.0000)--(736.0000,309.0000)--(736.0000,299.0000)--(726.0000,299.0000)--cycle;\end{scope}
\begin{scope}[scale=.6666667,draw=black,text=black]\path[pattern=north east lines,pattern color=black] (726.0000,299.0000)--(736.0000,299.0000)--(736.0000,289.0000)--(726.0000,289.0000)--cycle;\end{scope}
\begin{scope}[scale=.6666667,draw=black,text=black]\path[pattern=north east lines,pattern color=black] (726.0000,289.0000)--(736.0000,289.0000)--(736.0000,279.0000)--(726.0000,279.0000)--cycle;\end{scope}
\begin{scope}[scale=.6666667,draw=black,text=black]\path[pattern=north east lines,pattern color=black] (726.0000,279.0000)--(736.0000,279.0000)--(736.0000,269.0000)--(726.0000,269.0000)--cycle;\end{scope}
\begin{scope}[scale=.6666667,draw=black,text=black]\path[pattern=north east lines,pattern color=black] (726.0000,269.0000)--(736.0000,269.0000)--(736.0000,259.0000)--(726.0000,259.0000)--cycle;\end{scope}
\begin{scope}[scale=.6666667,draw=black,text=black]\path[pattern=north east lines,pattern color=black] (726.0000,259.0000)--(736.0000,259.0000)--(736.0000,249.0000)--(726.0000,249.0000)--cycle;\end{scope}
\begin{scope}[scale=.6666667,draw=black,text=black]\path[pattern=north east lines,pattern color=black] (726.0000,249.0000)--(736.0000,249.0000)--(736.0000,239.0000)--(726.0000,239.0000)--cycle;\end{scope}
\begin{scope}[scale=.6666667,draw=black,text=black]\path[pattern=north east lines,pattern color=black] (726.0000,239.0000)--(736.0000,239.0000)--(736.0000,229.0000)--(726.0000,229.0000)--cycle;\end{scope}
\begin{scope}[scale=.6666667,draw=black,text=black]\path[pattern=north east lines,pattern color=black] (726.0000,229.0000)--(736.0000,229.0000)--(736.0000,219.0000)--(726.0000,219.0000)--cycle;\end{scope}
\begin{scope}[scale=.6666667,draw=black,text=black]\path[pattern=north east lines,pattern color=black] (726.0000,219.0000)--(736.0000,219.0000)--(739.3000,215.0000)--(729.3000,215.0000)--cycle;\end{scope}
\begin{scope}[scale=.6666667,draw=black,text=black]\path[pattern=north east lines,pattern color=black] (729.3000,215.0000)--(739.3000,215.0000)--(742.6000,211.0000)--(732.6000,211.0000)--cycle;\end{scope}
\begin{scope}[scale=.6666667,draw=black,text=black]\path[pattern=north east lines,pattern color=black] (732.6000,211.0000)--(742.6000,211.0000)--(745.9000,207.0000)--(735.9000,207.0000)--cycle;\end{scope}
\begin{scope}[scale=.6666667,draw=black,text=black]\path[pattern=north east lines,pattern color=black] (735.9000,207.0000)--(745.9000,207.0000)--(749.2000,203.0000)--(739.2000,203.0000)--cycle;\end{scope}
\begin{scope}[scale=.6666667,draw=black,text=black]\path[pattern=north east lines,pattern color=black] (739.2000,203.0000)--(749.2000,203.0000)--(752.5000,199.0000)--(742.5000,199.0000)--cycle;\end{scope}
\begin{scope}[scale=.6666667,draw=black,text=black]\path[pattern=north east lines,pattern color=black] (742.5000,199.0000)--(752.5000,199.0000)--(755.8000,195.0000)--(745.8000,195.0000)--cycle;\end{scope}
\begin{scope}[scale=.6666667,draw=black,text=black]\path[pattern=north east lines,pattern color=black] (745.8000,195.0000)--(755.8000,195.0000)--(759.1000,191.0000)--(749.1000,191.0000)--cycle;\end{scope}
\begin{scope}[scale=.6666667,draw=black,text=black]\path[pattern=north east lines,pattern color=black] (749.1000,191.0000)--(759.1000,191.0000)--(762.4000,187.0000)--(752.4000,187.0000)--cycle;\end{scope}
\begin{scope}[scale=.6666667,draw=black,text=black]\path[pattern=north east lines,pattern color=black] (752.4000,187.0000)--(762.4000,187.0000)--(765.7000,183.0000)--(755.7000,183.0000)--cycle;\end{scope}
\begin{scope}[scale=.6666667,draw=black,text=black]\path[pattern=north east lines,pattern color=black] (755.7000,183.0000)--(765.7000,183.0000)--(769.0000,179.0000)--(759.0000,179.0000)--cycle;\end{scope}
\begin{scope}[scale=.6666667,draw=black,text=black]\draw[draw=black!85,line width=.32pt] (636.000,319.000)--(636.000,219.000);\end{scope}
\begin{scope}[scale=.6666667,draw=black,text=black]\draw[draw=black!85,line width=.32pt] (636.000,319.000)--(736.000,319.000);\end{scope}
\begin{scope}[scale=.6666667,draw=black,text=black]\draw[draw=black!85,line width=.32pt] (636.000,219.000)--(669.000,179.000);\end{scope}
\begin{scope}[scale=.6666667,draw=black,text=black]\draw[draw=black!85,line width=.32pt] (636.000,219.000)--(736.000,219.000);\end{scope}
\begin{scope}[scale=.6666667,draw=black,text=black]\draw[draw=black!85,line width=.32pt] (736.000,319.000)--(736.000,219.000);\end{scope}
\begin{scope}[scale=.6666667,draw=black,text=black]\draw[draw=black!85,line width=.32pt] (736.000,319.000)--(769.000,279.000);\end{scope}
\begin{scope}[scale=.6666667,draw=black,text=black]\draw[draw=black!85,line width=.32pt] (646.000,319.000)--(646.000,219.000);\end{scope}
\begin{scope}[scale=.6666667,draw=black,text=black]\draw[draw=black!85,line width=.32pt] (636.000,309.000)--(736.000,309.000);\end{scope}
\begin{scope}[scale=.6666667,draw=black,text=black]\draw[draw=black!85,line width=.32pt] (646.000,219.000)--(679.000,179.000);\end{scope}
\begin{scope}[scale=.6666667,draw=black,text=black]\draw[draw=black!85,line width=.32pt] (639.300,215.000)--(739.300,215.000);\end{scope}
\begin{scope}[scale=.6666667,draw=black,text=black]\draw[draw=black!85,line width=.32pt] (739.300,315.000)--(739.300,215.000);\end{scope}
\begin{scope}[scale=.6666667,draw=black,text=black]\draw[draw=black!85,line width=.32pt] (736.000,309.000)--(769.000,269.000);\end{scope}
\begin{scope}[scale=.6666667,draw=black,text=black]\draw[draw=black!85,line width=.32pt] (656.000,319.000)--(656.000,219.000);\end{scope}
\begin{scope}[scale=.6666667,draw=black,text=black]\draw[draw=black!85,line width=.32pt] (636.000,299.000)--(736.000,299.000);\end{scope}
\begin{scope}[scale=.6666667,draw=black,text=black]\draw[draw=black!85,line width=.32pt] (656.000,219.000)--(689.000,179.000);\end{scope}
\begin{scope}[scale=.6666667,draw=black,text=black]\draw[draw=black!85,line width=.32pt] (642.600,211.000)--(742.600,211.000);\end{scope}
\begin{scope}[scale=.6666667,draw=black,text=black]\draw[draw=black!85,line width=.32pt] (742.600,311.000)--(742.600,211.000);\end{scope}
\begin{scope}[scale=.6666667,draw=black,text=black]\draw[draw=black!85,line width=.32pt] (736.000,299.000)--(769.000,259.000);\end{scope}
\begin{scope}[scale=.6666667,draw=black,text=black]\draw[draw=black!85,line width=.32pt] (666.000,319.000)--(666.000,219.000);\end{scope}
\begin{scope}[scale=.6666667,draw=black,text=black]\draw[draw=black!85,line width=.32pt] (636.000,289.000)--(736.000,289.000);\end{scope}
\begin{scope}[scale=.6666667,draw=black,text=black]\draw[draw=black!85,line width=.32pt] (666.000,219.000)--(699.000,179.000);\end{scope}
\begin{scope}[scale=.6666667,draw=black,text=black]\draw[draw=black!85,line width=.32pt] (645.900,207.000)--(745.900,207.000);\end{scope}
\begin{scope}[scale=.6666667,draw=black,text=black]\draw[draw=black!85,line width=.32pt] (745.900,307.000)--(745.900,207.000);\end{scope}
\begin{scope}[scale=.6666667,draw=black,text=black]\draw[draw=black!85,line width=.32pt] (736.000,289.000)--(769.000,249.000);\end{scope}
\begin{scope}[scale=.6666667,draw=black,text=black]\draw[draw=black!85,line width=.32pt] (676.000,319.000)--(676.000,219.000);\end{scope}
\begin{scope}[scale=.6666667,draw=black,text=black]\draw[draw=black!85,line width=.32pt] (636.000,279.000)--(736.000,279.000);\end{scope}
\begin{scope}[scale=.6666667,draw=black,text=black]\draw[draw=black!85,line width=.32pt] (676.000,219.000)--(709.000,179.000);\end{scope}
\begin{scope}[scale=.6666667,draw=black,text=black]\draw[draw=black!85,line width=.32pt] (649.200,203.000)--(749.200,203.000);\end{scope}
\begin{scope}[scale=.6666667,draw=black,text=black]\draw[draw=black!85,line width=.32pt] (749.200,303.000)--(749.200,203.000);\end{scope}
\begin{scope}[scale=.6666667,draw=black,text=black]\draw[draw=black!85,line width=.32pt] (736.000,279.000)--(769.000,239.000);\end{scope}
\begin{scope}[scale=.6666667,draw=black,text=black]\draw[draw=black!85,line width=.32pt] (686.000,319.000)--(686.000,219.000);\end{scope}
\begin{scope}[scale=.6666667,draw=black,text=black]\draw[draw=black!85,line width=.32pt] (636.000,269.000)--(736.000,269.000);\end{scope}
\begin{scope}[scale=.6666667,draw=black,text=black]\draw[draw=black!85,line width=.32pt] (686.000,219.000)--(719.000,179.000);\end{scope}
\begin{scope}[scale=.6666667,draw=black,text=black]\draw[draw=black!85,line width=.32pt] (652.500,199.000)--(752.500,199.000);\end{scope}
\begin{scope}[scale=.6666667,draw=black,text=black]\draw[draw=black!85,line width=.32pt] (752.500,299.000)--(752.500,199.000);\end{scope}
\begin{scope}[scale=.6666667,draw=black,text=black]\draw[draw=black!85,line width=.32pt] (736.000,269.000)--(769.000,229.000);\end{scope}
\begin{scope}[scale=.6666667,draw=black,text=black]\draw[draw=black!85,line width=.32pt] (696.000,319.000)--(696.000,219.000);\end{scope}
\begin{scope}[scale=.6666667,draw=black,text=black]\draw[draw=black!85,line width=.32pt] (636.000,259.000)--(736.000,259.000);\end{scope}
\begin{scope}[scale=.6666667,draw=black,text=black]\draw[draw=black!85,line width=.32pt] (696.000,219.000)--(729.000,179.000);\end{scope}
\begin{scope}[scale=.6666667,draw=black,text=black]\draw[draw=black!85,line width=.32pt] (655.800,195.000)--(755.800,195.000);\end{scope}
\begin{scope}[scale=.6666667,draw=black,text=black]\draw[draw=black!85,line width=.32pt] (755.800,295.000)--(755.800,195.000);\end{scope}
\begin{scope}[scale=.6666667,draw=black,text=black]\draw[draw=black!85,line width=.32pt] (736.000,259.000)--(769.000,219.000);\end{scope}
\begin{scope}[scale=.6666667,draw=black,text=black]\draw[draw=black!85,line width=.32pt] (706.000,319.000)--(706.000,219.000);\end{scope}
\begin{scope}[scale=.6666667,draw=black,text=black]\draw[draw=black!85,line width=.32pt] (636.000,249.000)--(736.000,249.000);\end{scope}
\begin{scope}[scale=.6666667,draw=black,text=black]\draw[draw=black!85,line width=.32pt] (706.000,219.000)--(739.000,179.000);\end{scope}
\begin{scope}[scale=.6666667,draw=black,text=black]\draw[draw=black!85,line width=.32pt] (659.100,191.000)--(759.100,191.000);\end{scope}
\begin{scope}[scale=.6666667,draw=black,text=black]\draw[draw=black!85,line width=.32pt] (759.100,291.000)--(759.100,191.000);\end{scope}
\begin{scope}[scale=.6666667,draw=black,text=black]\draw[draw=black!85,line width=.32pt] (736.000,249.000)--(769.000,209.000);\end{scope}
\begin{scope}[scale=.6666667,draw=black,text=black]\draw[draw=black!85,line width=.32pt] (716.000,319.000)--(716.000,219.000);\end{scope}
\begin{scope}[scale=.6666667,draw=black,text=black]\draw[draw=black!85,line width=.32pt] (636.000,239.000)--(736.000,239.000);\end{scope}
\begin{scope}[scale=.6666667,draw=black,text=black]\draw[draw=black!85,line width=.32pt] (716.000,219.000)--(749.000,179.000);\end{scope}
\begin{scope}[scale=.6666667,draw=black,text=black]\draw[draw=black!85,line width=.32pt] (662.400,187.000)--(762.400,187.000);\end{scope}
\begin{scope}[scale=.6666667,draw=black,text=black]\draw[draw=black!85,line width=.32pt] (762.400,287.000)--(762.400,187.000);\end{scope}
\begin{scope}[scale=.6666667,draw=black,text=black]\draw[draw=black!85,line width=.32pt] (736.000,239.000)--(769.000,199.000);\end{scope}
\begin{scope}[scale=.6666667,draw=black,text=black]\draw[draw=black!85,line width=.32pt] (726.000,319.000)--(726.000,219.000);\end{scope}
\begin{scope}[scale=.6666667,draw=black,text=black]\draw[draw=black!85,line width=.32pt] (636.000,229.000)--(736.000,229.000);\end{scope}
\begin{scope}[scale=.6666667,draw=black,text=black]\draw[draw=black!85,line width=.32pt] (726.000,219.000)--(759.000,179.000);\end{scope}
\begin{scope}[scale=.6666667,draw=black,text=black]\draw[draw=black!85,line width=.32pt] (665.700,183.000)--(765.700,183.000);\end{scope}
\begin{scope}[scale=.6666667,draw=black,text=black]\draw[draw=black!85,line width=.32pt] (765.700,283.000)--(765.700,183.000);\end{scope}
\begin{scope}[scale=.6666667,draw=black,text=black]\draw[draw=black!85,line width=.32pt] (736.000,229.000)--(769.000,189.000);\end{scope}
\begin{scope}[scale=.6666667,draw=black,text=black]\draw[draw=black!85,line width=.32pt] (736.000,319.000)--(736.000,219.000);\end{scope}
\begin{scope}[scale=.6666667,draw=black,text=black]\draw[draw=black!85,line width=.32pt] (636.000,219.000)--(736.000,219.000);\end{scope}
\begin{scope}[scale=.6666667,draw=black,text=black]\draw[draw=black!85,line width=.32pt] (736.000,219.000)--(769.000,179.000);\end{scope}
\begin{scope}[scale=.6666667,draw=black,text=black]\draw[draw=black!85,line width=.32pt] (669.000,179.000)--(769.000,179.000);\end{scope}
\begin{scope}[scale=.6666667,draw=black,text=black]\draw[draw=black!85,line width=.32pt] (769.000,279.000)--(769.000,179.000);\end{scope}
\begin{scope}[scale=.6666667,draw=black,text=black]\draw[draw=black!85,line width=.32pt] (736.000,219.000)--(769.000,179.000);\end{scope}
\begin{scope}[scale=.6666667,draw=black,text=black]\draw[draw=black,line width=.55pt] (636.000,319.000)--(736.000,319.000)--(736.000,219.000)--(636.000,219.000)--cycle;\end{scope}
\begin{scope}[scale=.6666667,draw=black,text=black]\draw[draw=black,line width=.55pt] (636.000,219.000)--(736.000,219.000)--(769.000,179.000)--(669.000,179.000)--cycle;\end{scope}
\begin{scope}[scale=.6666667,draw=black,text=black]\draw[draw=black,line width=.55pt] (736.000,319.000)--(769.000,279.000)--(769.000,179.000)--(736.000,219.000)--cycle;\end{scope}
\begin{scope}[scale=.6666667,draw=black,text=black]\node[inner sep=0pt,font=\normalsize,text=black] at (123.000,189.000) {①};\end{scope}
\begin{scope}[scale=.6666667,draw=black,text=black]\node[inner sep=0pt,font=\normalsize,text=black] at (370.000,189.000) {②};\end{scope}
\begin{scope}[scale=.6666667,draw=black,text=black]\node[inner sep=0pt,font=\normalsize,text=black] at (617.000,190.000) {③};\end{scope}
\q{171}{339}{28}{13}{（}
\q{247}{339}{25}{13}{）}
\q{418}{339}{28}{13}{（}
\q{493}{339}{25}{13}{）}
\q{657}{338}{28}{13}{（}
\q{733}{338}{25}{13}{）}
\hdsection{2}{395}{给图形涂色表示它下面的小数。}
\begin{scope}[scale=.6666667,draw=black,text=black]\draw[draw=black,line width=.4pt] (143,444) rectangle (287,591);\foreach \i in {1,...,9}{\draw[draw=black,line width=.4pt] ({143+14.4*\i},444)--({143+14.4*\i},591);}\foreach \j in {1,...,0}{\draw[draw=black,line width=.4pt] (143,{444+147.0*\j})--(287,{444+147.0*\j});}\end{scope}
\begin{scope}[scale=.6666667,draw=black,text=black]\draw[draw=black,line width=.35pt] (389,444) rectangle (533,590);\foreach \i in {1,...,9}{\draw[draw=black,line width=.35pt] ({389+14.4*\i},444)--({389+14.4*\i},590);}\foreach \j in {1,...,9}{\draw[draw=black,line width=.35pt] (389,{444+14.6*\j})--(533,{444+14.6*\j});}\end{scope}
\begin{scope}[scale=.6666667,draw=black,text=black]\draw[draw=black,line width=.35pt] (636,444) rectangle (782,590);\foreach \i in {1,...,9}{\draw[draw=black,line width=.35pt] ({636+14.6*\i},444)--({636+14.6*\i},590);}\foreach \j in {1,...,9}{\draw[draw=black,line width=.35pt] (636,{444+14.6*\j})--(782,{444+14.6*\j});}\end{scope}
\begin{scope}[scale=.6666667,draw=black,text=black]\node[inner sep=0pt,font=\normalsize,text=black] at (123.000,451.000) {①};\end{scope}
\begin{scope}[scale=.6666667,draw=black,text=black]\node[inner sep=0pt,font=\normalsize,text=black] at (369.000,451.000) {②};\end{scope}
\begin{scope}[scale=.6666667,draw=black,text=black]\node[inner sep=0pt,font=\normalsize,text=black] at (616.000,451.000) {③};\end{scope}
\begin{scope}[scale=.6666667,draw=black,text=black]\node[inner sep=0pt,font=\normalsize,text=black] at (215.000,612.000) {$0.9$};\end{scope}
\begin{scope}[scale=.6666667,draw=black,text=black]\node[inner sep=0pt,font=\normalsize,text=black] at (461.000,612.000) {$0.09$};\end{scope}
\begin{scope}[scale=.6666667,draw=black,text=black]\node[inner sep=0pt,font=\normalsize,text=black] at (709.000,612.000) {$0.53$};\end{scope}
\hdsection{3}{657}{照样子填一填。}
\begin{scope}[scale=.6666667,draw=black,text=black]\draw[draw=black,line width=.55pt] (210.000,705.000)--(729.000,705.000)--(724.000,721.000)--(733.000,734.000)--(729.000,754.000)--(210.000,754.000)--cycle;\end{scope}
\begin{scope}[scale=.6666667,draw=black,text=black]\draw[draw=black,line width=.35pt] (219.000,705.000)--(219.000,719.000);\end{scope}
\begin{scope}[scale=.6666667,draw=black,text=black]\draw[draw=black,line width=.35pt] (223.767,705.000)--(223.767,711.000);\end{scope}
\begin{scope}[scale=.6666667,draw=black,text=black]\draw[draw=black,line width=.35pt] (228.533,705.000)--(228.533,711.000);\end{scope}
\begin{scope}[scale=.6666667,draw=black,text=black]\draw[draw=black,line width=.35pt] (233.300,705.000)--(233.300,711.000);\end{scope}
\begin{scope}[scale=.6666667,draw=black,text=black]\draw[draw=black,line width=.35pt] (238.067,705.000)--(238.067,711.000);\end{scope}
\begin{scope}[scale=.6666667,draw=black,text=black]\draw[draw=black,line width=.35pt] (242.833,705.000)--(242.833,714.000);\end{scope}
\begin{scope}[scale=.6666667,draw=black,text=black]\draw[draw=black,line width=.35pt] (247.600,705.000)--(247.600,711.000);\end{scope}
\begin{scope}[scale=.6666667,draw=black,text=black]\draw[draw=black,line width=.35pt] (252.367,705.000)--(252.367,711.000);\end{scope}
\begin{scope}[scale=.6666667,draw=black,text=black]\draw[draw=black,line width=.35pt] (257.133,705.000)--(257.133,711.000);\end{scope}
\begin{scope}[scale=.6666667,draw=black,text=black]\draw[draw=black,line width=.35pt] (261.900,705.000)--(261.900,711.000);\end{scope}
\begin{scope}[scale=.6666667,draw=black,text=black]\draw[draw=black,line width=.35pt] (266.667,705.000)--(266.667,719.000);\end{scope}
\begin{scope}[scale=.6666667,draw=black,text=black]\draw[draw=black,line width=.35pt] (271.433,705.000)--(271.433,711.000);\end{scope}
\begin{scope}[scale=.6666667,draw=black,text=black]\draw[draw=black,line width=.35pt] (276.200,705.000)--(276.200,711.000);\end{scope}
\begin{scope}[scale=.6666667,draw=black,text=black]\draw[draw=black,line width=.35pt] (280.967,705.000)--(280.967,711.000);\end{scope}
\begin{scope}[scale=.6666667,draw=black,text=black]\draw[draw=black,line width=.35pt] (285.733,705.000)--(285.733,711.000);\end{scope}
\begin{scope}[scale=.6666667,draw=black,text=black]\draw[draw=black,line width=.35pt] (290.500,705.000)--(290.500,714.000);\end{scope}
\begin{scope}[scale=.6666667,draw=black,text=black]\draw[draw=black,line width=.35pt] (295.267,705.000)--(295.267,711.000);\end{scope}
\begin{scope}[scale=.6666667,draw=black,text=black]\draw[draw=black,line width=.35pt] (300.033,705.000)--(300.033,711.000);\end{scope}
\begin{scope}[scale=.6666667,draw=black,text=black]\draw[draw=black,line width=.35pt] (304.800,705.000)--(304.800,711.000);\end{scope}
\begin{scope}[scale=.6666667,draw=black,text=black]\draw[draw=black,line width=.35pt] (309.567,705.000)--(309.567,711.000);\end{scope}
\begin{scope}[scale=.6666667,draw=black,text=black]\draw[draw=black,line width=.35pt] (314.333,705.000)--(314.333,719.000);\end{scope}
\begin{scope}[scale=.6666667,draw=black,text=black]\draw[draw=black,line width=.35pt] (319.100,705.000)--(319.100,711.000);\end{scope}
\begin{scope}[scale=.6666667,draw=black,text=black]\draw[draw=black,line width=.35pt] (323.867,705.000)--(323.867,711.000);\end{scope}
\begin{scope}[scale=.6666667,draw=black,text=black]\draw[draw=black,line width=.35pt] (328.633,705.000)--(328.633,711.000);\end{scope}
\begin{scope}[scale=.6666667,draw=black,text=black]\draw[draw=black,line width=.35pt] (333.400,705.000)--(333.400,711.000);\end{scope}
\begin{scope}[scale=.6666667,draw=black,text=black]\draw[draw=black,line width=.35pt] (338.167,705.000)--(338.167,714.000);\end{scope}
\begin{scope}[scale=.6666667,draw=black,text=black]\draw[draw=black,line width=.35pt] (342.933,705.000)--(342.933,711.000);\end{scope}
\begin{scope}[scale=.6666667,draw=black,text=black]\draw[draw=black,line width=.35pt] (347.700,705.000)--(347.700,711.000);\end{scope}
\begin{scope}[scale=.6666667,draw=black,text=black]\draw[draw=black,line width=.35pt] (352.467,705.000)--(352.467,711.000);\end{scope}
\begin{scope}[scale=.6666667,draw=black,text=black]\draw[draw=black,line width=.35pt] (357.233,705.000)--(357.233,711.000);\end{scope}
\begin{scope}[scale=.6666667,draw=black,text=black]\draw[draw=black,line width=.35pt] (362.000,705.000)--(362.000,719.000);\end{scope}
\begin{scope}[scale=.6666667,draw=black,text=black]\draw[draw=black,line width=.35pt] (366.767,705.000)--(366.767,711.000);\end{scope}
\begin{scope}[scale=.6666667,draw=black,text=black]\draw[draw=black,line width=.35pt] (371.533,705.000)--(371.533,711.000);\end{scope}
\begin{scope}[scale=.6666667,draw=black,text=black]\draw[draw=black,line width=.35pt] (376.300,705.000)--(376.300,711.000);\end{scope}
\begin{scope}[scale=.6666667,draw=black,text=black]\draw[draw=black,line width=.35pt] (381.067,705.000)--(381.067,711.000);\end{scope}
\begin{scope}[scale=.6666667,draw=black,text=black]\draw[draw=black,line width=.35pt] (385.833,705.000)--(385.833,714.000);\end{scope}
\begin{scope}[scale=.6666667,draw=black,text=black]\draw[draw=black,line width=.35pt] (390.600,705.000)--(390.600,711.000);\end{scope}
\begin{scope}[scale=.6666667,draw=black,text=black]\draw[draw=black,line width=.35pt] (395.367,705.000)--(395.367,711.000);\end{scope}
\begin{scope}[scale=.6666667,draw=black,text=black]\draw[draw=black,line width=.35pt] (400.133,705.000)--(400.133,711.000);\end{scope}
\begin{scope}[scale=.6666667,draw=black,text=black]\draw[draw=black,line width=.35pt] (404.900,705.000)--(404.900,711.000);\end{scope}
\begin{scope}[scale=.6666667,draw=black,text=black]\draw[draw=black,line width=.35pt] (409.667,705.000)--(409.667,719.000);\end{scope}
\begin{scope}[scale=.6666667,draw=black,text=black]\draw[draw=black,line width=.35pt] (414.433,705.000)--(414.433,711.000);\end{scope}
\begin{scope}[scale=.6666667,draw=black,text=black]\draw[draw=black,line width=.35pt] (419.200,705.000)--(419.200,711.000);\end{scope}
\begin{scope}[scale=.6666667,draw=black,text=black]\draw[draw=black,line width=.35pt] (423.967,705.000)--(423.967,711.000);\end{scope}
\begin{scope}[scale=.6666667,draw=black,text=black]\draw[draw=black,line width=.35pt] (428.733,705.000)--(428.733,711.000);\end{scope}
\begin{scope}[scale=.6666667,draw=black,text=black]\draw[draw=black,line width=.35pt] (433.500,705.000)--(433.500,714.000);\end{scope}
\begin{scope}[scale=.6666667,draw=black,text=black]\draw[draw=black,line width=.35pt] (438.267,705.000)--(438.267,711.000);\end{scope}
\begin{scope}[scale=.6666667,draw=black,text=black]\draw[draw=black,line width=.35pt] (443.033,705.000)--(443.033,711.000);\end{scope}
\begin{scope}[scale=.6666667,draw=black,text=black]\draw[draw=black,line width=.35pt] (447.800,705.000)--(447.800,711.000);\end{scope}
\begin{scope}[scale=.6666667,draw=black,text=black]\draw[draw=black,line width=.35pt] (452.567,705.000)--(452.567,711.000);\end{scope}
\begin{scope}[scale=.6666667,draw=black,text=black]\draw[draw=black,line width=.35pt] (457.333,705.000)--(457.333,719.000);\end{scope}
\begin{scope}[scale=.6666667,draw=black,text=black]\draw[draw=black,line width=.35pt] (462.100,705.000)--(462.100,711.000);\end{scope}
\begin{scope}[scale=.6666667,draw=black,text=black]\draw[draw=black,line width=.35pt] (466.867,705.000)--(466.867,711.000);\end{scope}
\begin{scope}[scale=.6666667,draw=black,text=black]\draw[draw=black,line width=.35pt] (471.633,705.000)--(471.633,711.000);\end{scope}
\begin{scope}[scale=.6666667,draw=black,text=black]\draw[draw=black,line width=.35pt] (476.400,705.000)--(476.400,711.000);\end{scope}
\begin{scope}[scale=.6666667,draw=black,text=black]\draw[draw=black,line width=.35pt] (481.167,705.000)--(481.167,714.000);\end{scope}
\begin{scope}[scale=.6666667,draw=black,text=black]\draw[draw=black,line width=.35pt] (485.933,705.000)--(485.933,711.000);\end{scope}
\begin{scope}[scale=.6666667,draw=black,text=black]\draw[draw=black,line width=.35pt] (490.700,705.000)--(490.700,711.000);\end{scope}
\begin{scope}[scale=.6666667,draw=black,text=black]\draw[draw=black,line width=.35pt] (495.467,705.000)--(495.467,711.000);\end{scope}
\begin{scope}[scale=.6666667,draw=black,text=black]\draw[draw=black,line width=.35pt] (500.233,705.000)--(500.233,711.000);\end{scope}
\begin{scope}[scale=.6666667,draw=black,text=black]\draw[draw=black,line width=.35pt] (505.000,705.000)--(505.000,719.000);\end{scope}
\begin{scope}[scale=.6666667,draw=black,text=black]\draw[draw=black,line width=.35pt] (509.767,705.000)--(509.767,711.000);\end{scope}
\begin{scope}[scale=.6666667,draw=black,text=black]\draw[draw=black,line width=.35pt] (514.533,705.000)--(514.533,711.000);\end{scope}
\begin{scope}[scale=.6666667,draw=black,text=black]\draw[draw=black,line width=.35pt] (519.300,705.000)--(519.300,711.000);\end{scope}
\begin{scope}[scale=.6666667,draw=black,text=black]\draw[draw=black,line width=.35pt] (524.067,705.000)--(524.067,711.000);\end{scope}
\begin{scope}[scale=.6666667,draw=black,text=black]\draw[draw=black,line width=.35pt] (528.833,705.000)--(528.833,714.000);\end{scope}
\begin{scope}[scale=.6666667,draw=black,text=black]\draw[draw=black,line width=.35pt] (533.600,705.000)--(533.600,711.000);\end{scope}
\begin{scope}[scale=.6666667,draw=black,text=black]\draw[draw=black,line width=.35pt] (538.367,705.000)--(538.367,711.000);\end{scope}
\begin{scope}[scale=.6666667,draw=black,text=black]\draw[draw=black,line width=.35pt] (543.133,705.000)--(543.133,711.000);\end{scope}
\begin{scope}[scale=.6666667,draw=black,text=black]\draw[draw=black,line width=.35pt] (547.900,705.000)--(547.900,711.000);\end{scope}
\begin{scope}[scale=.6666667,draw=black,text=black]\draw[draw=black,line width=.35pt] (552.667,705.000)--(552.667,719.000);\end{scope}
\begin{scope}[scale=.6666667,draw=black,text=black]\draw[draw=black,line width=.35pt] (557.433,705.000)--(557.433,711.000);\end{scope}
\begin{scope}[scale=.6666667,draw=black,text=black]\draw[draw=black,line width=.35pt] (562.200,705.000)--(562.200,711.000);\end{scope}
\begin{scope}[scale=.6666667,draw=black,text=black]\draw[draw=black,line width=.35pt] (566.967,705.000)--(566.967,711.000);\end{scope}
\begin{scope}[scale=.6666667,draw=black,text=black]\draw[draw=black,line width=.35pt] (571.733,705.000)--(571.733,711.000);\end{scope}
\begin{scope}[scale=.6666667,draw=black,text=black]\draw[draw=black,line width=.35pt] (576.500,705.000)--(576.500,714.000);\end{scope}
\begin{scope}[scale=.6666667,draw=black,text=black]\draw[draw=black,line width=.35pt] (581.267,705.000)--(581.267,711.000);\end{scope}
\begin{scope}[scale=.6666667,draw=black,text=black]\draw[draw=black,line width=.35pt] (586.033,705.000)--(586.033,711.000);\end{scope}
\begin{scope}[scale=.6666667,draw=black,text=black]\draw[draw=black,line width=.35pt] (590.800,705.000)--(590.800,711.000);\end{scope}
\begin{scope}[scale=.6666667,draw=black,text=black]\draw[draw=black,line width=.35pt] (595.567,705.000)--(595.567,711.000);\end{scope}
\begin{scope}[scale=.6666667,draw=black,text=black]\draw[draw=black,line width=.35pt] (600.333,705.000)--(600.333,719.000);\end{scope}
\begin{scope}[scale=.6666667,draw=black,text=black]\draw[draw=black,line width=.35pt] (605.100,705.000)--(605.100,711.000);\end{scope}
\begin{scope}[scale=.6666667,draw=black,text=black]\draw[draw=black,line width=.35pt] (609.867,705.000)--(609.867,711.000);\end{scope}
\begin{scope}[scale=.6666667,draw=black,text=black]\draw[draw=black,line width=.35pt] (614.633,705.000)--(614.633,711.000);\end{scope}
\begin{scope}[scale=.6666667,draw=black,text=black]\draw[draw=black,line width=.35pt] (619.400,705.000)--(619.400,711.000);\end{scope}
\begin{scope}[scale=.6666667,draw=black,text=black]\draw[draw=black,line width=.35pt] (624.167,705.000)--(624.167,714.000);\end{scope}
\begin{scope}[scale=.6666667,draw=black,text=black]\draw[draw=black,line width=.35pt] (628.933,705.000)--(628.933,711.000);\end{scope}
\begin{scope}[scale=.6666667,draw=black,text=black]\draw[draw=black,line width=.35pt] (633.700,705.000)--(633.700,711.000);\end{scope}
\begin{scope}[scale=.6666667,draw=black,text=black]\draw[draw=black,line width=.35pt] (638.467,705.000)--(638.467,711.000);\end{scope}
\begin{scope}[scale=.6666667,draw=black,text=black]\draw[draw=black,line width=.35pt] (643.233,705.000)--(643.233,711.000);\end{scope}
\begin{scope}[scale=.6666667,draw=black,text=black]\draw[draw=black,line width=.35pt] (648.000,705.000)--(648.000,719.000);\end{scope}
\begin{scope}[scale=.6666667,draw=black,text=black]\draw[draw=black,line width=.35pt] (652.767,705.000)--(652.767,711.000);\end{scope}
\begin{scope}[scale=.6666667,draw=black,text=black]\draw[draw=black,line width=.35pt] (657.533,705.000)--(657.533,711.000);\end{scope}
\begin{scope}[scale=.6666667,draw=black,text=black]\draw[draw=black,line width=.35pt] (662.300,705.000)--(662.300,711.000);\end{scope}
\begin{scope}[scale=.6666667,draw=black,text=black]\draw[draw=black,line width=.35pt] (667.067,705.000)--(667.067,711.000);\end{scope}
\begin{scope}[scale=.6666667,draw=black,text=black]\draw[draw=black,line width=.35pt] (671.833,705.000)--(671.833,714.000);\end{scope}
\begin{scope}[scale=.6666667,draw=black,text=black]\draw[draw=black,line width=.35pt] (676.600,705.000)--(676.600,711.000);\end{scope}
\begin{scope}[scale=.6666667,draw=black,text=black]\draw[draw=black,line width=.35pt] (681.367,705.000)--(681.367,711.000);\end{scope}
\begin{scope}[scale=.6666667,draw=black,text=black]\draw[draw=black,line width=.35pt] (686.133,705.000)--(686.133,711.000);\end{scope}
\begin{scope}[scale=.6666667,draw=black,text=black]\draw[draw=black,line width=.35pt] (690.900,705.000)--(690.900,711.000);\end{scope}
\begin{scope}[scale=.6666667,draw=black,text=black]\draw[draw=black,line width=.35pt] (695.667,705.000)--(695.667,719.000);\end{scope}
\begin{scope}[scale=.6666667,draw=black,text=black]\node[inner sep=0pt,font=\normalsize,text=black] at (219.000,738.000) {0};\end{scope}
\begin{scope}[scale=.6666667,draw=black,text=black]\node[inner sep=0pt,font=\normalsize,text=black] at (266.667,738.000) {1 cm};\end{scope}
\begin{scope}[scale=.6666667,draw=black,text=black]\node[inner sep=0pt,font=\normalsize,text=black] at (314.333,738.000) {2};\end{scope}
\begin{scope}[scale=.6666667,draw=black,text=black]\node[inner sep=0pt,font=\normalsize,text=black] at (362.000,738.000) {3};\end{scope}
\begin{scope}[scale=.6666667,draw=black,text=black]\node[inner sep=0pt,font=\normalsize,text=black] at (409.667,738.000) {4};\end{scope}
\begin{scope}[scale=.6666667,draw=black,text=black]\node[inner sep=0pt,font=\normalsize,text=black] at (457.333,738.000) {5};\end{scope}
\begin{scope}[scale=.6666667,draw=black,text=black]\node[inner sep=0pt,font=\normalsize,text=black] at (505.000,738.000) {6};\end{scope}
\begin{scope}[scale=.6666667,draw=black,text=black]\node[inner sep=0pt,font=\normalsize,text=black] at (552.667,738.000) {7};\end{scope}
\begin{scope}[scale=.6666667,draw=black,text=black]\node[inner sep=0pt,font=\normalsize,text=black] at (600.333,738.000) {8};\end{scope}
\begin{scope}[scale=.6666667,draw=black,text=black]\node[inner sep=0pt,font=\normalsize,text=black] at (648.000,738.000) {9};\end{scope}
\begin{scope}[scale=.6666667,draw=black,text=black]\node[inner sep=0pt,font=\normalsize,text=black] at (695.667,738.000) {10};\end{scope}
\begin{scope}[scale=.6666667,draw=black,text=black]\draw[draw=black,line width=.45pt,-{Stealth[length=3pt,width=2.4pt]}] (267.000,760.000)--(267.000,785.000);\end{scope}
\begin{scope}[scale=.6666667,draw=black,text=black]\draw[draw=black,line width=.45pt,-{Stealth[length=3pt,width=2.4pt]}] (362.000,760.000)--(362.000,785.000);\end{scope}
\begin{scope}[scale=.6666667,draw=black,text=black]\draw[draw=black,line width=.45pt,-{Stealth[length=3pt,width=2.4pt]}] (505.000,760.000)--(505.000,785.000);\end{scope}
\q{230}{790}{85}{12.5}{$\dfrac1{10}$分米}
\q{331}{790}{96}{12.5}{$\dfrac3{10}$分米}
\q{463}{804}{28}{13}{（}
\q{512}{804}{25}{13}{）}
\q{533}{804}{80}{12.5}{分米}
\q{225}{862}{103}{12.5}{0.1分米}
\q{324}{861}{28}{13}{（}
\q{378}{861}{25}{13}{）}
\q{399}{861}{76}{12.5}{分米}
\q{463}{861}{28}{13}{（}
\q{514}{861}{25}{13}{）}
\q{535}{861}{75}{12.5}{分米}
\hdsection{4}{921}{在计数器上表示下面各小数。}
\q{112}{967}{172}{13}{① 0.35}
\q{368}{967}{172}{13}{② 4.3}
\q{623}{967}{172}{13}{③ 0.732}
\begin{scope}[scale=.6666667,draw=black,text=black]\draw[line width=.5pt] (143,1015) rectangle (287,1079);\draw[line width=.5pt] (143,1019)--(287,1019);\end{scope}
\begin{scope}[scale=.6666667,draw=black,text=black]\draw[line width=.45pt] (165,1017)--(165,1079);\draw[line width=.45pt] (157,1015)..controls(166,995) and (171,1003)..(170,1015);\end{scope}
\begin{scope}[scale=.6666667,draw=black,text=black]\draw[line width=.45pt] (191,1017)--(191,1079);\draw[line width=.45pt] (183,1015)..controls(192,995) and (197,1003)..(196,1015);\end{scope}
\begin{scope}[scale=.6666667,draw=black,text=black]\draw[line width=.45pt] (224,1017)--(224,1079);\draw[line width=.45pt] (216,1015)..controls(225,995) and (230,1003)..(229,1015);\end{scope}
\begin{scope}[scale=.6666667,draw=black,text=black]\draw[line width=.45pt] (247,1017)--(247,1079);\draw[line width=.45pt] (239,1015)..controls(248,995) and (253,1003)..(252,1015);\end{scope}
\begin{scope}[scale=.6666667,draw=black,text=black]\draw[line width=.45pt] (273,1017)--(273,1079);\draw[line width=.45pt] (265,1015)..controls(274,995) and (279,1003)..(278,1015);\end{scope}
\begin{scope}[scale=.6666667,draw=black,text=black]\draw[draw=black,line width=.5pt,fill=white] (143.000,1078.000)--(151.000,1083.000)--(294.000,1083.000)--(287.000,1078.000)--cycle;\end{scope}
\begin{scope}[scale=.6666667,draw=black,text=black]\draw[draw=black,line width=.5pt,fill=white] (151.000,1083.000)--(294.000,1083.000)--(294.000,1147.000)--(151.000,1147.000)--cycle;\end{scope}
\begin{scope}[scale=.6666667,draw=black,text=black]\draw[draw=black,line width=.5pt,fill=white] (143.000,1078.000)--(151.000,1083.000)--(151.000,1147.000)--(143.000,1131.000)--cycle;\end{scope}
\begin{scope}[scale=.6666667,draw=black,text=black]\node[inner sep=0pt,font=\normalsize,text=black] at (165.000,1115.000) {\shortstack{十\\位}};\end{scope}
\begin{scope}[scale=.6666667,draw=black,text=black]\node[inner sep=0pt,font=\normalsize,text=black] at (191.000,1115.000) {\shortstack{个\\位}};\end{scope}
\begin{scope}[scale=.6666667,draw=black,text=black]\node[inner sep=0pt,font=\normalsize,text=black] at (224.000,1115.000) {\shortstack{十\\分\\位}};\end{scope}
\begin{scope}[scale=.6666667,draw=black,text=black]\node[inner sep=0pt,font=\normalsize,text=black] at (247.000,1115.000) {\shortstack{百\\分\\位}};\end{scope}
\begin{scope}[scale=.6666667,draw=black,text=black]\node[inner sep=0pt,font=\normalsize,text=black] at (273.000,1115.000) {\shortstack{千\\分\\位}};\end{scope}
\begin{scope}[scale=.6666667,draw=black,text=black]\draw[line width=.5pt] (398,1015) rectangle (542,1079);\draw[line width=.5pt] (398,1019)--(542,1019);\end{scope}
\begin{scope}[scale=.6666667,draw=black,text=black]\draw[line width=.45pt] (416,1017)--(416,1079);\draw[line width=.45pt] (408,1015)..controls(417,995) and (422,1003)..(421,1015);\end{scope}
\begin{scope}[scale=.6666667,draw=black,text=black]\draw[line width=.45pt] (441,1017)--(441,1079);\draw[line width=.45pt] (433,1015)..controls(442,995) and (447,1003)..(446,1015);\end{scope}
\begin{scope}[scale=.6666667,draw=black,text=black]\draw[line width=.45pt] (477,1017)--(477,1079);\draw[line width=.45pt] (469,1015)..controls(478,995) and (483,1003)..(482,1015);\end{scope}
\begin{scope}[scale=.6666667,draw=black,text=black]\draw[line width=.45pt] (502,1017)--(502,1079);\draw[line width=.45pt] (494,1015)..controls(503,995) and (508,1003)..(507,1015);\end{scope}
\begin{scope}[scale=.6666667,draw=black,text=black]\draw[line width=.45pt] (527,1017)--(527,1079);\draw[line width=.45pt] (519,1015)..controls(528,995) and (533,1003)..(532,1015);\end{scope}
\begin{scope}[scale=.6666667,draw=black,text=black]\draw[draw=black,line width=.5pt,fill=white] (398.000,1078.000)--(406.000,1083.000)--(549.000,1083.000)--(542.000,1078.000)--cycle;\end{scope}
\begin{scope}[scale=.6666667,draw=black,text=black]\draw[draw=black,line width=.5pt,fill=white] (406.000,1083.000)--(549.000,1083.000)--(549.000,1147.000)--(406.000,1147.000)--cycle;\end{scope}
\begin{scope}[scale=.6666667,draw=black,text=black]\draw[draw=black,line width=.5pt,fill=white] (398.000,1078.000)--(406.000,1083.000)--(406.000,1147.000)--(398.000,1131.000)--cycle;\end{scope}
\begin{scope}[scale=.6666667,draw=black,text=black]\node[inner sep=0pt,font=\normalsize,text=black] at (416.000,1115.000) {\shortstack{十\\位}};\end{scope}
\begin{scope}[scale=.6666667,draw=black,text=black]\node[inner sep=0pt,font=\normalsize,text=black] at (441.000,1115.000) {\shortstack{个\\位}};\end{scope}
\begin{scope}[scale=.6666667,draw=black,text=black]\node[inner sep=0pt,font=\normalsize,text=black] at (477.000,1115.000) {\shortstack{十\\分\\位}};\end{scope}
\begin{scope}[scale=.6666667,draw=black,text=black]\node[inner sep=0pt,font=\normalsize,text=black] at (502.000,1115.000) {\shortstack{百\\分\\位}};\end{scope}
\begin{scope}[scale=.6666667,draw=black,text=black]\node[inner sep=0pt,font=\normalsize,text=black] at (527.000,1115.000) {\shortstack{千\\分\\位}};\end{scope}
\begin{scope}[scale=.6666667,draw=black,text=black]\draw[line width=.5pt] (654,1015) rectangle (798,1079);\draw[line width=.5pt] (654,1019)--(798,1019);\end{scope}
\begin{scope}[scale=.6666667,draw=black,text=black]\draw[line width=.45pt] (675,1017)--(675,1079);\draw[line width=.45pt] (667,1015)..controls(676,995) and (681,1003)..(680,1015);\end{scope}
\begin{scope}[scale=.6666667,draw=black,text=black]\draw[line width=.45pt] (701,1017)--(701,1079);\draw[line width=.45pt] (693,1015)..controls(702,995) and (707,1003)..(706,1015);\end{scope}
\begin{scope}[scale=.6666667,draw=black,text=black]\draw[line width=.45pt] (731,1017)--(731,1079);\draw[line width=.45pt] (723,1015)..controls(732,995) and (737,1003)..(736,1015);\end{scope}
\begin{scope}[scale=.6666667,draw=black,text=black]\draw[line width=.45pt] (757,1017)--(757,1079);\draw[line width=.45pt] (749,1015)..controls(758,995) and (763,1003)..(762,1015);\end{scope}
\begin{scope}[scale=.6666667,draw=black,text=black]\draw[line width=.45pt] (781,1017)--(781,1079);\draw[line width=.45pt] (773,1015)..controls(782,995) and (787,1003)..(786,1015);\end{scope}
\begin{scope}[scale=.6666667,draw=black,text=black]\draw[draw=black,line width=.5pt,fill=white] (654.000,1078.000)--(662.000,1083.000)--(805.000,1083.000)--(798.000,1078.000)--cycle;\end{scope}
\begin{scope}[scale=.6666667,draw=black,text=black]\draw[draw=black,line width=.5pt,fill=white] (662.000,1083.000)--(805.000,1083.000)--(805.000,1147.000)--(662.000,1147.000)--cycle;\end{scope}
\begin{scope}[scale=.6666667,draw=black,text=black]\draw[draw=black,line width=.5pt,fill=white] (654.000,1078.000)--(662.000,1083.000)--(662.000,1147.000)--(654.000,1131.000)--cycle;\end{scope}
\begin{scope}[scale=.6666667,draw=black,text=black]\node[inner sep=0pt,font=\normalsize,text=black] at (675.000,1115.000) {\shortstack{十\\位}};\end{scope}
\begin{scope}[scale=.6666667,draw=black,text=black]\node[inner sep=0pt,font=\normalsize,text=black] at (701.000,1115.000) {\shortstack{个\\位}};\end{scope}
\begin{scope}[scale=.6666667,draw=black,text=black]\node[inner sep=0pt,font=\normalsize,text=black] at (731.000,1115.000) {\shortstack{十\\分\\位}};\end{scope}
\begin{scope}[scale=.6666667,draw=black,text=black]\node[inner sep=0pt,font=\normalsize,text=black] at (757.000,1115.000) {\shortstack{百\\分\\位}};\end{scope}
\begin{scope}[scale=.6666667,draw=black,text=black]\node[inner sep=0pt,font=\normalsize,text=black] at (781.000,1115.000) {\shortstack{千\\分\\位}};\end{scope}

\begin{scope}[color=answerred,draw=answerred,text=answerred]\ans{128.00}{226.67}{33.33}{10}{$\vphantom{0}0.4$}
\ans{294.00}{226.67}{33.33}{10}{$\vphantom{0}0.33$}
\ans{455.33}{225.33}{33.33}{10}{$\vphantom{0}0.203$}
\fill[answerred,opacity=.28](95.333,296.000) rectangle (181.733,394.000);
\fill[answerred,opacity=.28](259.333,296.000) rectangle (345.733,305.733);
\fill[answerred,opacity=.28](424.000,296.000) rectangle (472.667,393.333);
\fill[answerred,opacity=.28](472.667,296.000) rectangle (482.400,325.200);
\ans{228.67}{573.33}{48.00}{11}{$\vphantom{0}0.3$}
\ans{322.67}{534.00}{42.00}{11}{$\vphantom{0}\frac6{10}$}
\ans{321.33}{573.33}{48.00}{11}{$\vphantom{0}0.6$}
\filldraw[answerred,fill=answerred!65] (149.333,714.000) ellipse [x radius=4bp,y radius=2.2bp];
\filldraw[answerred,fill=answerred!65] (149.333,708.667) ellipse [x radius=4bp,y radius=2.2bp];
\filldraw[answerred,fill=answerred!65] (149.333,703.333) ellipse [x radius=4bp,y radius=2.2bp];
\filldraw[answerred,fill=answerred!65] (164.667,714.000) ellipse [x radius=4bp,y radius=2.2bp];
\filldraw[answerred,fill=answerred!65] (164.667,708.667) ellipse [x radius=4bp,y radius=2.2bp];
\filldraw[answerred,fill=answerred!65] (164.667,703.333) ellipse [x radius=4bp,y radius=2.2bp];
\filldraw[answerred,fill=answerred!65] (164.667,698.000) ellipse [x radius=4bp,y radius=2.2bp];
\filldraw[answerred,fill=answerred!65] (164.667,692.667) ellipse [x radius=4bp,y radius=2.2bp];
\filldraw[answerred,fill=answerred!65] (294.000,714.000) ellipse [x radius=4bp,y radius=2.2bp];
\filldraw[answerred,fill=answerred!65] (294.000,708.667) ellipse [x radius=4bp,y radius=2.2bp];
\filldraw[answerred,fill=answerred!65] (294.000,703.333) ellipse [x radius=4bp,y radius=2.2bp];
\filldraw[answerred,fill=answerred!65] (294.000,698.000) ellipse [x radius=4bp,y radius=2.2bp];
\filldraw[answerred,fill=answerred!65] (318.000,714.000) ellipse [x radius=4bp,y radius=2.2bp];
\filldraw[answerred,fill=answerred!65] (318.000,708.667) ellipse [x radius=4bp,y radius=2.2bp];
\filldraw[answerred,fill=answerred!65] (318.000,703.333) ellipse [x radius=4bp,y radius=2.2bp];
\filldraw[answerred,fill=answerred!65] (487.333,714.000) ellipse [x radius=4bp,y radius=2.2bp];
\filldraw[answerred,fill=answerred!65] (487.333,708.667) ellipse [x radius=4bp,y radius=2.2bp];
\filldraw[answerred,fill=answerred!65] (487.333,703.333) ellipse [x radius=4bp,y radius=2.2bp];
\filldraw[answerred,fill=answerred!65] (487.333,698.000) ellipse [x radius=4bp,y radius=2.2bp];
\filldraw[answerred,fill=answerred!65] (487.333,692.667) ellipse [x radius=4bp,y radius=2.2bp];
\filldraw[answerred,fill=answerred!65] (487.333,687.333) ellipse [x radius=4bp,y radius=2.2bp];
\filldraw[answerred,fill=answerred!65] (487.333,682.000) ellipse [x radius=4bp,y radius=2.2bp];
\filldraw[answerred,fill=answerred!65] (504.667,714.000) ellipse [x radius=4bp,y radius=2.2bp];
\filldraw[answerred,fill=answerred!65] (504.667,708.667) ellipse [x radius=4bp,y radius=2.2bp];
\filldraw[answerred,fill=answerred!65] (504.667,703.333) ellipse [x radius=4bp,y radius=2.2bp];
\filldraw[answerred,fill=answerred!65] (520.667,714.000) ellipse [x radius=4bp,y radius=2.2bp];
\filldraw[answerred,fill=answerred!65] (520.667,708.667) ellipse [x radius=4bp,y radius=2.2bp];
\end{scope}
\end{scope}\pgfresetboundingbox\path[use as bounding box] (56.0,288.0) rectangle (527.0,436.0);
\end{tikzpicture}}\par



\par





\par\noindent

\QHeading{第2题③　0.53的另一种涂色方式}

\Answer{方法一可涂满5列再加3格；方法二可涂满5行再加3格。只要把同一个整体平均分成100份并涂53份，两种方式都表示0.53。}
\par\smallskip\noindent\begin{adjustbox}{max width=\linewidth,center}
\begin{tikzpicture}[x=.28cm,y=.28cm,font=\normalsize]
\fill[answerred!22] (0,0) rectangle (5,10);\fill[answerred!22] (5,0) rectangle (6,3);
\draw[black!55,step=1] (0,0) grid (10,10);\node[below] at (5,0){按列涂53格};
\begin{scope}[xshift=3.9cm]
\fill[answerred!22] (0,0) rectangle (10,5);\fill[answerred!22] (0,5) rectangle (3,6);
\draw[black!55,step=1] (0,0) grid (10,10);\node[below] at (5,0){按行涂53格};
\end{scope}
\end{tikzpicture}
\end{adjustbox}\par\smallskip

\QHeading{第2题①②　其他正确的涂色方式}

\Answer{0.9表示10个同样大小的小格中的9格，可以留最左或最右的一格不涂；0.09表示100个同样大小小格中的9格，可以涂连续9格，也可以涂成3行3列的小方块。位置和形状可以变，所占整体的份数不能变。}

\QHeading{第3题　照样子填一填。}

\par\noindent\makebox[\linewidth][c]{
\begin{tikzpicture}[x=1bp,y=-1bp,color=black,draw=black,text=black]
\begin{scope}\clip (56.0,464.0) rectangle (494.0,612.0);
\hdhead{46}{小数的意义与性质——信息窗1}
\hdsection{1}{135}{用小数表示下图的涂色部分。}
\begin{scope}[scale=.6666667,draw=black,text=black]\path[pattern=north east lines,pattern color=black] (143,178) rectangle (200.6,323);\end{scope}
\begin{scope}[scale=.6666667,draw=black,text=black]\draw[draw=black,line width=.4pt] (143,178) rectangle (287,323);\foreach \i in {1,...,9}{\draw[draw=black,line width=.4pt] ({143+14.4*\i},178)--({143+14.4*\i},323);}\foreach \j in {1,...,0}{\draw[draw=black,line width=.4pt] (143,{178+145.0*\j})--(287,{178+145.0*\j});}\end{scope}
\begin{scope}[scale=.6666667,draw=black,text=black]\path[pattern=north east lines,pattern color=black] (389,178) rectangle (432.2,322);\end{scope}
\begin{scope}[scale=.6666667,draw=black,text=black]\path[pattern=north east lines,pattern color=black] (432.2,278.8) rectangle (446.59999999999997,322.0);\end{scope}
\begin{scope}[scale=.6666667,draw=black,text=black]\draw[draw=black,line width=.35pt] (389,178) rectangle (533,322);\foreach \i in {1,...,9}{\draw[draw=black,line width=.35pt] ({389+14.4*\i},178)--({389+14.4*\i},322);}\foreach \j in {1,...,9}{\draw[draw=black,line width=.35pt] (389,{178+14.4*\j})--(533,{178+14.4*\j});}\end{scope}
\begin{scope}[scale=.6666667,draw=black,text=black]\fill[white] (636.0000,319.0000)--(736.0000,319.0000)--(736.0000,219.0000)--(636.0000,219.0000)--cycle;\end{scope}
\begin{scope}[scale=.6666667,draw=black,text=black]\fill[white] (636.0000,219.0000)--(736.0000,219.0000)--(769.0000,179.0000)--(669.0000,179.0000)--cycle;\end{scope}
\begin{scope}[scale=.6666667,draw=black,text=black]\path[pattern=north east lines,pattern color=black] (736.0000,319.0000)--(769.0000,279.0000)--(769.0000,179.0000)--(736.0000,219.0000)--cycle;\end{scope}
\begin{scope}[scale=.6666667,draw=black,text=black]\path[pattern=north east lines,pattern color=black] (706.0000,249.0000)--(716.0000,249.0000)--(716.0000,239.0000)--(706.0000,239.0000)--cycle;\end{scope}
\begin{scope}[scale=.6666667,draw=black,text=black]\path[pattern=north east lines,pattern color=black] (706.0000,239.0000)--(716.0000,239.0000)--(716.0000,229.0000)--(706.0000,229.0000)--cycle;\end{scope}
\begin{scope}[scale=.6666667,draw=black,text=black]\path[pattern=north east lines,pattern color=black] (706.0000,229.0000)--(716.0000,229.0000)--(716.0000,219.0000)--(706.0000,219.0000)--cycle;\end{scope}
\begin{scope}[scale=.6666667,draw=black,text=black]\path[pattern=north east lines,pattern color=black] (706.0000,219.0000)--(716.0000,219.0000)--(719.3000,215.0000)--(709.3000,215.0000)--cycle;\end{scope}
\begin{scope}[scale=.6666667,draw=black,text=black]\path[pattern=north east lines,pattern color=black] (716.0000,319.0000)--(726.0000,319.0000)--(726.0000,309.0000)--(716.0000,309.0000)--cycle;\end{scope}
\begin{scope}[scale=.6666667,draw=black,text=black]\path[pattern=north east lines,pattern color=black] (716.0000,309.0000)--(726.0000,309.0000)--(726.0000,299.0000)--(716.0000,299.0000)--cycle;\end{scope}
\begin{scope}[scale=.6666667,draw=black,text=black]\path[pattern=north east lines,pattern color=black] (716.0000,299.0000)--(726.0000,299.0000)--(726.0000,289.0000)--(716.0000,289.0000)--cycle;\end{scope}
\begin{scope}[scale=.6666667,draw=black,text=black]\path[pattern=north east lines,pattern color=black] (716.0000,289.0000)--(726.0000,289.0000)--(726.0000,279.0000)--(716.0000,279.0000)--cycle;\end{scope}
\begin{scope}[scale=.6666667,draw=black,text=black]\path[pattern=north east lines,pattern color=black] (716.0000,279.0000)--(726.0000,279.0000)--(726.0000,269.0000)--(716.0000,269.0000)--cycle;\end{scope}
\begin{scope}[scale=.6666667,draw=black,text=black]\path[pattern=north east lines,pattern color=black] (716.0000,269.0000)--(726.0000,269.0000)--(726.0000,259.0000)--(716.0000,259.0000)--cycle;\end{scope}
\begin{scope}[scale=.6666667,draw=black,text=black]\path[pattern=north east lines,pattern color=black] (716.0000,259.0000)--(726.0000,259.0000)--(726.0000,249.0000)--(716.0000,249.0000)--cycle;\end{scope}
\begin{scope}[scale=.6666667,draw=black,text=black]\path[pattern=north east lines,pattern color=black] (716.0000,249.0000)--(726.0000,249.0000)--(726.0000,239.0000)--(716.0000,239.0000)--cycle;\end{scope}
\begin{scope}[scale=.6666667,draw=black,text=black]\path[pattern=north east lines,pattern color=black] (716.0000,239.0000)--(726.0000,239.0000)--(726.0000,229.0000)--(716.0000,229.0000)--cycle;\end{scope}
\begin{scope}[scale=.6666667,draw=black,text=black]\path[pattern=north east lines,pattern color=black] (716.0000,229.0000)--(726.0000,229.0000)--(726.0000,219.0000)--(716.0000,219.0000)--cycle;\end{scope}
\begin{scope}[scale=.6666667,draw=black,text=black]\path[pattern=north east lines,pattern color=black] (716.0000,219.0000)--(726.0000,219.0000)--(729.3000,215.0000)--(719.3000,215.0000)--cycle;\end{scope}
\begin{scope}[scale=.6666667,draw=black,text=black]\path[pattern=north east lines,pattern color=black] (719.3000,215.0000)--(729.3000,215.0000)--(732.6000,211.0000)--(722.6000,211.0000)--cycle;\end{scope}
\begin{scope}[scale=.6666667,draw=black,text=black]\path[pattern=north east lines,pattern color=black] (722.6000,211.0000)--(732.6000,211.0000)--(735.9000,207.0000)--(725.9000,207.0000)--cycle;\end{scope}
\begin{scope}[scale=.6666667,draw=black,text=black]\path[pattern=north east lines,pattern color=black] (725.9000,207.0000)--(735.9000,207.0000)--(739.2000,203.0000)--(729.2000,203.0000)--cycle;\end{scope}
\begin{scope}[scale=.6666667,draw=black,text=black]\path[pattern=north east lines,pattern color=black] (729.2000,203.0000)--(739.2000,203.0000)--(742.5000,199.0000)--(732.5000,199.0000)--cycle;\end{scope}
\begin{scope}[scale=.6666667,draw=black,text=black]\path[pattern=north east lines,pattern color=black] (732.5000,199.0000)--(742.5000,199.0000)--(745.8000,195.0000)--(735.8000,195.0000)--cycle;\end{scope}
\begin{scope}[scale=.6666667,draw=black,text=black]\path[pattern=north east lines,pattern color=black] (735.8000,195.0000)--(745.8000,195.0000)--(749.1000,191.0000)--(739.1000,191.0000)--cycle;\end{scope}
\begin{scope}[scale=.6666667,draw=black,text=black]\path[pattern=north east lines,pattern color=black] (739.1000,191.0000)--(749.1000,191.0000)--(752.4000,187.0000)--(742.4000,187.0000)--cycle;\end{scope}
\begin{scope}[scale=.6666667,draw=black,text=black]\path[pattern=north east lines,pattern color=black] (742.4000,187.0000)--(752.4000,187.0000)--(755.7000,183.0000)--(745.7000,183.0000)--cycle;\end{scope}
\begin{scope}[scale=.6666667,draw=black,text=black]\path[pattern=north east lines,pattern color=black] (745.7000,183.0000)--(755.7000,183.0000)--(759.0000,179.0000)--(749.0000,179.0000)--cycle;\end{scope}
\begin{scope}[scale=.6666667,draw=black,text=black]\path[pattern=north east lines,pattern color=black] (726.0000,319.0000)--(736.0000,319.0000)--(736.0000,309.0000)--(726.0000,309.0000)--cycle;\end{scope}
\begin{scope}[scale=.6666667,draw=black,text=black]\path[pattern=north east lines,pattern color=black] (726.0000,309.0000)--(736.0000,309.0000)--(736.0000,299.0000)--(726.0000,299.0000)--cycle;\end{scope}
\begin{scope}[scale=.6666667,draw=black,text=black]\path[pattern=north east lines,pattern color=black] (726.0000,299.0000)--(736.0000,299.0000)--(736.0000,289.0000)--(726.0000,289.0000)--cycle;\end{scope}
\begin{scope}[scale=.6666667,draw=black,text=black]\path[pattern=north east lines,pattern color=black] (726.0000,289.0000)--(736.0000,289.0000)--(736.0000,279.0000)--(726.0000,279.0000)--cycle;\end{scope}
\begin{scope}[scale=.6666667,draw=black,text=black]\path[pattern=north east lines,pattern color=black] (726.0000,279.0000)--(736.0000,279.0000)--(736.0000,269.0000)--(726.0000,269.0000)--cycle;\end{scope}
\begin{scope}[scale=.6666667,draw=black,text=black]\path[pattern=north east lines,pattern color=black] (726.0000,269.0000)--(736.0000,269.0000)--(736.0000,259.0000)--(726.0000,259.0000)--cycle;\end{scope}
\begin{scope}[scale=.6666667,draw=black,text=black]\path[pattern=north east lines,pattern color=black] (726.0000,259.0000)--(736.0000,259.0000)--(736.0000,249.0000)--(726.0000,249.0000)--cycle;\end{scope}
\begin{scope}[scale=.6666667,draw=black,text=black]\path[pattern=north east lines,pattern color=black] (726.0000,249.0000)--(736.0000,249.0000)--(736.0000,239.0000)--(726.0000,239.0000)--cycle;\end{scope}
\begin{scope}[scale=.6666667,draw=black,text=black]\path[pattern=north east lines,pattern color=black] (726.0000,239.0000)--(736.0000,239.0000)--(736.0000,229.0000)--(726.0000,229.0000)--cycle;\end{scope}
\begin{scope}[scale=.6666667,draw=black,text=black]\path[pattern=north east lines,pattern color=black] (726.0000,229.0000)--(736.0000,229.0000)--(736.0000,219.0000)--(726.0000,219.0000)--cycle;\end{scope}
\begin{scope}[scale=.6666667,draw=black,text=black]\path[pattern=north east lines,pattern color=black] (726.0000,219.0000)--(736.0000,219.0000)--(739.3000,215.0000)--(729.3000,215.0000)--cycle;\end{scope}
\begin{scope}[scale=.6666667,draw=black,text=black]\path[pattern=north east lines,pattern color=black] (729.3000,215.0000)--(739.3000,215.0000)--(742.6000,211.0000)--(732.6000,211.0000)--cycle;\end{scope}
\begin{scope}[scale=.6666667,draw=black,text=black]\path[pattern=north east lines,pattern color=black] (732.6000,211.0000)--(742.6000,211.0000)--(745.9000,207.0000)--(735.9000,207.0000)--cycle;\end{scope}
\begin{scope}[scale=.6666667,draw=black,text=black]\path[pattern=north east lines,pattern color=black] (735.9000,207.0000)--(745.9000,207.0000)--(749.2000,203.0000)--(739.2000,203.0000)--cycle;\end{scope}
\begin{scope}[scale=.6666667,draw=black,text=black]\path[pattern=north east lines,pattern color=black] (739.2000,203.0000)--(749.2000,203.0000)--(752.5000,199.0000)--(742.5000,199.0000)--cycle;\end{scope}
\begin{scope}[scale=.6666667,draw=black,text=black]\path[pattern=north east lines,pattern color=black] (742.5000,199.0000)--(752.5000,199.0000)--(755.8000,195.0000)--(745.8000,195.0000)--cycle;\end{scope}
\begin{scope}[scale=.6666667,draw=black,text=black]\path[pattern=north east lines,pattern color=black] (745.8000,195.0000)--(755.8000,195.0000)--(759.1000,191.0000)--(749.1000,191.0000)--cycle;\end{scope}
\begin{scope}[scale=.6666667,draw=black,text=black]\path[pattern=north east lines,pattern color=black] (749.1000,191.0000)--(759.1000,191.0000)--(762.4000,187.0000)--(752.4000,187.0000)--cycle;\end{scope}
\begin{scope}[scale=.6666667,draw=black,text=black]\path[pattern=north east lines,pattern color=black] (752.4000,187.0000)--(762.4000,187.0000)--(765.7000,183.0000)--(755.7000,183.0000)--cycle;\end{scope}
\begin{scope}[scale=.6666667,draw=black,text=black]\path[pattern=north east lines,pattern color=black] (755.7000,183.0000)--(765.7000,183.0000)--(769.0000,179.0000)--(759.0000,179.0000)--cycle;\end{scope}
\begin{scope}[scale=.6666667,draw=black,text=black]\draw[draw=black!85,line width=.32pt] (636.000,319.000)--(636.000,219.000);\end{scope}
\begin{scope}[scale=.6666667,draw=black,text=black]\draw[draw=black!85,line width=.32pt] (636.000,319.000)--(736.000,319.000);\end{scope}
\begin{scope}[scale=.6666667,draw=black,text=black]\draw[draw=black!85,line width=.32pt] (636.000,219.000)--(669.000,179.000);\end{scope}
\begin{scope}[scale=.6666667,draw=black,text=black]\draw[draw=black!85,line width=.32pt] (636.000,219.000)--(736.000,219.000);\end{scope}
\begin{scope}[scale=.6666667,draw=black,text=black]\draw[draw=black!85,line width=.32pt] (736.000,319.000)--(736.000,219.000);\end{scope}
\begin{scope}[scale=.6666667,draw=black,text=black]\draw[draw=black!85,line width=.32pt] (736.000,319.000)--(769.000,279.000);\end{scope}
\begin{scope}[scale=.6666667,draw=black,text=black]\draw[draw=black!85,line width=.32pt] (646.000,319.000)--(646.000,219.000);\end{scope}
\begin{scope}[scale=.6666667,draw=black,text=black]\draw[draw=black!85,line width=.32pt] (636.000,309.000)--(736.000,309.000);\end{scope}
\begin{scope}[scale=.6666667,draw=black,text=black]\draw[draw=black!85,line width=.32pt] (646.000,219.000)--(679.000,179.000);\end{scope}
\begin{scope}[scale=.6666667,draw=black,text=black]\draw[draw=black!85,line width=.32pt] (639.300,215.000)--(739.300,215.000);\end{scope}
\begin{scope}[scale=.6666667,draw=black,text=black]\draw[draw=black!85,line width=.32pt] (739.300,315.000)--(739.300,215.000);\end{scope}
\begin{scope}[scale=.6666667,draw=black,text=black]\draw[draw=black!85,line width=.32pt] (736.000,309.000)--(769.000,269.000);\end{scope}
\begin{scope}[scale=.6666667,draw=black,text=black]\draw[draw=black!85,line width=.32pt] (656.000,319.000)--(656.000,219.000);\end{scope}
\begin{scope}[scale=.6666667,draw=black,text=black]\draw[draw=black!85,line width=.32pt] (636.000,299.000)--(736.000,299.000);\end{scope}
\begin{scope}[scale=.6666667,draw=black,text=black]\draw[draw=black!85,line width=.32pt] (656.000,219.000)--(689.000,179.000);\end{scope}
\begin{scope}[scale=.6666667,draw=black,text=black]\draw[draw=black!85,line width=.32pt] (642.600,211.000)--(742.600,211.000);\end{scope}
\begin{scope}[scale=.6666667,draw=black,text=black]\draw[draw=black!85,line width=.32pt] (742.600,311.000)--(742.600,211.000);\end{scope}
\begin{scope}[scale=.6666667,draw=black,text=black]\draw[draw=black!85,line width=.32pt] (736.000,299.000)--(769.000,259.000);\end{scope}
\begin{scope}[scale=.6666667,draw=black,text=black]\draw[draw=black!85,line width=.32pt] (666.000,319.000)--(666.000,219.000);\end{scope}
\begin{scope}[scale=.6666667,draw=black,text=black]\draw[draw=black!85,line width=.32pt] (636.000,289.000)--(736.000,289.000);\end{scope}
\begin{scope}[scale=.6666667,draw=black,text=black]\draw[draw=black!85,line width=.32pt] (666.000,219.000)--(699.000,179.000);\end{scope}
\begin{scope}[scale=.6666667,draw=black,text=black]\draw[draw=black!85,line width=.32pt] (645.900,207.000)--(745.900,207.000);\end{scope}
\begin{scope}[scale=.6666667,draw=black,text=black]\draw[draw=black!85,line width=.32pt] (745.900,307.000)--(745.900,207.000);\end{scope}
\begin{scope}[scale=.6666667,draw=black,text=black]\draw[draw=black!85,line width=.32pt] (736.000,289.000)--(769.000,249.000);\end{scope}
\begin{scope}[scale=.6666667,draw=black,text=black]\draw[draw=black!85,line width=.32pt] (676.000,319.000)--(676.000,219.000);\end{scope}
\begin{scope}[scale=.6666667,draw=black,text=black]\draw[draw=black!85,line width=.32pt] (636.000,279.000)--(736.000,279.000);\end{scope}
\begin{scope}[scale=.6666667,draw=black,text=black]\draw[draw=black!85,line width=.32pt] (676.000,219.000)--(709.000,179.000);\end{scope}
\begin{scope}[scale=.6666667,draw=black,text=black]\draw[draw=black!85,line width=.32pt] (649.200,203.000)--(749.200,203.000);\end{scope}
\begin{scope}[scale=.6666667,draw=black,text=black]\draw[draw=black!85,line width=.32pt] (749.200,303.000)--(749.200,203.000);\end{scope}
\begin{scope}[scale=.6666667,draw=black,text=black]\draw[draw=black!85,line width=.32pt] (736.000,279.000)--(769.000,239.000);\end{scope}
\begin{scope}[scale=.6666667,draw=black,text=black]\draw[draw=black!85,line width=.32pt] (686.000,319.000)--(686.000,219.000);\end{scope}
\begin{scope}[scale=.6666667,draw=black,text=black]\draw[draw=black!85,line width=.32pt] (636.000,269.000)--(736.000,269.000);\end{scope}
\begin{scope}[scale=.6666667,draw=black,text=black]\draw[draw=black!85,line width=.32pt] (686.000,219.000)--(719.000,179.000);\end{scope}
\begin{scope}[scale=.6666667,draw=black,text=black]\draw[draw=black!85,line width=.32pt] (652.500,199.000)--(752.500,199.000);\end{scope}
\begin{scope}[scale=.6666667,draw=black,text=black]\draw[draw=black!85,line width=.32pt] (752.500,299.000)--(752.500,199.000);\end{scope}
\begin{scope}[scale=.6666667,draw=black,text=black]\draw[draw=black!85,line width=.32pt] (736.000,269.000)--(769.000,229.000);\end{scope}
\begin{scope}[scale=.6666667,draw=black,text=black]\draw[draw=black!85,line width=.32pt] (696.000,319.000)--(696.000,219.000);\end{scope}
\begin{scope}[scale=.6666667,draw=black,text=black]\draw[draw=black!85,line width=.32pt] (636.000,259.000)--(736.000,259.000);\end{scope}
\begin{scope}[scale=.6666667,draw=black,text=black]\draw[draw=black!85,line width=.32pt] (696.000,219.000)--(729.000,179.000);\end{scope}
\begin{scope}[scale=.6666667,draw=black,text=black]\draw[draw=black!85,line width=.32pt] (655.800,195.000)--(755.800,195.000);\end{scope}
\begin{scope}[scale=.6666667,draw=black,text=black]\draw[draw=black!85,line width=.32pt] (755.800,295.000)--(755.800,195.000);\end{scope}
\begin{scope}[scale=.6666667,draw=black,text=black]\draw[draw=black!85,line width=.32pt] (736.000,259.000)--(769.000,219.000);\end{scope}
\begin{scope}[scale=.6666667,draw=black,text=black]\draw[draw=black!85,line width=.32pt] (706.000,319.000)--(706.000,219.000);\end{scope}
\begin{scope}[scale=.6666667,draw=black,text=black]\draw[draw=black!85,line width=.32pt] (636.000,249.000)--(736.000,249.000);\end{scope}
\begin{scope}[scale=.6666667,draw=black,text=black]\draw[draw=black!85,line width=.32pt] (706.000,219.000)--(739.000,179.000);\end{scope}
\begin{scope}[scale=.6666667,draw=black,text=black]\draw[draw=black!85,line width=.32pt] (659.100,191.000)--(759.100,191.000);\end{scope}
\begin{scope}[scale=.6666667,draw=black,text=black]\draw[draw=black!85,line width=.32pt] (759.100,291.000)--(759.100,191.000);\end{scope}
\begin{scope}[scale=.6666667,draw=black,text=black]\draw[draw=black!85,line width=.32pt] (736.000,249.000)--(769.000,209.000);\end{scope}
\begin{scope}[scale=.6666667,draw=black,text=black]\draw[draw=black!85,line width=.32pt] (716.000,319.000)--(716.000,219.000);\end{scope}
\begin{scope}[scale=.6666667,draw=black,text=black]\draw[draw=black!85,line width=.32pt] (636.000,239.000)--(736.000,239.000);\end{scope}
\begin{scope}[scale=.6666667,draw=black,text=black]\draw[draw=black!85,line width=.32pt] (716.000,219.000)--(749.000,179.000);\end{scope}
\begin{scope}[scale=.6666667,draw=black,text=black]\draw[draw=black!85,line width=.32pt] (662.400,187.000)--(762.400,187.000);\end{scope}
\begin{scope}[scale=.6666667,draw=black,text=black]\draw[draw=black!85,line width=.32pt] (762.400,287.000)--(762.400,187.000);\end{scope}
\begin{scope}[scale=.6666667,draw=black,text=black]\draw[draw=black!85,line width=.32pt] (736.000,239.000)--(769.000,199.000);\end{scope}
\begin{scope}[scale=.6666667,draw=black,text=black]\draw[draw=black!85,line width=.32pt] (726.000,319.000)--(726.000,219.000);\end{scope}
\begin{scope}[scale=.6666667,draw=black,text=black]\draw[draw=black!85,line width=.32pt] (636.000,229.000)--(736.000,229.000);\end{scope}
\begin{scope}[scale=.6666667,draw=black,text=black]\draw[draw=black!85,line width=.32pt] (726.000,219.000)--(759.000,179.000);\end{scope}
\begin{scope}[scale=.6666667,draw=black,text=black]\draw[draw=black!85,line width=.32pt] (665.700,183.000)--(765.700,183.000);\end{scope}
\begin{scope}[scale=.6666667,draw=black,text=black]\draw[draw=black!85,line width=.32pt] (765.700,283.000)--(765.700,183.000);\end{scope}
\begin{scope}[scale=.6666667,draw=black,text=black]\draw[draw=black!85,line width=.32pt] (736.000,229.000)--(769.000,189.000);\end{scope}
\begin{scope}[scale=.6666667,draw=black,text=black]\draw[draw=black!85,line width=.32pt] (736.000,319.000)--(736.000,219.000);\end{scope}
\begin{scope}[scale=.6666667,draw=black,text=black]\draw[draw=black!85,line width=.32pt] (636.000,219.000)--(736.000,219.000);\end{scope}
\begin{scope}[scale=.6666667,draw=black,text=black]\draw[draw=black!85,line width=.32pt] (736.000,219.000)--(769.000,179.000);\end{scope}
\begin{scope}[scale=.6666667,draw=black,text=black]\draw[draw=black!85,line width=.32pt] (669.000,179.000)--(769.000,179.000);\end{scope}
\begin{scope}[scale=.6666667,draw=black,text=black]\draw[draw=black!85,line width=.32pt] (769.000,279.000)--(769.000,179.000);\end{scope}
\begin{scope}[scale=.6666667,draw=black,text=black]\draw[draw=black!85,line width=.32pt] (736.000,219.000)--(769.000,179.000);\end{scope}
\begin{scope}[scale=.6666667,draw=black,text=black]\draw[draw=black,line width=.55pt] (636.000,319.000)--(736.000,319.000)--(736.000,219.000)--(636.000,219.000)--cycle;\end{scope}
\begin{scope}[scale=.6666667,draw=black,text=black]\draw[draw=black,line width=.55pt] (636.000,219.000)--(736.000,219.000)--(769.000,179.000)--(669.000,179.000)--cycle;\end{scope}
\begin{scope}[scale=.6666667,draw=black,text=black]\draw[draw=black,line width=.55pt] (736.000,319.000)--(769.000,279.000)--(769.000,179.000)--(736.000,219.000)--cycle;\end{scope}
\begin{scope}[scale=.6666667,draw=black,text=black]\node[inner sep=0pt,font=\normalsize,text=black] at (123.000,189.000) {①};\end{scope}
\begin{scope}[scale=.6666667,draw=black,text=black]\node[inner sep=0pt,font=\normalsize,text=black] at (370.000,189.000) {②};\end{scope}
\begin{scope}[scale=.6666667,draw=black,text=black]\node[inner sep=0pt,font=\normalsize,text=black] at (617.000,190.000) {③};\end{scope}
\q{171}{339}{28}{13}{（}
\q{247}{339}{25}{13}{）}
\q{418}{339}{28}{13}{（}
\q{493}{339}{25}{13}{）}
\q{657}{338}{28}{13}{（}
\q{733}{338}{25}{13}{）}
\hdsection{2}{395}{给图形涂色表示它下面的小数。}
\begin{scope}[scale=.6666667,draw=black,text=black]\draw[draw=black,line width=.4pt] (143,444) rectangle (287,591);\foreach \i in {1,...,9}{\draw[draw=black,line width=.4pt] ({143+14.4*\i},444)--({143+14.4*\i},591);}\foreach \j in {1,...,0}{\draw[draw=black,line width=.4pt] (143,{444+147.0*\j})--(287,{444+147.0*\j});}\end{scope}
\begin{scope}[scale=.6666667,draw=black,text=black]\draw[draw=black,line width=.35pt] (389,444) rectangle (533,590);\foreach \i in {1,...,9}{\draw[draw=black,line width=.35pt] ({389+14.4*\i},444)--({389+14.4*\i},590);}\foreach \j in {1,...,9}{\draw[draw=black,line width=.35pt] (389,{444+14.6*\j})--(533,{444+14.6*\j});}\end{scope}
\begin{scope}[scale=.6666667,draw=black,text=black]\draw[draw=black,line width=.35pt] (636,444) rectangle (782,590);\foreach \i in {1,...,9}{\draw[draw=black,line width=.35pt] ({636+14.6*\i},444)--({636+14.6*\i},590);}\foreach \j in {1,...,9}{\draw[draw=black,line width=.35pt] (636,{444+14.6*\j})--(782,{444+14.6*\j});}\end{scope}
\begin{scope}[scale=.6666667,draw=black,text=black]\node[inner sep=0pt,font=\normalsize,text=black] at (123.000,451.000) {①};\end{scope}
\begin{scope}[scale=.6666667,draw=black,text=black]\node[inner sep=0pt,font=\normalsize,text=black] at (369.000,451.000) {②};\end{scope}
\begin{scope}[scale=.6666667,draw=black,text=black]\node[inner sep=0pt,font=\normalsize,text=black] at (616.000,451.000) {③};\end{scope}
\begin{scope}[scale=.6666667,draw=black,text=black]\node[inner sep=0pt,font=\normalsize,text=black] at (215.000,612.000) {$0.9$};\end{scope}
\begin{scope}[scale=.6666667,draw=black,text=black]\node[inner sep=0pt,font=\normalsize,text=black] at (461.000,612.000) {$0.09$};\end{scope}
\begin{scope}[scale=.6666667,draw=black,text=black]\node[inner sep=0pt,font=\normalsize,text=black] at (709.000,612.000) {$0.53$};\end{scope}
\hdsection{3}{657}{照样子填一填。}
\begin{scope}[scale=.6666667,draw=black,text=black]\draw[draw=black,line width=.55pt] (210.000,705.000)--(729.000,705.000)--(724.000,721.000)--(733.000,734.000)--(729.000,754.000)--(210.000,754.000)--cycle;\end{scope}
\begin{scope}[scale=.6666667,draw=black,text=black]\draw[draw=black,line width=.35pt] (219.000,705.000)--(219.000,719.000);\end{scope}
\begin{scope}[scale=.6666667,draw=black,text=black]\draw[draw=black,line width=.35pt] (223.767,705.000)--(223.767,711.000);\end{scope}
\begin{scope}[scale=.6666667,draw=black,text=black]\draw[draw=black,line width=.35pt] (228.533,705.000)--(228.533,711.000);\end{scope}
\begin{scope}[scale=.6666667,draw=black,text=black]\draw[draw=black,line width=.35pt] (233.300,705.000)--(233.300,711.000);\end{scope}
\begin{scope}[scale=.6666667,draw=black,text=black]\draw[draw=black,line width=.35pt] (238.067,705.000)--(238.067,711.000);\end{scope}
\begin{scope}[scale=.6666667,draw=black,text=black]\draw[draw=black,line width=.35pt] (242.833,705.000)--(242.833,714.000);\end{scope}
\begin{scope}[scale=.6666667,draw=black,text=black]\draw[draw=black,line width=.35pt] (247.600,705.000)--(247.600,711.000);\end{scope}
\begin{scope}[scale=.6666667,draw=black,text=black]\draw[draw=black,line width=.35pt] (252.367,705.000)--(252.367,711.000);\end{scope}
\begin{scope}[scale=.6666667,draw=black,text=black]\draw[draw=black,line width=.35pt] (257.133,705.000)--(257.133,711.000);\end{scope}
\begin{scope}[scale=.6666667,draw=black,text=black]\draw[draw=black,line width=.35pt] (261.900,705.000)--(261.900,711.000);\end{scope}
\begin{scope}[scale=.6666667,draw=black,text=black]\draw[draw=black,line width=.35pt] (266.667,705.000)--(266.667,719.000);\end{scope}
\begin{scope}[scale=.6666667,draw=black,text=black]\draw[draw=black,line width=.35pt] (271.433,705.000)--(271.433,711.000);\end{scope}
\begin{scope}[scale=.6666667,draw=black,text=black]\draw[draw=black,line width=.35pt] (276.200,705.000)--(276.200,711.000);\end{scope}
\begin{scope}[scale=.6666667,draw=black,text=black]\draw[draw=black,line width=.35pt] (280.967,705.000)--(280.967,711.000);\end{scope}
\begin{scope}[scale=.6666667,draw=black,text=black]\draw[draw=black,line width=.35pt] (285.733,705.000)--(285.733,711.000);\end{scope}
\begin{scope}[scale=.6666667,draw=black,text=black]\draw[draw=black,line width=.35pt] (290.500,705.000)--(290.500,714.000);\end{scope}
\begin{scope}[scale=.6666667,draw=black,text=black]\draw[draw=black,line width=.35pt] (295.267,705.000)--(295.267,711.000);\end{scope}
\begin{scope}[scale=.6666667,draw=black,text=black]\draw[draw=black,line width=.35pt] (300.033,705.000)--(300.033,711.000);\end{scope}
\begin{scope}[scale=.6666667,draw=black,text=black]\draw[draw=black,line width=.35pt] (304.800,705.000)--(304.800,711.000);\end{scope}
\begin{scope}[scale=.6666667,draw=black,text=black]\draw[draw=black,line width=.35pt] (309.567,705.000)--(309.567,711.000);\end{scope}
\begin{scope}[scale=.6666667,draw=black,text=black]\draw[draw=black,line width=.35pt] (314.333,705.000)--(314.333,719.000);\end{scope}
\begin{scope}[scale=.6666667,draw=black,text=black]\draw[draw=black,line width=.35pt] (319.100,705.000)--(319.100,711.000);\end{scope}
\begin{scope}[scale=.6666667,draw=black,text=black]\draw[draw=black,line width=.35pt] (323.867,705.000)--(323.867,711.000);\end{scope}
\begin{scope}[scale=.6666667,draw=black,text=black]\draw[draw=black,line width=.35pt] (328.633,705.000)--(328.633,711.000);\end{scope}
\begin{scope}[scale=.6666667,draw=black,text=black]\draw[draw=black,line width=.35pt] (333.400,705.000)--(333.400,711.000);\end{scope}
\begin{scope}[scale=.6666667,draw=black,text=black]\draw[draw=black,line width=.35pt] (338.167,705.000)--(338.167,714.000);\end{scope}
\begin{scope}[scale=.6666667,draw=black,text=black]\draw[draw=black,line width=.35pt] (342.933,705.000)--(342.933,711.000);\end{scope}
\begin{scope}[scale=.6666667,draw=black,text=black]\draw[draw=black,line width=.35pt] (347.700,705.000)--(347.700,711.000);\end{scope}
\begin{scope}[scale=.6666667,draw=black,text=black]\draw[draw=black,line width=.35pt] (352.467,705.000)--(352.467,711.000);\end{scope}
\begin{scope}[scale=.6666667,draw=black,text=black]\draw[draw=black,line width=.35pt] (357.233,705.000)--(357.233,711.000);\end{scope}
\begin{scope}[scale=.6666667,draw=black,text=black]\draw[draw=black,line width=.35pt] (362.000,705.000)--(362.000,719.000);\end{scope}
\begin{scope}[scale=.6666667,draw=black,text=black]\draw[draw=black,line width=.35pt] (366.767,705.000)--(366.767,711.000);\end{scope}
\begin{scope}[scale=.6666667,draw=black,text=black]\draw[draw=black,line width=.35pt] (371.533,705.000)--(371.533,711.000);\end{scope}
\begin{scope}[scale=.6666667,draw=black,text=black]\draw[draw=black,line width=.35pt] (376.300,705.000)--(376.300,711.000);\end{scope}
\begin{scope}[scale=.6666667,draw=black,text=black]\draw[draw=black,line width=.35pt] (381.067,705.000)--(381.067,711.000);\end{scope}
\begin{scope}[scale=.6666667,draw=black,text=black]\draw[draw=black,line width=.35pt] (385.833,705.000)--(385.833,714.000);\end{scope}
\begin{scope}[scale=.6666667,draw=black,text=black]\draw[draw=black,line width=.35pt] (390.600,705.000)--(390.600,711.000);\end{scope}
\begin{scope}[scale=.6666667,draw=black,text=black]\draw[draw=black,line width=.35pt] (395.367,705.000)--(395.367,711.000);\end{scope}
\begin{scope}[scale=.6666667,draw=black,text=black]\draw[draw=black,line width=.35pt] (400.133,705.000)--(400.133,711.000);\end{scope}
\begin{scope}[scale=.6666667,draw=black,text=black]\draw[draw=black,line width=.35pt] (404.900,705.000)--(404.900,711.000);\end{scope}
\begin{scope}[scale=.6666667,draw=black,text=black]\draw[draw=black,line width=.35pt] (409.667,705.000)--(409.667,719.000);\end{scope}
\begin{scope}[scale=.6666667,draw=black,text=black]\draw[draw=black,line width=.35pt] (414.433,705.000)--(414.433,711.000);\end{scope}
\begin{scope}[scale=.6666667,draw=black,text=black]\draw[draw=black,line width=.35pt] (419.200,705.000)--(419.200,711.000);\end{scope}
\begin{scope}[scale=.6666667,draw=black,text=black]\draw[draw=black,line width=.35pt] (423.967,705.000)--(423.967,711.000);\end{scope}
\begin{scope}[scale=.6666667,draw=black,text=black]\draw[draw=black,line width=.35pt] (428.733,705.000)--(428.733,711.000);\end{scope}
\begin{scope}[scale=.6666667,draw=black,text=black]\draw[draw=black,line width=.35pt] (433.500,705.000)--(433.500,714.000);\end{scope}
\begin{scope}[scale=.6666667,draw=black,text=black]\draw[draw=black,line width=.35pt] (438.267,705.000)--(438.267,711.000);\end{scope}
\begin{scope}[scale=.6666667,draw=black,text=black]\draw[draw=black,line width=.35pt] (443.033,705.000)--(443.033,711.000);\end{scope}
\begin{scope}[scale=.6666667,draw=black,text=black]\draw[draw=black,line width=.35pt] (447.800,705.000)--(447.800,711.000);\end{scope}
\begin{scope}[scale=.6666667,draw=black,text=black]\draw[draw=black,line width=.35pt] (452.567,705.000)--(452.567,711.000);\end{scope}
\begin{scope}[scale=.6666667,draw=black,text=black]\draw[draw=black,line width=.35pt] (457.333,705.000)--(457.333,719.000);\end{scope}
\begin{scope}[scale=.6666667,draw=black,text=black]\draw[draw=black,line width=.35pt] (462.100,705.000)--(462.100,711.000);\end{scope}
\begin{scope}[scale=.6666667,draw=black,text=black]\draw[draw=black,line width=.35pt] (466.867,705.000)--(466.867,711.000);\end{scope}
\begin{scope}[scale=.6666667,draw=black,text=black]\draw[draw=black,line width=.35pt] (471.633,705.000)--(471.633,711.000);\end{scope}
\begin{scope}[scale=.6666667,draw=black,text=black]\draw[draw=black,line width=.35pt] (476.400,705.000)--(476.400,711.000);\end{scope}
\begin{scope}[scale=.6666667,draw=black,text=black]\draw[draw=black,line width=.35pt] (481.167,705.000)--(481.167,714.000);\end{scope}
\begin{scope}[scale=.6666667,draw=black,text=black]\draw[draw=black,line width=.35pt] (485.933,705.000)--(485.933,711.000);\end{scope}
\begin{scope}[scale=.6666667,draw=black,text=black]\draw[draw=black,line width=.35pt] (490.700,705.000)--(490.700,711.000);\end{scope}
\begin{scope}[scale=.6666667,draw=black,text=black]\draw[draw=black,line width=.35pt] (495.467,705.000)--(495.467,711.000);\end{scope}
\begin{scope}[scale=.6666667,draw=black,text=black]\draw[draw=black,line width=.35pt] (500.233,705.000)--(500.233,711.000);\end{scope}
\begin{scope}[scale=.6666667,draw=black,text=black]\draw[draw=black,line width=.35pt] (505.000,705.000)--(505.000,719.000);\end{scope}
\begin{scope}[scale=.6666667,draw=black,text=black]\draw[draw=black,line width=.35pt] (509.767,705.000)--(509.767,711.000);\end{scope}
\begin{scope}[scale=.6666667,draw=black,text=black]\draw[draw=black,line width=.35pt] (514.533,705.000)--(514.533,711.000);\end{scope}
\begin{scope}[scale=.6666667,draw=black,text=black]\draw[draw=black,line width=.35pt] (519.300,705.000)--(519.300,711.000);\end{scope}
\begin{scope}[scale=.6666667,draw=black,text=black]\draw[draw=black,line width=.35pt] (524.067,705.000)--(524.067,711.000);\end{scope}
\begin{scope}[scale=.6666667,draw=black,text=black]\draw[draw=black,line width=.35pt] (528.833,705.000)--(528.833,714.000);\end{scope}
\begin{scope}[scale=.6666667,draw=black,text=black]\draw[draw=black,line width=.35pt] (533.600,705.000)--(533.600,711.000);\end{scope}
\begin{scope}[scale=.6666667,draw=black,text=black]\draw[draw=black,line width=.35pt] (538.367,705.000)--(538.367,711.000);\end{scope}
\begin{scope}[scale=.6666667,draw=black,text=black]\draw[draw=black,line width=.35pt] (543.133,705.000)--(543.133,711.000);\end{scope}
\begin{scope}[scale=.6666667,draw=black,text=black]\draw[draw=black,line width=.35pt] (547.900,705.000)--(547.900,711.000);\end{scope}
\begin{scope}[scale=.6666667,draw=black,text=black]\draw[draw=black,line width=.35pt] (552.667,705.000)--(552.667,719.000);\end{scope}
\begin{scope}[scale=.6666667,draw=black,text=black]\draw[draw=black,line width=.35pt] (557.433,705.000)--(557.433,711.000);\end{scope}
\begin{scope}[scale=.6666667,draw=black,text=black]\draw[draw=black,line width=.35pt] (562.200,705.000)--(562.200,711.000);\end{scope}
\begin{scope}[scale=.6666667,draw=black,text=black]\draw[draw=black,line width=.35pt] (566.967,705.000)--(566.967,711.000);\end{scope}
\begin{scope}[scale=.6666667,draw=black,text=black]\draw[draw=black,line width=.35pt] (571.733,705.000)--(571.733,711.000);\end{scope}
\begin{scope}[scale=.6666667,draw=black,text=black]\draw[draw=black,line width=.35pt] (576.500,705.000)--(576.500,714.000);\end{scope}
\begin{scope}[scale=.6666667,draw=black,text=black]\draw[draw=black,line width=.35pt] (581.267,705.000)--(581.267,711.000);\end{scope}
\begin{scope}[scale=.6666667,draw=black,text=black]\draw[draw=black,line width=.35pt] (586.033,705.000)--(586.033,711.000);\end{scope}
\begin{scope}[scale=.6666667,draw=black,text=black]\draw[draw=black,line width=.35pt] (590.800,705.000)--(590.800,711.000);\end{scope}
\begin{scope}[scale=.6666667,draw=black,text=black]\draw[draw=black,line width=.35pt] (595.567,705.000)--(595.567,711.000);\end{scope}
\begin{scope}[scale=.6666667,draw=black,text=black]\draw[draw=black,line width=.35pt] (600.333,705.000)--(600.333,719.000);\end{scope}
\begin{scope}[scale=.6666667,draw=black,text=black]\draw[draw=black,line width=.35pt] (605.100,705.000)--(605.100,711.000);\end{scope}
\begin{scope}[scale=.6666667,draw=black,text=black]\draw[draw=black,line width=.35pt] (609.867,705.000)--(609.867,711.000);\end{scope}
\begin{scope}[scale=.6666667,draw=black,text=black]\draw[draw=black,line width=.35pt] (614.633,705.000)--(614.633,711.000);\end{scope}
\begin{scope}[scale=.6666667,draw=black,text=black]\draw[draw=black,line width=.35pt] (619.400,705.000)--(619.400,711.000);\end{scope}
\begin{scope}[scale=.6666667,draw=black,text=black]\draw[draw=black,line width=.35pt] (624.167,705.000)--(624.167,714.000);\end{scope}
\begin{scope}[scale=.6666667,draw=black,text=black]\draw[draw=black,line width=.35pt] (628.933,705.000)--(628.933,711.000);\end{scope}
\begin{scope}[scale=.6666667,draw=black,text=black]\draw[draw=black,line width=.35pt] (633.700,705.000)--(633.700,711.000);\end{scope}
\begin{scope}[scale=.6666667,draw=black,text=black]\draw[draw=black,line width=.35pt] (638.467,705.000)--(638.467,711.000);\end{scope}
\begin{scope}[scale=.6666667,draw=black,text=black]\draw[draw=black,line width=.35pt] (643.233,705.000)--(643.233,711.000);\end{scope}
\begin{scope}[scale=.6666667,draw=black,text=black]\draw[draw=black,line width=.35pt] (648.000,705.000)--(648.000,719.000);\end{scope}
\begin{scope}[scale=.6666667,draw=black,text=black]\draw[draw=black,line width=.35pt] (652.767,705.000)--(652.767,711.000);\end{scope}
\begin{scope}[scale=.6666667,draw=black,text=black]\draw[draw=black,line width=.35pt] (657.533,705.000)--(657.533,711.000);\end{scope}
\begin{scope}[scale=.6666667,draw=black,text=black]\draw[draw=black,line width=.35pt] (662.300,705.000)--(662.300,711.000);\end{scope}
\begin{scope}[scale=.6666667,draw=black,text=black]\draw[draw=black,line width=.35pt] (667.067,705.000)--(667.067,711.000);\end{scope}
\begin{scope}[scale=.6666667,draw=black,text=black]\draw[draw=black,line width=.35pt] (671.833,705.000)--(671.833,714.000);\end{scope}
\begin{scope}[scale=.6666667,draw=black,text=black]\draw[draw=black,line width=.35pt] (676.600,705.000)--(676.600,711.000);\end{scope}
\begin{scope}[scale=.6666667,draw=black,text=black]\draw[draw=black,line width=.35pt] (681.367,705.000)--(681.367,711.000);\end{scope}
\begin{scope}[scale=.6666667,draw=black,text=black]\draw[draw=black,line width=.35pt] (686.133,705.000)--(686.133,711.000);\end{scope}
\begin{scope}[scale=.6666667,draw=black,text=black]\draw[draw=black,line width=.35pt] (690.900,705.000)--(690.900,711.000);\end{scope}
\begin{scope}[scale=.6666667,draw=black,text=black]\draw[draw=black,line width=.35pt] (695.667,705.000)--(695.667,719.000);\end{scope}
\begin{scope}[scale=.6666667,draw=black,text=black]\node[inner sep=0pt,font=\normalsize,text=black] at (219.000,738.000) {0};\end{scope}
\begin{scope}[scale=.6666667,draw=black,text=black]\node[inner sep=0pt,font=\normalsize,text=black] at (266.667,738.000) {1 cm};\end{scope}
\begin{scope}[scale=.6666667,draw=black,text=black]\node[inner sep=0pt,font=\normalsize,text=black] at (314.333,738.000) {2};\end{scope}
\begin{scope}[scale=.6666667,draw=black,text=black]\node[inner sep=0pt,font=\normalsize,text=black] at (362.000,738.000) {3};\end{scope}
\begin{scope}[scale=.6666667,draw=black,text=black]\node[inner sep=0pt,font=\normalsize,text=black] at (409.667,738.000) {4};\end{scope}
\begin{scope}[scale=.6666667,draw=black,text=black]\node[inner sep=0pt,font=\normalsize,text=black] at (457.333,738.000) {5};\end{scope}
\begin{scope}[scale=.6666667,draw=black,text=black]\node[inner sep=0pt,font=\normalsize,text=black] at (505.000,738.000) {6};\end{scope}
\begin{scope}[scale=.6666667,draw=black,text=black]\node[inner sep=0pt,font=\normalsize,text=black] at (552.667,738.000) {7};\end{scope}
\begin{scope}[scale=.6666667,draw=black,text=black]\node[inner sep=0pt,font=\normalsize,text=black] at (600.333,738.000) {8};\end{scope}
\begin{scope}[scale=.6666667,draw=black,text=black]\node[inner sep=0pt,font=\normalsize,text=black] at (648.000,738.000) {9};\end{scope}
\begin{scope}[scale=.6666667,draw=black,text=black]\node[inner sep=0pt,font=\normalsize,text=black] at (695.667,738.000) {10};\end{scope}
\begin{scope}[scale=.6666667,draw=black,text=black]\draw[draw=black,line width=.45pt,-{Stealth[length=3pt,width=2.4pt]}] (267.000,760.000)--(267.000,785.000);\end{scope}
\begin{scope}[scale=.6666667,draw=black,text=black]\draw[draw=black,line width=.45pt,-{Stealth[length=3pt,width=2.4pt]}] (362.000,760.000)--(362.000,785.000);\end{scope}
\begin{scope}[scale=.6666667,draw=black,text=black]\draw[draw=black,line width=.45pt,-{Stealth[length=3pt,width=2.4pt]}] (505.000,760.000)--(505.000,785.000);\end{scope}
\q{230}{790}{85}{12.5}{$\dfrac1{10}$分米}
\q{331}{790}{96}{12.5}{$\dfrac3{10}$分米}
\q{463}{804}{28}{13}{（}
\q{512}{804}{25}{13}{）}
\q{533}{804}{80}{12.5}{分米}
\q{225}{862}{103}{12.5}{0.1分米}
\q{324}{861}{28}{13}{（}
\q{378}{861}{25}{13}{）}
\q{399}{861}{76}{12.5}{分米}
\q{463}{861}{28}{13}{（}
\q{514}{861}{25}{13}{）}
\q{535}{861}{75}{12.5}{分米}
\hdsection{4}{921}{在计数器上表示下面各小数。}
\q{112}{967}{172}{13}{① 0.35}
\q{368}{967}{172}{13}{② 4.3}
\q{623}{967}{172}{13}{③ 0.732}
\begin{scope}[scale=.6666667,draw=black,text=black]\draw[line width=.5pt] (143,1015) rectangle (287,1079);\draw[line width=.5pt] (143,1019)--(287,1019);\end{scope}
\begin{scope}[scale=.6666667,draw=black,text=black]\draw[line width=.45pt] (165,1017)--(165,1079);\draw[line width=.45pt] (157,1015)..controls(166,995) and (171,1003)..(170,1015);\end{scope}
\begin{scope}[scale=.6666667,draw=black,text=black]\draw[line width=.45pt] (191,1017)--(191,1079);\draw[line width=.45pt] (183,1015)..controls(192,995) and (197,1003)..(196,1015);\end{scope}
\begin{scope}[scale=.6666667,draw=black,text=black]\draw[line width=.45pt] (224,1017)--(224,1079);\draw[line width=.45pt] (216,1015)..controls(225,995) and (230,1003)..(229,1015);\end{scope}
\begin{scope}[scale=.6666667,draw=black,text=black]\draw[line width=.45pt] (247,1017)--(247,1079);\draw[line width=.45pt] (239,1015)..controls(248,995) and (253,1003)..(252,1015);\end{scope}
\begin{scope}[scale=.6666667,draw=black,text=black]\draw[line width=.45pt] (273,1017)--(273,1079);\draw[line width=.45pt] (265,1015)..controls(274,995) and (279,1003)..(278,1015);\end{scope}
\begin{scope}[scale=.6666667,draw=black,text=black]\draw[draw=black,line width=.5pt,fill=white] (143.000,1078.000)--(151.000,1083.000)--(294.000,1083.000)--(287.000,1078.000)--cycle;\end{scope}
\begin{scope}[scale=.6666667,draw=black,text=black]\draw[draw=black,line width=.5pt,fill=white] (151.000,1083.000)--(294.000,1083.000)--(294.000,1147.000)--(151.000,1147.000)--cycle;\end{scope}
\begin{scope}[scale=.6666667,draw=black,text=black]\draw[draw=black,line width=.5pt,fill=white] (143.000,1078.000)--(151.000,1083.000)--(151.000,1147.000)--(143.000,1131.000)--cycle;\end{scope}
\begin{scope}[scale=.6666667,draw=black,text=black]\node[inner sep=0pt,font=\normalsize,text=black] at (165.000,1115.000) {\shortstack{十\\位}};\end{scope}
\begin{scope}[scale=.6666667,draw=black,text=black]\node[inner sep=0pt,font=\normalsize,text=black] at (191.000,1115.000) {\shortstack{个\\位}};\end{scope}
\begin{scope}[scale=.6666667,draw=black,text=black]\node[inner sep=0pt,font=\normalsize,text=black] at (224.000,1115.000) {\shortstack{十\\分\\位}};\end{scope}
\begin{scope}[scale=.6666667,draw=black,text=black]\node[inner sep=0pt,font=\normalsize,text=black] at (247.000,1115.000) {\shortstack{百\\分\\位}};\end{scope}
\begin{scope}[scale=.6666667,draw=black,text=black]\node[inner sep=0pt,font=\normalsize,text=black] at (273.000,1115.000) {\shortstack{千\\分\\位}};\end{scope}
\begin{scope}[scale=.6666667,draw=black,text=black]\draw[line width=.5pt] (398,1015) rectangle (542,1079);\draw[line width=.5pt] (398,1019)--(542,1019);\end{scope}
\begin{scope}[scale=.6666667,draw=black,text=black]\draw[line width=.45pt] (416,1017)--(416,1079);\draw[line width=.45pt] (408,1015)..controls(417,995) and (422,1003)..(421,1015);\end{scope}
\begin{scope}[scale=.6666667,draw=black,text=black]\draw[line width=.45pt] (441,1017)--(441,1079);\draw[line width=.45pt] (433,1015)..controls(442,995) and (447,1003)..(446,1015);\end{scope}
\begin{scope}[scale=.6666667,draw=black,text=black]\draw[line width=.45pt] (477,1017)--(477,1079);\draw[line width=.45pt] (469,1015)..controls(478,995) and (483,1003)..(482,1015);\end{scope}
\begin{scope}[scale=.6666667,draw=black,text=black]\draw[line width=.45pt] (502,1017)--(502,1079);\draw[line width=.45pt] (494,1015)..controls(503,995) and (508,1003)..(507,1015);\end{scope}
\begin{scope}[scale=.6666667,draw=black,text=black]\draw[line width=.45pt] (527,1017)--(527,1079);\draw[line width=.45pt] (519,1015)..controls(528,995) and (533,1003)..(532,1015);\end{scope}
\begin{scope}[scale=.6666667,draw=black,text=black]\draw[draw=black,line width=.5pt,fill=white] (398.000,1078.000)--(406.000,1083.000)--(549.000,1083.000)--(542.000,1078.000)--cycle;\end{scope}
\begin{scope}[scale=.6666667,draw=black,text=black]\draw[draw=black,line width=.5pt,fill=white] (406.000,1083.000)--(549.000,1083.000)--(549.000,1147.000)--(406.000,1147.000)--cycle;\end{scope}
\begin{scope}[scale=.6666667,draw=black,text=black]\draw[draw=black,line width=.5pt,fill=white] (398.000,1078.000)--(406.000,1083.000)--(406.000,1147.000)--(398.000,1131.000)--cycle;\end{scope}
\begin{scope}[scale=.6666667,draw=black,text=black]\node[inner sep=0pt,font=\normalsize,text=black] at (416.000,1115.000) {\shortstack{十\\位}};\end{scope}
\begin{scope}[scale=.6666667,draw=black,text=black]\node[inner sep=0pt,font=\normalsize,text=black] at (441.000,1115.000) {\shortstack{个\\位}};\end{scope}
\begin{scope}[scale=.6666667,draw=black,text=black]\node[inner sep=0pt,font=\normalsize,text=black] at (477.000,1115.000) {\shortstack{十\\分\\位}};\end{scope}
\begin{scope}[scale=.6666667,draw=black,text=black]\node[inner sep=0pt,font=\normalsize,text=black] at (502.000,1115.000) {\shortstack{百\\分\\位}};\end{scope}
\begin{scope}[scale=.6666667,draw=black,text=black]\node[inner sep=0pt,font=\normalsize,text=black] at (527.000,1115.000) {\shortstack{千\\分\\位}};\end{scope}
\begin{scope}[scale=.6666667,draw=black,text=black]\draw[line width=.5pt] (654,1015) rectangle (798,1079);\draw[line width=.5pt] (654,1019)--(798,1019);\end{scope}
\begin{scope}[scale=.6666667,draw=black,text=black]\draw[line width=.45pt] (675,1017)--(675,1079);\draw[line width=.45pt] (667,1015)..controls(676,995) and (681,1003)..(680,1015);\end{scope}
\begin{scope}[scale=.6666667,draw=black,text=black]\draw[line width=.45pt] (701,1017)--(701,1079);\draw[line width=.45pt] (693,1015)..controls(702,995) and (707,1003)..(706,1015);\end{scope}
\begin{scope}[scale=.6666667,draw=black,text=black]\draw[line width=.45pt] (731,1017)--(731,1079);\draw[line width=.45pt] (723,1015)..controls(732,995) and (737,1003)..(736,1015);\end{scope}
\begin{scope}[scale=.6666667,draw=black,text=black]\draw[line width=.45pt] (757,1017)--(757,1079);\draw[line width=.45pt] (749,1015)..controls(758,995) and (763,1003)..(762,1015);\end{scope}
\begin{scope}[scale=.6666667,draw=black,text=black]\draw[line width=.45pt] (781,1017)--(781,1079);\draw[line width=.45pt] (773,1015)..controls(782,995) and (787,1003)..(786,1015);\end{scope}
\begin{scope}[scale=.6666667,draw=black,text=black]\draw[draw=black,line width=.5pt,fill=white] (654.000,1078.000)--(662.000,1083.000)--(805.000,1083.000)--(798.000,1078.000)--cycle;\end{scope}
\begin{scope}[scale=.6666667,draw=black,text=black]\draw[draw=black,line width=.5pt,fill=white] (662.000,1083.000)--(805.000,1083.000)--(805.000,1147.000)--(662.000,1147.000)--cycle;\end{scope}
\begin{scope}[scale=.6666667,draw=black,text=black]\draw[draw=black,line width=.5pt,fill=white] (654.000,1078.000)--(662.000,1083.000)--(662.000,1147.000)--(654.000,1131.000)--cycle;\end{scope}
\begin{scope}[scale=.6666667,draw=black,text=black]\node[inner sep=0pt,font=\normalsize,text=black] at (675.000,1115.000) {\shortstack{十\\位}};\end{scope}
\begin{scope}[scale=.6666667,draw=black,text=black]\node[inner sep=0pt,font=\normalsize,text=black] at (701.000,1115.000) {\shortstack{个\\位}};\end{scope}
\begin{scope}[scale=.6666667,draw=black,text=black]\node[inner sep=0pt,font=\normalsize,text=black] at (731.000,1115.000) {\shortstack{十\\分\\位}};\end{scope}
\begin{scope}[scale=.6666667,draw=black,text=black]\node[inner sep=0pt,font=\normalsize,text=black] at (757.000,1115.000) {\shortstack{百\\分\\位}};\end{scope}
\begin{scope}[scale=.6666667,draw=black,text=black]\node[inner sep=0pt,font=\normalsize,text=black] at (781.000,1115.000) {\shortstack{千\\分\\位}};\end{scope}

\begin{scope}[color=answerred,draw=answerred,text=answerred]\ans{128.00}{226.67}{33.33}{10}{$\vphantom{0}0.4$}
\ans{294.00}{226.67}{33.33}{10}{$\vphantom{0}0.33$}
\ans{455.33}{225.33}{33.33}{10}{$\vphantom{0}0.203$}
\fill[answerred,opacity=.28](95.333,296.000) rectangle (181.733,394.000);
\fill[answerred,opacity=.28](259.333,296.000) rectangle (345.733,305.733);
\fill[answerred,opacity=.28](424.000,296.000) rectangle (472.667,393.333);
\fill[answerred,opacity=.28](472.667,296.000) rectangle (482.400,325.200);
\ans{228.67}{573.33}{48.00}{11}{$\vphantom{0}0.3$}
\ans{322.67}{534.00}{42.00}{11}{$\vphantom{0}\frac6{10}$}
\ans{321.33}{573.33}{48.00}{11}{$\vphantom{0}0.6$}
\filldraw[answerred,fill=answerred!65] (149.333,714.000) ellipse [x radius=4bp,y radius=2.2bp];
\filldraw[answerred,fill=answerred!65] (149.333,708.667) ellipse [x radius=4bp,y radius=2.2bp];
\filldraw[answerred,fill=answerred!65] (149.333,703.333) ellipse [x radius=4bp,y radius=2.2bp];
\filldraw[answerred,fill=answerred!65] (164.667,714.000) ellipse [x radius=4bp,y radius=2.2bp];
\filldraw[answerred,fill=answerred!65] (164.667,708.667) ellipse [x radius=4bp,y radius=2.2bp];
\filldraw[answerred,fill=answerred!65] (164.667,703.333) ellipse [x radius=4bp,y radius=2.2bp];
\filldraw[answerred,fill=answerred!65] (164.667,698.000) ellipse [x radius=4bp,y radius=2.2bp];
\filldraw[answerred,fill=answerred!65] (164.667,692.667) ellipse [x radius=4bp,y radius=2.2bp];
\filldraw[answerred,fill=answerred!65] (294.000,714.000) ellipse [x radius=4bp,y radius=2.2bp];
\filldraw[answerred,fill=answerred!65] (294.000,708.667) ellipse [x radius=4bp,y radius=2.2bp];
\filldraw[answerred,fill=answerred!65] (294.000,703.333) ellipse [x radius=4bp,y radius=2.2bp];
\filldraw[answerred,fill=answerred!65] (294.000,698.000) ellipse [x radius=4bp,y radius=2.2bp];
\filldraw[answerred,fill=answerred!65] (318.000,714.000) ellipse [x radius=4bp,y radius=2.2bp];
\filldraw[answerred,fill=answerred!65] (318.000,708.667) ellipse [x radius=4bp,y radius=2.2bp];
\filldraw[answerred,fill=answerred!65] (318.000,703.333) ellipse [x radius=4bp,y radius=2.2bp];
\filldraw[answerred,fill=answerred!65] (487.333,714.000) ellipse [x radius=4bp,y radius=2.2bp];
\filldraw[answerred,fill=answerred!65] (487.333,708.667) ellipse [x radius=4bp,y radius=2.2bp];
\filldraw[answerred,fill=answerred!65] (487.333,703.333) ellipse [x radius=4bp,y radius=2.2bp];
\filldraw[answerred,fill=answerred!65] (487.333,698.000) ellipse [x radius=4bp,y radius=2.2bp];
\filldraw[answerred,fill=answerred!65] (487.333,692.667) ellipse [x radius=4bp,y radius=2.2bp];
\filldraw[answerred,fill=answerred!65] (487.333,687.333) ellipse [x radius=4bp,y radius=2.2bp];
\filldraw[answerred,fill=answerred!65] (487.333,682.000) ellipse [x radius=4bp,y radius=2.2bp];
\filldraw[answerred,fill=answerred!65] (504.667,714.000) ellipse [x radius=4bp,y radius=2.2bp];
\filldraw[answerred,fill=answerred!65] (504.667,708.667) ellipse [x radius=4bp,y radius=2.2bp];
\filldraw[answerred,fill=answerred!65] (504.667,703.333) ellipse [x radius=4bp,y radius=2.2bp];
\filldraw[answerred,fill=answerred!65] (520.667,714.000) ellipse [x radius=4bp,y radius=2.2bp];
\filldraw[answerred,fill=answerred!65] (520.667,708.667) ellipse [x radius=4bp,y radius=2.2bp];
\end{scope}
\end{scope}\pgfresetboundingbox\path[use as bounding box] (56.0,464.0) rectangle (494.0,612.0);
\end{tikzpicture}}\par



\par

\par\noindent

\QHeading{第4题　在计数器上表示下面各小数。}

\par\noindent\makebox[\linewidth][c]{
\begin{tikzpicture}[x=1bp,y=-1bp,color=black,draw=black,text=black]
\begin{scope}\clip (69.0,639.0) rectangle (542.0,770.0);
\hdhead{46}{小数的意义与性质——信息窗1}
\hdsection{1}{135}{用小数表示下图的涂色部分。}
\begin{scope}[scale=.6666667,draw=black,text=black]\path[pattern=north east lines,pattern color=black] (143,178) rectangle (200.6,323);\end{scope}
\begin{scope}[scale=.6666667,draw=black,text=black]\draw[draw=black,line width=.4pt] (143,178) rectangle (287,323);\foreach \i in {1,...,9}{\draw[draw=black,line width=.4pt] ({143+14.4*\i},178)--({143+14.4*\i},323);}\foreach \j in {1,...,0}{\draw[draw=black,line width=.4pt] (143,{178+145.0*\j})--(287,{178+145.0*\j});}\end{scope}
\begin{scope}[scale=.6666667,draw=black,text=black]\path[pattern=north east lines,pattern color=black] (389,178) rectangle (432.2,322);\end{scope}
\begin{scope}[scale=.6666667,draw=black,text=black]\path[pattern=north east lines,pattern color=black] (432.2,278.8) rectangle (446.59999999999997,322.0);\end{scope}
\begin{scope}[scale=.6666667,draw=black,text=black]\draw[draw=black,line width=.35pt] (389,178) rectangle (533,322);\foreach \i in {1,...,9}{\draw[draw=black,line width=.35pt] ({389+14.4*\i},178)--({389+14.4*\i},322);}\foreach \j in {1,...,9}{\draw[draw=black,line width=.35pt] (389,{178+14.4*\j})--(533,{178+14.4*\j});}\end{scope}
\begin{scope}[scale=.6666667,draw=black,text=black]\fill[white] (636.0000,319.0000)--(736.0000,319.0000)--(736.0000,219.0000)--(636.0000,219.0000)--cycle;\end{scope}
\begin{scope}[scale=.6666667,draw=black,text=black]\fill[white] (636.0000,219.0000)--(736.0000,219.0000)--(769.0000,179.0000)--(669.0000,179.0000)--cycle;\end{scope}
\begin{scope}[scale=.6666667,draw=black,text=black]\path[pattern=north east lines,pattern color=black] (736.0000,319.0000)--(769.0000,279.0000)--(769.0000,179.0000)--(736.0000,219.0000)--cycle;\end{scope}
\begin{scope}[scale=.6666667,draw=black,text=black]\path[pattern=north east lines,pattern color=black] (706.0000,249.0000)--(716.0000,249.0000)--(716.0000,239.0000)--(706.0000,239.0000)--cycle;\end{scope}
\begin{scope}[scale=.6666667,draw=black,text=black]\path[pattern=north east lines,pattern color=black] (706.0000,239.0000)--(716.0000,239.0000)--(716.0000,229.0000)--(706.0000,229.0000)--cycle;\end{scope}
\begin{scope}[scale=.6666667,draw=black,text=black]\path[pattern=north east lines,pattern color=black] (706.0000,229.0000)--(716.0000,229.0000)--(716.0000,219.0000)--(706.0000,219.0000)--cycle;\end{scope}
\begin{scope}[scale=.6666667,draw=black,text=black]\path[pattern=north east lines,pattern color=black] (706.0000,219.0000)--(716.0000,219.0000)--(719.3000,215.0000)--(709.3000,215.0000)--cycle;\end{scope}
\begin{scope}[scale=.6666667,draw=black,text=black]\path[pattern=north east lines,pattern color=black] (716.0000,319.0000)--(726.0000,319.0000)--(726.0000,309.0000)--(716.0000,309.0000)--cycle;\end{scope}
\begin{scope}[scale=.6666667,draw=black,text=black]\path[pattern=north east lines,pattern color=black] (716.0000,309.0000)--(726.0000,309.0000)--(726.0000,299.0000)--(716.0000,299.0000)--cycle;\end{scope}
\begin{scope}[scale=.6666667,draw=black,text=black]\path[pattern=north east lines,pattern color=black] (716.0000,299.0000)--(726.0000,299.0000)--(726.0000,289.0000)--(716.0000,289.0000)--cycle;\end{scope}
\begin{scope}[scale=.6666667,draw=black,text=black]\path[pattern=north east lines,pattern color=black] (716.0000,289.0000)--(726.0000,289.0000)--(726.0000,279.0000)--(716.0000,279.0000)--cycle;\end{scope}
\begin{scope}[scale=.6666667,draw=black,text=black]\path[pattern=north east lines,pattern color=black] (716.0000,279.0000)--(726.0000,279.0000)--(726.0000,269.0000)--(716.0000,269.0000)--cycle;\end{scope}
\begin{scope}[scale=.6666667,draw=black,text=black]\path[pattern=north east lines,pattern color=black] (716.0000,269.0000)--(726.0000,269.0000)--(726.0000,259.0000)--(716.0000,259.0000)--cycle;\end{scope}
\begin{scope}[scale=.6666667,draw=black,text=black]\path[pattern=north east lines,pattern color=black] (716.0000,259.0000)--(726.0000,259.0000)--(726.0000,249.0000)--(716.0000,249.0000)--cycle;\end{scope}
\begin{scope}[scale=.6666667,draw=black,text=black]\path[pattern=north east lines,pattern color=black] (716.0000,249.0000)--(726.0000,249.0000)--(726.0000,239.0000)--(716.0000,239.0000)--cycle;\end{scope}
\begin{scope}[scale=.6666667,draw=black,text=black]\path[pattern=north east lines,pattern color=black] (716.0000,239.0000)--(726.0000,239.0000)--(726.0000,229.0000)--(716.0000,229.0000)--cycle;\end{scope}
\begin{scope}[scale=.6666667,draw=black,text=black]\path[pattern=north east lines,pattern color=black] (716.0000,229.0000)--(726.0000,229.0000)--(726.0000,219.0000)--(716.0000,219.0000)--cycle;\end{scope}
\begin{scope}[scale=.6666667,draw=black,text=black]\path[pattern=north east lines,pattern color=black] (716.0000,219.0000)--(726.0000,219.0000)--(729.3000,215.0000)--(719.3000,215.0000)--cycle;\end{scope}
\begin{scope}[scale=.6666667,draw=black,text=black]\path[pattern=north east lines,pattern color=black] (719.3000,215.0000)--(729.3000,215.0000)--(732.6000,211.0000)--(722.6000,211.0000)--cycle;\end{scope}
\begin{scope}[scale=.6666667,draw=black,text=black]\path[pattern=north east lines,pattern color=black] (722.6000,211.0000)--(732.6000,211.0000)--(735.9000,207.0000)--(725.9000,207.0000)--cycle;\end{scope}
\begin{scope}[scale=.6666667,draw=black,text=black]\path[pattern=north east lines,pattern color=black] (725.9000,207.0000)--(735.9000,207.0000)--(739.2000,203.0000)--(729.2000,203.0000)--cycle;\end{scope}
\begin{scope}[scale=.6666667,draw=black,text=black]\path[pattern=north east lines,pattern color=black] (729.2000,203.0000)--(739.2000,203.0000)--(742.5000,199.0000)--(732.5000,199.0000)--cycle;\end{scope}
\begin{scope}[scale=.6666667,draw=black,text=black]\path[pattern=north east lines,pattern color=black] (732.5000,199.0000)--(742.5000,199.0000)--(745.8000,195.0000)--(735.8000,195.0000)--cycle;\end{scope}
\begin{scope}[scale=.6666667,draw=black,text=black]\path[pattern=north east lines,pattern color=black] (735.8000,195.0000)--(745.8000,195.0000)--(749.1000,191.0000)--(739.1000,191.0000)--cycle;\end{scope}
\begin{scope}[scale=.6666667,draw=black,text=black]\path[pattern=north east lines,pattern color=black] (739.1000,191.0000)--(749.1000,191.0000)--(752.4000,187.0000)--(742.4000,187.0000)--cycle;\end{scope}
\begin{scope}[scale=.6666667,draw=black,text=black]\path[pattern=north east lines,pattern color=black] (742.4000,187.0000)--(752.4000,187.0000)--(755.7000,183.0000)--(745.7000,183.0000)--cycle;\end{scope}
\begin{scope}[scale=.6666667,draw=black,text=black]\path[pattern=north east lines,pattern color=black] (745.7000,183.0000)--(755.7000,183.0000)--(759.0000,179.0000)--(749.0000,179.0000)--cycle;\end{scope}
\begin{scope}[scale=.6666667,draw=black,text=black]\path[pattern=north east lines,pattern color=black] (726.0000,319.0000)--(736.0000,319.0000)--(736.0000,309.0000)--(726.0000,309.0000)--cycle;\end{scope}
\begin{scope}[scale=.6666667,draw=black,text=black]\path[pattern=north east lines,pattern color=black] (726.0000,309.0000)--(736.0000,309.0000)--(736.0000,299.0000)--(726.0000,299.0000)--cycle;\end{scope}
\begin{scope}[scale=.6666667,draw=black,text=black]\path[pattern=north east lines,pattern color=black] (726.0000,299.0000)--(736.0000,299.0000)--(736.0000,289.0000)--(726.0000,289.0000)--cycle;\end{scope}
\begin{scope}[scale=.6666667,draw=black,text=black]\path[pattern=north east lines,pattern color=black] (726.0000,289.0000)--(736.0000,289.0000)--(736.0000,279.0000)--(726.0000,279.0000)--cycle;\end{scope}
\begin{scope}[scale=.6666667,draw=black,text=black]\path[pattern=north east lines,pattern color=black] (726.0000,279.0000)--(736.0000,279.0000)--(736.0000,269.0000)--(726.0000,269.0000)--cycle;\end{scope}
\begin{scope}[scale=.6666667,draw=black,text=black]\path[pattern=north east lines,pattern color=black] (726.0000,269.0000)--(736.0000,269.0000)--(736.0000,259.0000)--(726.0000,259.0000)--cycle;\end{scope}
\begin{scope}[scale=.6666667,draw=black,text=black]\path[pattern=north east lines,pattern color=black] (726.0000,259.0000)--(736.0000,259.0000)--(736.0000,249.0000)--(726.0000,249.0000)--cycle;\end{scope}
\begin{scope}[scale=.6666667,draw=black,text=black]\path[pattern=north east lines,pattern color=black] (726.0000,249.0000)--(736.0000,249.0000)--(736.0000,239.0000)--(726.0000,239.0000)--cycle;\end{scope}
\begin{scope}[scale=.6666667,draw=black,text=black]\path[pattern=north east lines,pattern color=black] (726.0000,239.0000)--(736.0000,239.0000)--(736.0000,229.0000)--(726.0000,229.0000)--cycle;\end{scope}
\begin{scope}[scale=.6666667,draw=black,text=black]\path[pattern=north east lines,pattern color=black] (726.0000,229.0000)--(736.0000,229.0000)--(736.0000,219.0000)--(726.0000,219.0000)--cycle;\end{scope}
\begin{scope}[scale=.6666667,draw=black,text=black]\path[pattern=north east lines,pattern color=black] (726.0000,219.0000)--(736.0000,219.0000)--(739.3000,215.0000)--(729.3000,215.0000)--cycle;\end{scope}
\begin{scope}[scale=.6666667,draw=black,text=black]\path[pattern=north east lines,pattern color=black] (729.3000,215.0000)--(739.3000,215.0000)--(742.6000,211.0000)--(732.6000,211.0000)--cycle;\end{scope}
\begin{scope}[scale=.6666667,draw=black,text=black]\path[pattern=north east lines,pattern color=black] (732.6000,211.0000)--(742.6000,211.0000)--(745.9000,207.0000)--(735.9000,207.0000)--cycle;\end{scope}
\begin{scope}[scale=.6666667,draw=black,text=black]\path[pattern=north east lines,pattern color=black] (735.9000,207.0000)--(745.9000,207.0000)--(749.2000,203.0000)--(739.2000,203.0000)--cycle;\end{scope}
\begin{scope}[scale=.6666667,draw=black,text=black]\path[pattern=north east lines,pattern color=black] (739.2000,203.0000)--(749.2000,203.0000)--(752.5000,199.0000)--(742.5000,199.0000)--cycle;\end{scope}
\begin{scope}[scale=.6666667,draw=black,text=black]\path[pattern=north east lines,pattern color=black] (742.5000,199.0000)--(752.5000,199.0000)--(755.8000,195.0000)--(745.8000,195.0000)--cycle;\end{scope}
\begin{scope}[scale=.6666667,draw=black,text=black]\path[pattern=north east lines,pattern color=black] (745.8000,195.0000)--(755.8000,195.0000)--(759.1000,191.0000)--(749.1000,191.0000)--cycle;\end{scope}
\begin{scope}[scale=.6666667,draw=black,text=black]\path[pattern=north east lines,pattern color=black] (749.1000,191.0000)--(759.1000,191.0000)--(762.4000,187.0000)--(752.4000,187.0000)--cycle;\end{scope}
\begin{scope}[scale=.6666667,draw=black,text=black]\path[pattern=north east lines,pattern color=black] (752.4000,187.0000)--(762.4000,187.0000)--(765.7000,183.0000)--(755.7000,183.0000)--cycle;\end{scope}
\begin{scope}[scale=.6666667,draw=black,text=black]\path[pattern=north east lines,pattern color=black] (755.7000,183.0000)--(765.7000,183.0000)--(769.0000,179.0000)--(759.0000,179.0000)--cycle;\end{scope}
\begin{scope}[scale=.6666667,draw=black,text=black]\draw[draw=black!85,line width=.32pt] (636.000,319.000)--(636.000,219.000);\end{scope}
\begin{scope}[scale=.6666667,draw=black,text=black]\draw[draw=black!85,line width=.32pt] (636.000,319.000)--(736.000,319.000);\end{scope}
\begin{scope}[scale=.6666667,draw=black,text=black]\draw[draw=black!85,line width=.32pt] (636.000,219.000)--(669.000,179.000);\end{scope}
\begin{scope}[scale=.6666667,draw=black,text=black]\draw[draw=black!85,line width=.32pt] (636.000,219.000)--(736.000,219.000);\end{scope}
\begin{scope}[scale=.6666667,draw=black,text=black]\draw[draw=black!85,line width=.32pt] (736.000,319.000)--(736.000,219.000);\end{scope}
\begin{scope}[scale=.6666667,draw=black,text=black]\draw[draw=black!85,line width=.32pt] (736.000,319.000)--(769.000,279.000);\end{scope}
\begin{scope}[scale=.6666667,draw=black,text=black]\draw[draw=black!85,line width=.32pt] (646.000,319.000)--(646.000,219.000);\end{scope}
\begin{scope}[scale=.6666667,draw=black,text=black]\draw[draw=black!85,line width=.32pt] (636.000,309.000)--(736.000,309.000);\end{scope}
\begin{scope}[scale=.6666667,draw=black,text=black]\draw[draw=black!85,line width=.32pt] (646.000,219.000)--(679.000,179.000);\end{scope}
\begin{scope}[scale=.6666667,draw=black,text=black]\draw[draw=black!85,line width=.32pt] (639.300,215.000)--(739.300,215.000);\end{scope}
\begin{scope}[scale=.6666667,draw=black,text=black]\draw[draw=black!85,line width=.32pt] (739.300,315.000)--(739.300,215.000);\end{scope}
\begin{scope}[scale=.6666667,draw=black,text=black]\draw[draw=black!85,line width=.32pt] (736.000,309.000)--(769.000,269.000);\end{scope}
\begin{scope}[scale=.6666667,draw=black,text=black]\draw[draw=black!85,line width=.32pt] (656.000,319.000)--(656.000,219.000);\end{scope}
\begin{scope}[scale=.6666667,draw=black,text=black]\draw[draw=black!85,line width=.32pt] (636.000,299.000)--(736.000,299.000);\end{scope}
\begin{scope}[scale=.6666667,draw=black,text=black]\draw[draw=black!85,line width=.32pt] (656.000,219.000)--(689.000,179.000);\end{scope}
\begin{scope}[scale=.6666667,draw=black,text=black]\draw[draw=black!85,line width=.32pt] (642.600,211.000)--(742.600,211.000);\end{scope}
\begin{scope}[scale=.6666667,draw=black,text=black]\draw[draw=black!85,line width=.32pt] (742.600,311.000)--(742.600,211.000);\end{scope}
\begin{scope}[scale=.6666667,draw=black,text=black]\draw[draw=black!85,line width=.32pt] (736.000,299.000)--(769.000,259.000);\end{scope}
\begin{scope}[scale=.6666667,draw=black,text=black]\draw[draw=black!85,line width=.32pt] (666.000,319.000)--(666.000,219.000);\end{scope}
\begin{scope}[scale=.6666667,draw=black,text=black]\draw[draw=black!85,line width=.32pt] (636.000,289.000)--(736.000,289.000);\end{scope}
\begin{scope}[scale=.6666667,draw=black,text=black]\draw[draw=black!85,line width=.32pt] (666.000,219.000)--(699.000,179.000);\end{scope}
\begin{scope}[scale=.6666667,draw=black,text=black]\draw[draw=black!85,line width=.32pt] (645.900,207.000)--(745.900,207.000);\end{scope}
\begin{scope}[scale=.6666667,draw=black,text=black]\draw[draw=black!85,line width=.32pt] (745.900,307.000)--(745.900,207.000);\end{scope}
\begin{scope}[scale=.6666667,draw=black,text=black]\draw[draw=black!85,line width=.32pt] (736.000,289.000)--(769.000,249.000);\end{scope}
\begin{scope}[scale=.6666667,draw=black,text=black]\draw[draw=black!85,line width=.32pt] (676.000,319.000)--(676.000,219.000);\end{scope}
\begin{scope}[scale=.6666667,draw=black,text=black]\draw[draw=black!85,line width=.32pt] (636.000,279.000)--(736.000,279.000);\end{scope}
\begin{scope}[scale=.6666667,draw=black,text=black]\draw[draw=black!85,line width=.32pt] (676.000,219.000)--(709.000,179.000);\end{scope}
\begin{scope}[scale=.6666667,draw=black,text=black]\draw[draw=black!85,line width=.32pt] (649.200,203.000)--(749.200,203.000);\end{scope}
\begin{scope}[scale=.6666667,draw=black,text=black]\draw[draw=black!85,line width=.32pt] (749.200,303.000)--(749.200,203.000);\end{scope}
\begin{scope}[scale=.6666667,draw=black,text=black]\draw[draw=black!85,line width=.32pt] (736.000,279.000)--(769.000,239.000);\end{scope}
\begin{scope}[scale=.6666667,draw=black,text=black]\draw[draw=black!85,line width=.32pt] (686.000,319.000)--(686.000,219.000);\end{scope}
\begin{scope}[scale=.6666667,draw=black,text=black]\draw[draw=black!85,line width=.32pt] (636.000,269.000)--(736.000,269.000);\end{scope}
\begin{scope}[scale=.6666667,draw=black,text=black]\draw[draw=black!85,line width=.32pt] (686.000,219.000)--(719.000,179.000);\end{scope}
\begin{scope}[scale=.6666667,draw=black,text=black]\draw[draw=black!85,line width=.32pt] (652.500,199.000)--(752.500,199.000);\end{scope}
\begin{scope}[scale=.6666667,draw=black,text=black]\draw[draw=black!85,line width=.32pt] (752.500,299.000)--(752.500,199.000);\end{scope}
\begin{scope}[scale=.6666667,draw=black,text=black]\draw[draw=black!85,line width=.32pt] (736.000,269.000)--(769.000,229.000);\end{scope}
\begin{scope}[scale=.6666667,draw=black,text=black]\draw[draw=black!85,line width=.32pt] (696.000,319.000)--(696.000,219.000);\end{scope}
\begin{scope}[scale=.6666667,draw=black,text=black]\draw[draw=black!85,line width=.32pt] (636.000,259.000)--(736.000,259.000);\end{scope}
\begin{scope}[scale=.6666667,draw=black,text=black]\draw[draw=black!85,line width=.32pt] (696.000,219.000)--(729.000,179.000);\end{scope}
\begin{scope}[scale=.6666667,draw=black,text=black]\draw[draw=black!85,line width=.32pt] (655.800,195.000)--(755.800,195.000);\end{scope}
\begin{scope}[scale=.6666667,draw=black,text=black]\draw[draw=black!85,line width=.32pt] (755.800,295.000)--(755.800,195.000);\end{scope}
\begin{scope}[scale=.6666667,draw=black,text=black]\draw[draw=black!85,line width=.32pt] (736.000,259.000)--(769.000,219.000);\end{scope}
\begin{scope}[scale=.6666667,draw=black,text=black]\draw[draw=black!85,line width=.32pt] (706.000,319.000)--(706.000,219.000);\end{scope}
\begin{scope}[scale=.6666667,draw=black,text=black]\draw[draw=black!85,line width=.32pt] (636.000,249.000)--(736.000,249.000);\end{scope}
\begin{scope}[scale=.6666667,draw=black,text=black]\draw[draw=black!85,line width=.32pt] (706.000,219.000)--(739.000,179.000);\end{scope}
\begin{scope}[scale=.6666667,draw=black,text=black]\draw[draw=black!85,line width=.32pt] (659.100,191.000)--(759.100,191.000);\end{scope}
\begin{scope}[scale=.6666667,draw=black,text=black]\draw[draw=black!85,line width=.32pt] (759.100,291.000)--(759.100,191.000);\end{scope}
\begin{scope}[scale=.6666667,draw=black,text=black]\draw[draw=black!85,line width=.32pt] (736.000,249.000)--(769.000,209.000);\end{scope}
\begin{scope}[scale=.6666667,draw=black,text=black]\draw[draw=black!85,line width=.32pt] (716.000,319.000)--(716.000,219.000);\end{scope}
\begin{scope}[scale=.6666667,draw=black,text=black]\draw[draw=black!85,line width=.32pt] (636.000,239.000)--(736.000,239.000);\end{scope}
\begin{scope}[scale=.6666667,draw=black,text=black]\draw[draw=black!85,line width=.32pt] (716.000,219.000)--(749.000,179.000);\end{scope}
\begin{scope}[scale=.6666667,draw=black,text=black]\draw[draw=black!85,line width=.32pt] (662.400,187.000)--(762.400,187.000);\end{scope}
\begin{scope}[scale=.6666667,draw=black,text=black]\draw[draw=black!85,line width=.32pt] (762.400,287.000)--(762.400,187.000);\end{scope}
\begin{scope}[scale=.6666667,draw=black,text=black]\draw[draw=black!85,line width=.32pt] (736.000,239.000)--(769.000,199.000);\end{scope}
\begin{scope}[scale=.6666667,draw=black,text=black]\draw[draw=black!85,line width=.32pt] (726.000,319.000)--(726.000,219.000);\end{scope}
\begin{scope}[scale=.6666667,draw=black,text=black]\draw[draw=black!85,line width=.32pt] (636.000,229.000)--(736.000,229.000);\end{scope}
\begin{scope}[scale=.6666667,draw=black,text=black]\draw[draw=black!85,line width=.32pt] (726.000,219.000)--(759.000,179.000);\end{scope}
\begin{scope}[scale=.6666667,draw=black,text=black]\draw[draw=black!85,line width=.32pt] (665.700,183.000)--(765.700,183.000);\end{scope}
\begin{scope}[scale=.6666667,draw=black,text=black]\draw[draw=black!85,line width=.32pt] (765.700,283.000)--(765.700,183.000);\end{scope}
\begin{scope}[scale=.6666667,draw=black,text=black]\draw[draw=black!85,line width=.32pt] (736.000,229.000)--(769.000,189.000);\end{scope}
\begin{scope}[scale=.6666667,draw=black,text=black]\draw[draw=black!85,line width=.32pt] (736.000,319.000)--(736.000,219.000);\end{scope}
\begin{scope}[scale=.6666667,draw=black,text=black]\draw[draw=black!85,line width=.32pt] (636.000,219.000)--(736.000,219.000);\end{scope}
\begin{scope}[scale=.6666667,draw=black,text=black]\draw[draw=black!85,line width=.32pt] (736.000,219.000)--(769.000,179.000);\end{scope}
\begin{scope}[scale=.6666667,draw=black,text=black]\draw[draw=black!85,line width=.32pt] (669.000,179.000)--(769.000,179.000);\end{scope}
\begin{scope}[scale=.6666667,draw=black,text=black]\draw[draw=black!85,line width=.32pt] (769.000,279.000)--(769.000,179.000);\end{scope}
\begin{scope}[scale=.6666667,draw=black,text=black]\draw[draw=black!85,line width=.32pt] (736.000,219.000)--(769.000,179.000);\end{scope}
\begin{scope}[scale=.6666667,draw=black,text=black]\draw[draw=black,line width=.55pt] (636.000,319.000)--(736.000,319.000)--(736.000,219.000)--(636.000,219.000)--cycle;\end{scope}
\begin{scope}[scale=.6666667,draw=black,text=black]\draw[draw=black,line width=.55pt] (636.000,219.000)--(736.000,219.000)--(769.000,179.000)--(669.000,179.000)--cycle;\end{scope}
\begin{scope}[scale=.6666667,draw=black,text=black]\draw[draw=black,line width=.55pt] (736.000,319.000)--(769.000,279.000)--(769.000,179.000)--(736.000,219.000)--cycle;\end{scope}
\begin{scope}[scale=.6666667,draw=black,text=black]\node[inner sep=0pt,font=\normalsize,text=black] at (123.000,189.000) {①};\end{scope}
\begin{scope}[scale=.6666667,draw=black,text=black]\node[inner sep=0pt,font=\normalsize,text=black] at (370.000,189.000) {②};\end{scope}
\begin{scope}[scale=.6666667,draw=black,text=black]\node[inner sep=0pt,font=\normalsize,text=black] at (617.000,190.000) {③};\end{scope}
\q{171}{339}{28}{13}{（}
\q{247}{339}{25}{13}{）}
\q{418}{339}{28}{13}{（}
\q{493}{339}{25}{13}{）}
\q{657}{338}{28}{13}{（}
\q{733}{338}{25}{13}{）}
\hdsection{2}{395}{给图形涂色表示它下面的小数。}
\begin{scope}[scale=.6666667,draw=black,text=black]\draw[draw=black,line width=.4pt] (143,444) rectangle (287,591);\foreach \i in {1,...,9}{\draw[draw=black,line width=.4pt] ({143+14.4*\i},444)--({143+14.4*\i},591);}\foreach \j in {1,...,0}{\draw[draw=black,line width=.4pt] (143,{444+147.0*\j})--(287,{444+147.0*\j});}\end{scope}
\begin{scope}[scale=.6666667,draw=black,text=black]\draw[draw=black,line width=.35pt] (389,444) rectangle (533,590);\foreach \i in {1,...,9}{\draw[draw=black,line width=.35pt] ({389+14.4*\i},444)--({389+14.4*\i},590);}\foreach \j in {1,...,9}{\draw[draw=black,line width=.35pt] (389,{444+14.6*\j})--(533,{444+14.6*\j});}\end{scope}
\begin{scope}[scale=.6666667,draw=black,text=black]\draw[draw=black,line width=.35pt] (636,444) rectangle (782,590);\foreach \i in {1,...,9}{\draw[draw=black,line width=.35pt] ({636+14.6*\i},444)--({636+14.6*\i},590);}\foreach \j in {1,...,9}{\draw[draw=black,line width=.35pt] (636,{444+14.6*\j})--(782,{444+14.6*\j});}\end{scope}
\begin{scope}[scale=.6666667,draw=black,text=black]\node[inner sep=0pt,font=\normalsize,text=black] at (123.000,451.000) {①};\end{scope}
\begin{scope}[scale=.6666667,draw=black,text=black]\node[inner sep=0pt,font=\normalsize,text=black] at (369.000,451.000) {②};\end{scope}
\begin{scope}[scale=.6666667,draw=black,text=black]\node[inner sep=0pt,font=\normalsize,text=black] at (616.000,451.000) {③};\end{scope}
\begin{scope}[scale=.6666667,draw=black,text=black]\node[inner sep=0pt,font=\normalsize,text=black] at (215.000,612.000) {$0.9$};\end{scope}
\begin{scope}[scale=.6666667,draw=black,text=black]\node[inner sep=0pt,font=\normalsize,text=black] at (461.000,612.000) {$0.09$};\end{scope}
\begin{scope}[scale=.6666667,draw=black,text=black]\node[inner sep=0pt,font=\normalsize,text=black] at (709.000,612.000) {$0.53$};\end{scope}
\hdsection{3}{657}{照样子填一填。}
\begin{scope}[scale=.6666667,draw=black,text=black]\draw[draw=black,line width=.55pt] (210.000,705.000)--(729.000,705.000)--(724.000,721.000)--(733.000,734.000)--(729.000,754.000)--(210.000,754.000)--cycle;\end{scope}
\begin{scope}[scale=.6666667,draw=black,text=black]\draw[draw=black,line width=.35pt] (219.000,705.000)--(219.000,719.000);\end{scope}
\begin{scope}[scale=.6666667,draw=black,text=black]\draw[draw=black,line width=.35pt] (223.767,705.000)--(223.767,711.000);\end{scope}
\begin{scope}[scale=.6666667,draw=black,text=black]\draw[draw=black,line width=.35pt] (228.533,705.000)--(228.533,711.000);\end{scope}
\begin{scope}[scale=.6666667,draw=black,text=black]\draw[draw=black,line width=.35pt] (233.300,705.000)--(233.300,711.000);\end{scope}
\begin{scope}[scale=.6666667,draw=black,text=black]\draw[draw=black,line width=.35pt] (238.067,705.000)--(238.067,711.000);\end{scope}
\begin{scope}[scale=.6666667,draw=black,text=black]\draw[draw=black,line width=.35pt] (242.833,705.000)--(242.833,714.000);\end{scope}
\begin{scope}[scale=.6666667,draw=black,text=black]\draw[draw=black,line width=.35pt] (247.600,705.000)--(247.600,711.000);\end{scope}
\begin{scope}[scale=.6666667,draw=black,text=black]\draw[draw=black,line width=.35pt] (252.367,705.000)--(252.367,711.000);\end{scope}
\begin{scope}[scale=.6666667,draw=black,text=black]\draw[draw=black,line width=.35pt] (257.133,705.000)--(257.133,711.000);\end{scope}
\begin{scope}[scale=.6666667,draw=black,text=black]\draw[draw=black,line width=.35pt] (261.900,705.000)--(261.900,711.000);\end{scope}
\begin{scope}[scale=.6666667,draw=black,text=black]\draw[draw=black,line width=.35pt] (266.667,705.000)--(266.667,719.000);\end{scope}
\begin{scope}[scale=.6666667,draw=black,text=black]\draw[draw=black,line width=.35pt] (271.433,705.000)--(271.433,711.000);\end{scope}
\begin{scope}[scale=.6666667,draw=black,text=black]\draw[draw=black,line width=.35pt] (276.200,705.000)--(276.200,711.000);\end{scope}
\begin{scope}[scale=.6666667,draw=black,text=black]\draw[draw=black,line width=.35pt] (280.967,705.000)--(280.967,711.000);\end{scope}
\begin{scope}[scale=.6666667,draw=black,text=black]\draw[draw=black,line width=.35pt] (285.733,705.000)--(285.733,711.000);\end{scope}
\begin{scope}[scale=.6666667,draw=black,text=black]\draw[draw=black,line width=.35pt] (290.500,705.000)--(290.500,714.000);\end{scope}
\begin{scope}[scale=.6666667,draw=black,text=black]\draw[draw=black,line width=.35pt] (295.267,705.000)--(295.267,711.000);\end{scope}
\begin{scope}[scale=.6666667,draw=black,text=black]\draw[draw=black,line width=.35pt] (300.033,705.000)--(300.033,711.000);\end{scope}
\begin{scope}[scale=.6666667,draw=black,text=black]\draw[draw=black,line width=.35pt] (304.800,705.000)--(304.800,711.000);\end{scope}
\begin{scope}[scale=.6666667,draw=black,text=black]\draw[draw=black,line width=.35pt] (309.567,705.000)--(309.567,711.000);\end{scope}
\begin{scope}[scale=.6666667,draw=black,text=black]\draw[draw=black,line width=.35pt] (314.333,705.000)--(314.333,719.000);\end{scope}
\begin{scope}[scale=.6666667,draw=black,text=black]\draw[draw=black,line width=.35pt] (319.100,705.000)--(319.100,711.000);\end{scope}
\begin{scope}[scale=.6666667,draw=black,text=black]\draw[draw=black,line width=.35pt] (323.867,705.000)--(323.867,711.000);\end{scope}
\begin{scope}[scale=.6666667,draw=black,text=black]\draw[draw=black,line width=.35pt] (328.633,705.000)--(328.633,711.000);\end{scope}
\begin{scope}[scale=.6666667,draw=black,text=black]\draw[draw=black,line width=.35pt] (333.400,705.000)--(333.400,711.000);\end{scope}
\begin{scope}[scale=.6666667,draw=black,text=black]\draw[draw=black,line width=.35pt] (338.167,705.000)--(338.167,714.000);\end{scope}
\begin{scope}[scale=.6666667,draw=black,text=black]\draw[draw=black,line width=.35pt] (342.933,705.000)--(342.933,711.000);\end{scope}
\begin{scope}[scale=.6666667,draw=black,text=black]\draw[draw=black,line width=.35pt] (347.700,705.000)--(347.700,711.000);\end{scope}
\begin{scope}[scale=.6666667,draw=black,text=black]\draw[draw=black,line width=.35pt] (352.467,705.000)--(352.467,711.000);\end{scope}
\begin{scope}[scale=.6666667,draw=black,text=black]\draw[draw=black,line width=.35pt] (357.233,705.000)--(357.233,711.000);\end{scope}
\begin{scope}[scale=.6666667,draw=black,text=black]\draw[draw=black,line width=.35pt] (362.000,705.000)--(362.000,719.000);\end{scope}
\begin{scope}[scale=.6666667,draw=black,text=black]\draw[draw=black,line width=.35pt] (366.767,705.000)--(366.767,711.000);\end{scope}
\begin{scope}[scale=.6666667,draw=black,text=black]\draw[draw=black,line width=.35pt] (371.533,705.000)--(371.533,711.000);\end{scope}
\begin{scope}[scale=.6666667,draw=black,text=black]\draw[draw=black,line width=.35pt] (376.300,705.000)--(376.300,711.000);\end{scope}
\begin{scope}[scale=.6666667,draw=black,text=black]\draw[draw=black,line width=.35pt] (381.067,705.000)--(381.067,711.000);\end{scope}
\begin{scope}[scale=.6666667,draw=black,text=black]\draw[draw=black,line width=.35pt] (385.833,705.000)--(385.833,714.000);\end{scope}
\begin{scope}[scale=.6666667,draw=black,text=black]\draw[draw=black,line width=.35pt] (390.600,705.000)--(390.600,711.000);\end{scope}
\begin{scope}[scale=.6666667,draw=black,text=black]\draw[draw=black,line width=.35pt] (395.367,705.000)--(395.367,711.000);\end{scope}
\begin{scope}[scale=.6666667,draw=black,text=black]\draw[draw=black,line width=.35pt] (400.133,705.000)--(400.133,711.000);\end{scope}
\begin{scope}[scale=.6666667,draw=black,text=black]\draw[draw=black,line width=.35pt] (404.900,705.000)--(404.900,711.000);\end{scope}
\begin{scope}[scale=.6666667,draw=black,text=black]\draw[draw=black,line width=.35pt] (409.667,705.000)--(409.667,719.000);\end{scope}
\begin{scope}[scale=.6666667,draw=black,text=black]\draw[draw=black,line width=.35pt] (414.433,705.000)--(414.433,711.000);\end{scope}
\begin{scope}[scale=.6666667,draw=black,text=black]\draw[draw=black,line width=.35pt] (419.200,705.000)--(419.200,711.000);\end{scope}
\begin{scope}[scale=.6666667,draw=black,text=black]\draw[draw=black,line width=.35pt] (423.967,705.000)--(423.967,711.000);\end{scope}
\begin{scope}[scale=.6666667,draw=black,text=black]\draw[draw=black,line width=.35pt] (428.733,705.000)--(428.733,711.000);\end{scope}
\begin{scope}[scale=.6666667,draw=black,text=black]\draw[draw=black,line width=.35pt] (433.500,705.000)--(433.500,714.000);\end{scope}
\begin{scope}[scale=.6666667,draw=black,text=black]\draw[draw=black,line width=.35pt] (438.267,705.000)--(438.267,711.000);\end{scope}
\begin{scope}[scale=.6666667,draw=black,text=black]\draw[draw=black,line width=.35pt] (443.033,705.000)--(443.033,711.000);\end{scope}
\begin{scope}[scale=.6666667,draw=black,text=black]\draw[draw=black,line width=.35pt] (447.800,705.000)--(447.800,711.000);\end{scope}
\begin{scope}[scale=.6666667,draw=black,text=black]\draw[draw=black,line width=.35pt] (452.567,705.000)--(452.567,711.000);\end{scope}
\begin{scope}[scale=.6666667,draw=black,text=black]\draw[draw=black,line width=.35pt] (457.333,705.000)--(457.333,719.000);\end{scope}
\begin{scope}[scale=.6666667,draw=black,text=black]\draw[draw=black,line width=.35pt] (462.100,705.000)--(462.100,711.000);\end{scope}
\begin{scope}[scale=.6666667,draw=black,text=black]\draw[draw=black,line width=.35pt] (466.867,705.000)--(466.867,711.000);\end{scope}
\begin{scope}[scale=.6666667,draw=black,text=black]\draw[draw=black,line width=.35pt] (471.633,705.000)--(471.633,711.000);\end{scope}
\begin{scope}[scale=.6666667,draw=black,text=black]\draw[draw=black,line width=.35pt] (476.400,705.000)--(476.400,711.000);\end{scope}
\begin{scope}[scale=.6666667,draw=black,text=black]\draw[draw=black,line width=.35pt] (481.167,705.000)--(481.167,714.000);\end{scope}
\begin{scope}[scale=.6666667,draw=black,text=black]\draw[draw=black,line width=.35pt] (485.933,705.000)--(485.933,711.000);\end{scope}
\begin{scope}[scale=.6666667,draw=black,text=black]\draw[draw=black,line width=.35pt] (490.700,705.000)--(490.700,711.000);\end{scope}
\begin{scope}[scale=.6666667,draw=black,text=black]\draw[draw=black,line width=.35pt] (495.467,705.000)--(495.467,711.000);\end{scope}
\begin{scope}[scale=.6666667,draw=black,text=black]\draw[draw=black,line width=.35pt] (500.233,705.000)--(500.233,711.000);\end{scope}
\begin{scope}[scale=.6666667,draw=black,text=black]\draw[draw=black,line width=.35pt] (505.000,705.000)--(505.000,719.000);\end{scope}
\begin{scope}[scale=.6666667,draw=black,text=black]\draw[draw=black,line width=.35pt] (509.767,705.000)--(509.767,711.000);\end{scope}
\begin{scope}[scale=.6666667,draw=black,text=black]\draw[draw=black,line width=.35pt] (514.533,705.000)--(514.533,711.000);\end{scope}
\begin{scope}[scale=.6666667,draw=black,text=black]\draw[draw=black,line width=.35pt] (519.300,705.000)--(519.300,711.000);\end{scope}
\begin{scope}[scale=.6666667,draw=black,text=black]\draw[draw=black,line width=.35pt] (524.067,705.000)--(524.067,711.000);\end{scope}
\begin{scope}[scale=.6666667,draw=black,text=black]\draw[draw=black,line width=.35pt] (528.833,705.000)--(528.833,714.000);\end{scope}
\begin{scope}[scale=.6666667,draw=black,text=black]\draw[draw=black,line width=.35pt] (533.600,705.000)--(533.600,711.000);\end{scope}
\begin{scope}[scale=.6666667,draw=black,text=black]\draw[draw=black,line width=.35pt] (538.367,705.000)--(538.367,711.000);\end{scope}
\begin{scope}[scale=.6666667,draw=black,text=black]\draw[draw=black,line width=.35pt] (543.133,705.000)--(543.133,711.000);\end{scope}
\begin{scope}[scale=.6666667,draw=black,text=black]\draw[draw=black,line width=.35pt] (547.900,705.000)--(547.900,711.000);\end{scope}
\begin{scope}[scale=.6666667,draw=black,text=black]\draw[draw=black,line width=.35pt] (552.667,705.000)--(552.667,719.000);\end{scope}
\begin{scope}[scale=.6666667,draw=black,text=black]\draw[draw=black,line width=.35pt] (557.433,705.000)--(557.433,711.000);\end{scope}
\begin{scope}[scale=.6666667,draw=black,text=black]\draw[draw=black,line width=.35pt] (562.200,705.000)--(562.200,711.000);\end{scope}
\begin{scope}[scale=.6666667,draw=black,text=black]\draw[draw=black,line width=.35pt] (566.967,705.000)--(566.967,711.000);\end{scope}
\begin{scope}[scale=.6666667,draw=black,text=black]\draw[draw=black,line width=.35pt] (571.733,705.000)--(571.733,711.000);\end{scope}
\begin{scope}[scale=.6666667,draw=black,text=black]\draw[draw=black,line width=.35pt] (576.500,705.000)--(576.500,714.000);\end{scope}
\begin{scope}[scale=.6666667,draw=black,text=black]\draw[draw=black,line width=.35pt] (581.267,705.000)--(581.267,711.000);\end{scope}
\begin{scope}[scale=.6666667,draw=black,text=black]\draw[draw=black,line width=.35pt] (586.033,705.000)--(586.033,711.000);\end{scope}
\begin{scope}[scale=.6666667,draw=black,text=black]\draw[draw=black,line width=.35pt] (590.800,705.000)--(590.800,711.000);\end{scope}
\begin{scope}[scale=.6666667,draw=black,text=black]\draw[draw=black,line width=.35pt] (595.567,705.000)--(595.567,711.000);\end{scope}
\begin{scope}[scale=.6666667,draw=black,text=black]\draw[draw=black,line width=.35pt] (600.333,705.000)--(600.333,719.000);\end{scope}
\begin{scope}[scale=.6666667,draw=black,text=black]\draw[draw=black,line width=.35pt] (605.100,705.000)--(605.100,711.000);\end{scope}
\begin{scope}[scale=.6666667,draw=black,text=black]\draw[draw=black,line width=.35pt] (609.867,705.000)--(609.867,711.000);\end{scope}
\begin{scope}[scale=.6666667,draw=black,text=black]\draw[draw=black,line width=.35pt] (614.633,705.000)--(614.633,711.000);\end{scope}
\begin{scope}[scale=.6666667,draw=black,text=black]\draw[draw=black,line width=.35pt] (619.400,705.000)--(619.400,711.000);\end{scope}
\begin{scope}[scale=.6666667,draw=black,text=black]\draw[draw=black,line width=.35pt] (624.167,705.000)--(624.167,714.000);\end{scope}
\begin{scope}[scale=.6666667,draw=black,text=black]\draw[draw=black,line width=.35pt] (628.933,705.000)--(628.933,711.000);\end{scope}
\begin{scope}[scale=.6666667,draw=black,text=black]\draw[draw=black,line width=.35pt] (633.700,705.000)--(633.700,711.000);\end{scope}
\begin{scope}[scale=.6666667,draw=black,text=black]\draw[draw=black,line width=.35pt] (638.467,705.000)--(638.467,711.000);\end{scope}
\begin{scope}[scale=.6666667,draw=black,text=black]\draw[draw=black,line width=.35pt] (643.233,705.000)--(643.233,711.000);\end{scope}
\begin{scope}[scale=.6666667,draw=black,text=black]\draw[draw=black,line width=.35pt] (648.000,705.000)--(648.000,719.000);\end{scope}
\begin{scope}[scale=.6666667,draw=black,text=black]\draw[draw=black,line width=.35pt] (652.767,705.000)--(652.767,711.000);\end{scope}
\begin{scope}[scale=.6666667,draw=black,text=black]\draw[draw=black,line width=.35pt] (657.533,705.000)--(657.533,711.000);\end{scope}
\begin{scope}[scale=.6666667,draw=black,text=black]\draw[draw=black,line width=.35pt] (662.300,705.000)--(662.300,711.000);\end{scope}
\begin{scope}[scale=.6666667,draw=black,text=black]\draw[draw=black,line width=.35pt] (667.067,705.000)--(667.067,711.000);\end{scope}
\begin{scope}[scale=.6666667,draw=black,text=black]\draw[draw=black,line width=.35pt] (671.833,705.000)--(671.833,714.000);\end{scope}
\begin{scope}[scale=.6666667,draw=black,text=black]\draw[draw=black,line width=.35pt] (676.600,705.000)--(676.600,711.000);\end{scope}
\begin{scope}[scale=.6666667,draw=black,text=black]\draw[draw=black,line width=.35pt] (681.367,705.000)--(681.367,711.000);\end{scope}
\begin{scope}[scale=.6666667,draw=black,text=black]\draw[draw=black,line width=.35pt] (686.133,705.000)--(686.133,711.000);\end{scope}
\begin{scope}[scale=.6666667,draw=black,text=black]\draw[draw=black,line width=.35pt] (690.900,705.000)--(690.900,711.000);\end{scope}
\begin{scope}[scale=.6666667,draw=black,text=black]\draw[draw=black,line width=.35pt] (695.667,705.000)--(695.667,719.000);\end{scope}
\begin{scope}[scale=.6666667,draw=black,text=black]\node[inner sep=0pt,font=\normalsize,text=black] at (219.000,738.000) {0};\end{scope}
\begin{scope}[scale=.6666667,draw=black,text=black]\node[inner sep=0pt,font=\normalsize,text=black] at (266.667,738.000) {1 cm};\end{scope}
\begin{scope}[scale=.6666667,draw=black,text=black]\node[inner sep=0pt,font=\normalsize,text=black] at (314.333,738.000) {2};\end{scope}
\begin{scope}[scale=.6666667,draw=black,text=black]\node[inner sep=0pt,font=\normalsize,text=black] at (362.000,738.000) {3};\end{scope}
\begin{scope}[scale=.6666667,draw=black,text=black]\node[inner sep=0pt,font=\normalsize,text=black] at (409.667,738.000) {4};\end{scope}
\begin{scope}[scale=.6666667,draw=black,text=black]\node[inner sep=0pt,font=\normalsize,text=black] at (457.333,738.000) {5};\end{scope}
\begin{scope}[scale=.6666667,draw=black,text=black]\node[inner sep=0pt,font=\normalsize,text=black] at (505.000,738.000) {6};\end{scope}
\begin{scope}[scale=.6666667,draw=black,text=black]\node[inner sep=0pt,font=\normalsize,text=black] at (552.667,738.000) {7};\end{scope}
\begin{scope}[scale=.6666667,draw=black,text=black]\node[inner sep=0pt,font=\normalsize,text=black] at (600.333,738.000) {8};\end{scope}
\begin{scope}[scale=.6666667,draw=black,text=black]\node[inner sep=0pt,font=\normalsize,text=black] at (648.000,738.000) {9};\end{scope}
\begin{scope}[scale=.6666667,draw=black,text=black]\node[inner sep=0pt,font=\normalsize,text=black] at (695.667,738.000) {10};\end{scope}
\begin{scope}[scale=.6666667,draw=black,text=black]\draw[draw=black,line width=.45pt,-{Stealth[length=3pt,width=2.4pt]}] (267.000,760.000)--(267.000,785.000);\end{scope}
\begin{scope}[scale=.6666667,draw=black,text=black]\draw[draw=black,line width=.45pt,-{Stealth[length=3pt,width=2.4pt]}] (362.000,760.000)--(362.000,785.000);\end{scope}
\begin{scope}[scale=.6666667,draw=black,text=black]\draw[draw=black,line width=.45pt,-{Stealth[length=3pt,width=2.4pt]}] (505.000,760.000)--(505.000,785.000);\end{scope}
\q{230}{790}{85}{12.5}{$\dfrac1{10}$分米}
\q{331}{790}{96}{12.5}{$\dfrac3{10}$分米}
\q{463}{804}{28}{13}{（}
\q{512}{804}{25}{13}{）}
\q{533}{804}{80}{12.5}{分米}
\q{225}{862}{103}{12.5}{0.1分米}
\q{324}{861}{28}{13}{（}
\q{378}{861}{25}{13}{）}
\q{399}{861}{76}{12.5}{分米}
\q{463}{861}{28}{13}{（}
\q{514}{861}{25}{13}{）}
\q{535}{861}{75}{12.5}{分米}
\hdsection{4}{921}{在计数器上表示下面各小数。}
\q{112}{967}{172}{13}{① 0.35}
\q{368}{967}{172}{13}{② 4.3}
\q{623}{967}{172}{13}{③ 0.732}
\begin{scope}[scale=.6666667,draw=black,text=black]\draw[line width=.5pt] (143,1015) rectangle (287,1079);\draw[line width=.5pt] (143,1019)--(287,1019);\end{scope}
\begin{scope}[scale=.6666667,draw=black,text=black]\draw[line width=.45pt] (165,1017)--(165,1079);\draw[line width=.45pt] (157,1015)..controls(166,995) and (171,1003)..(170,1015);\end{scope}
\begin{scope}[scale=.6666667,draw=black,text=black]\draw[line width=.45pt] (191,1017)--(191,1079);\draw[line width=.45pt] (183,1015)..controls(192,995) and (197,1003)..(196,1015);\end{scope}
\begin{scope}[scale=.6666667,draw=black,text=black]\draw[line width=.45pt] (224,1017)--(224,1079);\draw[line width=.45pt] (216,1015)..controls(225,995) and (230,1003)..(229,1015);\end{scope}
\begin{scope}[scale=.6666667,draw=black,text=black]\draw[line width=.45pt] (247,1017)--(247,1079);\draw[line width=.45pt] (239,1015)..controls(248,995) and (253,1003)..(252,1015);\end{scope}
\begin{scope}[scale=.6666667,draw=black,text=black]\draw[line width=.45pt] (273,1017)--(273,1079);\draw[line width=.45pt] (265,1015)..controls(274,995) and (279,1003)..(278,1015);\end{scope}
\begin{scope}[scale=.6666667,draw=black,text=black]\draw[draw=black,line width=.5pt,fill=white] (143.000,1078.000)--(151.000,1083.000)--(294.000,1083.000)--(287.000,1078.000)--cycle;\end{scope}
\begin{scope}[scale=.6666667,draw=black,text=black]\draw[draw=black,line width=.5pt,fill=white] (151.000,1083.000)--(294.000,1083.000)--(294.000,1147.000)--(151.000,1147.000)--cycle;\end{scope}
\begin{scope}[scale=.6666667,draw=black,text=black]\draw[draw=black,line width=.5pt,fill=white] (143.000,1078.000)--(151.000,1083.000)--(151.000,1147.000)--(143.000,1131.000)--cycle;\end{scope}
\begin{scope}[scale=.6666667,draw=black,text=black]\node[inner sep=0pt,font=\normalsize,text=black] at (165.000,1115.000) {\shortstack{十\\位}};\end{scope}
\begin{scope}[scale=.6666667,draw=black,text=black]\node[inner sep=0pt,font=\normalsize,text=black] at (191.000,1115.000) {\shortstack{个\\位}};\end{scope}
\begin{scope}[scale=.6666667,draw=black,text=black]\node[inner sep=0pt,font=\normalsize,text=black] at (224.000,1115.000) {\shortstack{十\\分\\位}};\end{scope}
\begin{scope}[scale=.6666667,draw=black,text=black]\node[inner sep=0pt,font=\normalsize,text=black] at (247.000,1115.000) {\shortstack{百\\分\\位}};\end{scope}
\begin{scope}[scale=.6666667,draw=black,text=black]\node[inner sep=0pt,font=\normalsize,text=black] at (273.000,1115.000) {\shortstack{千\\分\\位}};\end{scope}
\begin{scope}[scale=.6666667,draw=black,text=black]\draw[line width=.5pt] (398,1015) rectangle (542,1079);\draw[line width=.5pt] (398,1019)--(542,1019);\end{scope}
\begin{scope}[scale=.6666667,draw=black,text=black]\draw[line width=.45pt] (416,1017)--(416,1079);\draw[line width=.45pt] (408,1015)..controls(417,995) and (422,1003)..(421,1015);\end{scope}
\begin{scope}[scale=.6666667,draw=black,text=black]\draw[line width=.45pt] (441,1017)--(441,1079);\draw[line width=.45pt] (433,1015)..controls(442,995) and (447,1003)..(446,1015);\end{scope}
\begin{scope}[scale=.6666667,draw=black,text=black]\draw[line width=.45pt] (477,1017)--(477,1079);\draw[line width=.45pt] (469,1015)..controls(478,995) and (483,1003)..(482,1015);\end{scope}
\begin{scope}[scale=.6666667,draw=black,text=black]\draw[line width=.45pt] (502,1017)--(502,1079);\draw[line width=.45pt] (494,1015)..controls(503,995) and (508,1003)..(507,1015);\end{scope}
\begin{scope}[scale=.6666667,draw=black,text=black]\draw[line width=.45pt] (527,1017)--(527,1079);\draw[line width=.45pt] (519,1015)..controls(528,995) and (533,1003)..(532,1015);\end{scope}
\begin{scope}[scale=.6666667,draw=black,text=black]\draw[draw=black,line width=.5pt,fill=white] (398.000,1078.000)--(406.000,1083.000)--(549.000,1083.000)--(542.000,1078.000)--cycle;\end{scope}
\begin{scope}[scale=.6666667,draw=black,text=black]\draw[draw=black,line width=.5pt,fill=white] (406.000,1083.000)--(549.000,1083.000)--(549.000,1147.000)--(406.000,1147.000)--cycle;\end{scope}
\begin{scope}[scale=.6666667,draw=black,text=black]\draw[draw=black,line width=.5pt,fill=white] (398.000,1078.000)--(406.000,1083.000)--(406.000,1147.000)--(398.000,1131.000)--cycle;\end{scope}
\begin{scope}[scale=.6666667,draw=black,text=black]\node[inner sep=0pt,font=\normalsize,text=black] at (416.000,1115.000) {\shortstack{十\\位}};\end{scope}
\begin{scope}[scale=.6666667,draw=black,text=black]\node[inner sep=0pt,font=\normalsize,text=black] at (441.000,1115.000) {\shortstack{个\\位}};\end{scope}
\begin{scope}[scale=.6666667,draw=black,text=black]\node[inner sep=0pt,font=\normalsize,text=black] at (477.000,1115.000) {\shortstack{十\\分\\位}};\end{scope}
\begin{scope}[scale=.6666667,draw=black,text=black]\node[inner sep=0pt,font=\normalsize,text=black] at (502.000,1115.000) {\shortstack{百\\分\\位}};\end{scope}
\begin{scope}[scale=.6666667,draw=black,text=black]\node[inner sep=0pt,font=\normalsize,text=black] at (527.000,1115.000) {\shortstack{千\\分\\位}};\end{scope}
\begin{scope}[scale=.6666667,draw=black,text=black]\draw[line width=.5pt] (654,1015) rectangle (798,1079);\draw[line width=.5pt] (654,1019)--(798,1019);\end{scope}
\begin{scope}[scale=.6666667,draw=black,text=black]\draw[line width=.45pt] (675,1017)--(675,1079);\draw[line width=.45pt] (667,1015)..controls(676,995) and (681,1003)..(680,1015);\end{scope}
\begin{scope}[scale=.6666667,draw=black,text=black]\draw[line width=.45pt] (701,1017)--(701,1079);\draw[line width=.45pt] (693,1015)..controls(702,995) and (707,1003)..(706,1015);\end{scope}
\begin{scope}[scale=.6666667,draw=black,text=black]\draw[line width=.45pt] (731,1017)--(731,1079);\draw[line width=.45pt] (723,1015)..controls(732,995) and (737,1003)..(736,1015);\end{scope}
\begin{scope}[scale=.6666667,draw=black,text=black]\draw[line width=.45pt] (757,1017)--(757,1079);\draw[line width=.45pt] (749,1015)..controls(758,995) and (763,1003)..(762,1015);\end{scope}
\begin{scope}[scale=.6666667,draw=black,text=black]\draw[line width=.45pt] (781,1017)--(781,1079);\draw[line width=.45pt] (773,1015)..controls(782,995) and (787,1003)..(786,1015);\end{scope}
\begin{scope}[scale=.6666667,draw=black,text=black]\draw[draw=black,line width=.5pt,fill=white] (654.000,1078.000)--(662.000,1083.000)--(805.000,1083.000)--(798.000,1078.000)--cycle;\end{scope}
\begin{scope}[scale=.6666667,draw=black,text=black]\draw[draw=black,line width=.5pt,fill=white] (662.000,1083.000)--(805.000,1083.000)--(805.000,1147.000)--(662.000,1147.000)--cycle;\end{scope}
\begin{scope}[scale=.6666667,draw=black,text=black]\draw[draw=black,line width=.5pt,fill=white] (654.000,1078.000)--(662.000,1083.000)--(662.000,1147.000)--(654.000,1131.000)--cycle;\end{scope}
\begin{scope}[scale=.6666667,draw=black,text=black]\node[inner sep=0pt,font=\normalsize,text=black] at (675.000,1115.000) {\shortstack{十\\位}};\end{scope}
\begin{scope}[scale=.6666667,draw=black,text=black]\node[inner sep=0pt,font=\normalsize,text=black] at (701.000,1115.000) {\shortstack{个\\位}};\end{scope}
\begin{scope}[scale=.6666667,draw=black,text=black]\node[inner sep=0pt,font=\normalsize,text=black] at (731.000,1115.000) {\shortstack{十\\分\\位}};\end{scope}
\begin{scope}[scale=.6666667,draw=black,text=black]\node[inner sep=0pt,font=\normalsize,text=black] at (757.000,1115.000) {\shortstack{百\\分\\位}};\end{scope}
\begin{scope}[scale=.6666667,draw=black,text=black]\node[inner sep=0pt,font=\normalsize,text=black] at (781.000,1115.000) {\shortstack{千\\分\\位}};\end{scope}

\begin{scope}[color=answerred,draw=answerred,text=answerred]\ans{128.00}{226.67}{33.33}{10}{$\vphantom{0}0.4$}
\ans{294.00}{226.67}{33.33}{10}{$\vphantom{0}0.33$}
\ans{455.33}{225.33}{33.33}{10}{$\vphantom{0}0.203$}
\fill[answerred,opacity=.28](95.333,296.000) rectangle (181.733,394.000);
\fill[answerred,opacity=.28](259.333,296.000) rectangle (345.733,305.733);
\fill[answerred,opacity=.28](424.000,296.000) rectangle (472.667,393.333);
\fill[answerred,opacity=.28](472.667,296.000) rectangle (482.400,325.200);
\ans{228.67}{573.33}{48.00}{11}{$\vphantom{0}0.3$}
\ans{322.67}{534.00}{42.00}{11}{$\vphantom{0}\frac6{10}$}
\ans{321.33}{573.33}{48.00}{11}{$\vphantom{0}0.6$}
\filldraw[answerred,fill=answerred!65] (149.333,714.000) ellipse [x radius=4bp,y radius=2.2bp];
\filldraw[answerred,fill=answerred!65] (149.333,708.667) ellipse [x radius=4bp,y radius=2.2bp];
\filldraw[answerred,fill=answerred!65] (149.333,703.333) ellipse [x radius=4bp,y radius=2.2bp];
\filldraw[answerred,fill=answerred!65] (164.667,714.000) ellipse [x radius=4bp,y radius=2.2bp];
\filldraw[answerred,fill=answerred!65] (164.667,708.667) ellipse [x radius=4bp,y radius=2.2bp];
\filldraw[answerred,fill=answerred!65] (164.667,703.333) ellipse [x radius=4bp,y radius=2.2bp];
\filldraw[answerred,fill=answerred!65] (164.667,698.000) ellipse [x radius=4bp,y radius=2.2bp];
\filldraw[answerred,fill=answerred!65] (164.667,692.667) ellipse [x radius=4bp,y radius=2.2bp];
\filldraw[answerred,fill=answerred!65] (294.000,714.000) ellipse [x radius=4bp,y radius=2.2bp];
\filldraw[answerred,fill=answerred!65] (294.000,708.667) ellipse [x radius=4bp,y radius=2.2bp];
\filldraw[answerred,fill=answerred!65] (294.000,703.333) ellipse [x radius=4bp,y radius=2.2bp];
\filldraw[answerred,fill=answerred!65] (294.000,698.000) ellipse [x radius=4bp,y radius=2.2bp];
\filldraw[answerred,fill=answerred!65] (318.000,714.000) ellipse [x radius=4bp,y radius=2.2bp];
\filldraw[answerred,fill=answerred!65] (318.000,708.667) ellipse [x radius=4bp,y radius=2.2bp];
\filldraw[answerred,fill=answerred!65] (318.000,703.333) ellipse [x radius=4bp,y radius=2.2bp];
\filldraw[answerred,fill=answerred!65] (487.333,714.000) ellipse [x radius=4bp,y radius=2.2bp];
\filldraw[answerred,fill=answerred!65] (487.333,708.667) ellipse [x radius=4bp,y radius=2.2bp];
\filldraw[answerred,fill=answerred!65] (487.333,703.333) ellipse [x radius=4bp,y radius=2.2bp];
\filldraw[answerred,fill=answerred!65] (487.333,698.000) ellipse [x radius=4bp,y radius=2.2bp];
\filldraw[answerred,fill=answerred!65] (487.333,692.667) ellipse [x radius=4bp,y radius=2.2bp];
\filldraw[answerred,fill=answerred!65] (487.333,687.333) ellipse [x radius=4bp,y radius=2.2bp];
\filldraw[answerred,fill=answerred!65] (487.333,682.000) ellipse [x radius=4bp,y radius=2.2bp];
\filldraw[answerred,fill=answerred!65] (504.667,714.000) ellipse [x radius=4bp,y radius=2.2bp];
\filldraw[answerred,fill=answerred!65] (504.667,708.667) ellipse [x radius=4bp,y radius=2.2bp];
\filldraw[answerred,fill=answerred!65] (504.667,703.333) ellipse [x radius=4bp,y radius=2.2bp];
\filldraw[answerred,fill=answerred!65] (520.667,714.000) ellipse [x radius=4bp,y radius=2.2bp];
\filldraw[answerred,fill=answerred!65] (520.667,708.667) ellipse [x radius=4bp,y radius=2.2bp];
\end{scope}
\end{scope}\pgfresetboundingbox\path[use as bounding box] (69.0,639.0) rectangle (542.0,770.0);
\end{tikzpicture}}\par



\par






